% Style

\documentclass[journal, 12pt, onecolumn]{IEEEtran}
\usepackage[utf8]{inputenc}
\usepackage[T1]{fontenc}
\usepackage{amsmath}
\usepackage{graphicx}
\usepackage{tikz}
\usepackage{pgfplots}
\usepackage{bbm}
\usepackage{environ}
% Functionallity

% LaTeX package for machine learning research papers.

\usepackage{amsmath,amsfonts} % Mathematical equations, font and symbols
\usepackage{amssymb} % Additional mathematical symbols
\usepackage{graphicx} % Including figures and plots
\usepackage{booktabs} % Professional-looking tables
\usepackage{algorithm2e} % Formatting machine learning algorithms
\usepackage{listings} % Code snippet formatting
\usepackage{caption} % Customized captions for figures and tables
\usepackage{color} % Adding color to text and elements
%\usepackage{biblatex} % Advanced citation management
\usepackage{textcomp} % Different symbols
\usepackage{mathrsfs}
\usepackage{enumitem} % alphabet enumerating

\usepackage{mathtools}
\mathtoolsset{showonlyrefs}
\usepackage{etoolbox}
%%%%%%%%%%%%%%%%%%%%%%%%%%%%%%%%%%%%%%
%        macros.tex                  %
%%%%%%%%%%%%%%%%%%%%%%%%%%%%%%%%%%%%%%

%From PhD thesis
\usepackage{amsmath, amssymb, bbm, xspace}
\usepackage{epsfig}
\usepackage{longtable}
\usepackage{color}
\usepackage{mathrsfs}
\usepackage{comment}
\usepackage{ifthen}
\newboolean{showcomments}
\setboolean{showcomments}{true}
%\newcommand{\aak}[1]{  \ifthenelse{\boolean{showcomments}}
%{ \textcolor{blue}{(AAK says:  #1)}} {}  }
\def\sthat{\ {\rm such\ that}\ }
%\usepackage{fancyhdr}
% \usepackage[T1]{fontenc}
% \usepackage{arev}
%\usepackage{helvet}
\usepackage{courier}
%\usepackage{courier}
% \usepackage[T1]{fontenc}
% \usepackage[sc]{mathpazo}
% \linespread{1.05}         % Palatino needs more leading (space between lines)
\usepackage{xcolor}
\usepackage{todonotes}
\usepackage[framemethod=tikz]{mdframed}
\usepackage{lineno}
\newmdenv[leftmargin=\dimexpr-0.4em, innerleftmargin=0.5em,
rightmargin=\dimexpr-0.4em, innerrightmargin=0.5em,
linewidth=2pt,linecolor=red, topline=false, bottomline=false,
innertopmargin=0pt,innerbottommargin=0pt,skipbelow=0pt,skipabove=0pt,%
]{notex}

\newenvironment{note}%
{\vskip\dimexpr\dp\strutbox-\prevdepth\relax\notex\strut\ignorespaces}%
{\xdef\notetpd{\the\prevdepth}\endnotex\vskip-\notetpd\relax}

\let\oldtodo\todo

\makeatletter%
\DeclareDocumentCommand{\todo}{ O{} +g +d<> }{%
		\setlength{\marginparwidth}{1.5cm}%
	%
	\IfNoValueTF{#2}{\relax}{%
		\oldtodo[caption={#2},size=\scriptsize,#1]{\renewcommand{\baselinestretch}{0.8}\selectfont\sffamily#2\par}%
	}%
	\IfNoValueTF{#3}{\relax}{%
		\IfNoValueTF{#2}{% when parmark but no todo
			\begin{note}%
				\begin{internallinenumbers}%
					\indent%
					#3%
				\end{internallinenumbers}%
			\end{note}%
		}{% when parmark and todo
			\vspace{-0\baselineskip}%
			\begin{note}%
				\begin{internallinenumbers}%
					\indent%
					#3%
				\end{internallinenumbers}%
			\end{note}%
		}%
	}%
}%
\makeatother
\usepackage{soul}

\newcommand{\hlc}[2][yellow]{{%
		\colorlet{foo}{#1}%
		\sethlcolor{foo}\hl{#2}}%
}

\newtheorem{conjecture}{Conjecture}[section]

\newcommand{\removetodo}[2]{\todo[color=pink]{\textbf{delete:} ``#1'' #2}\hlc[pink]{#1}}
\newcommand{\inserttodo}[1]{\todo[color=green!40]{\textbf{insert:} #1}}


\newcommand{\hltodoy}[2]{\todo[color=yellow!40]{#2}\hl{#1} }

\newcommand{\hltodoc}[3]{\todo[color=#3!40]{#2}\hlc[#3]{#1} }

\newcommand{\hltodo}[2]{\todo[color=orange!40]{#2}\hlc[orange!40]{#1} }
\newcommand{\replacetodo}[2]{\todo[color=pink!40]{\textbf{replace with:}``#2'' }\hl{#1} }

\usepackage{marginnote}
\newcommand{\todol}[1]{{%
		\let\marginpar\marginnote
		\reversemarginpar
		\renewcommand{\baselinestretch}{0.8}%
		\todo{#1}}}

\newcommand{\inserttodol}[1]{{%
		\let\marginpar\marginnote
		\reversemarginpar
		\renewcommand{\baselinestretch}{0.8}%
		\inserttodo{#1}}}

\newcommand{\removetodol}[2]{{%
		\let\marginpar\marginnote
		\reversemarginpar
		\renewcommand{\baselinestretch}{0.8}%
		\removetodo{#1}{#2}}}

\newcommand{\hltodol}[2]{{%
		\let\marginpar\marginnote
		\reversemarginpar
		\renewcommand{\baselinestretch}{0.8}%
		\hltodo{#1}{#2}}}

\newcommand{\replacetodol}[2]{{%
		\let\marginpar\marginnote
		\reversemarginpar
		\renewcommand{\baselinestretch}{0.8}%
		\replacetodo{#1}{#2}}}

\newcommand{\hltodoyl}[2]{{%
		\let\marginpar\marginnote
		\reversemarginpar
		\renewcommand{\baselinestretch}{0.8}%
		\hltodoy{#1}{#2}}}

\newcommand{\hltodocl}[3]{{		\let\marginpar\marginnote
		\reversemarginpar
		\renewcommand{\baselinestretch}{0.8}%
		\hltodoc{#1}{#2}{#3}}}

\newcommand{\tododone}[1]{\todo[nolist,color=gray!20,bordercolor=gray!60]{#1}}
% \usepackage{titlesec}
% \titlelabel{\qed \thetitle \quad}
% \titleformat*{\section}{\itshape}

% This abstract environment is a little narrower
% than the default.  Note the capital A...
% \newenvironment{Abstract}
%   {\begin{center}
%     \textbf{Abstract} \\
%     \medskip
%   \begin{minipage}{4.8in}}
%   {\end{minipage}
% \end{center}}



\newtheorem{theorem}{Theorem}[section]
\newtheorem{assumption}{Assumption}[section]
\newtheorem{lemma}[theorem]{Lemma}
\newtheorem{claim}[theorem]{Claim}
\newtheorem{proposition}[theorem]{Proposition}
\newtheorem{prop}[theorem]{Proposition}
\newtheorem{corollary}[theorem]{Corollary}
\newtheorem{cor}[theorem]{Corollary}
\newtheorem{result}[theorem]{Result}
\newtheorem{definition}{Definition}[section]

\newcommand{\comb}[2]{\prescript{#1}{}C_{#2}}

%\newcounter{theorem}[chapter]
% \renewcommand{\thetheorem}{\thechapter.\arabic{theorem}}
%
% \renewenvironment{theorem}{\refstepcounter{theorem} {\ \\[2ex]}{\bf \large \textsf{Theorem \thetheorem}} \newline \it}{}

\def\bkE{{\rm I\kern-.17em E}}
\def\bk1{{\rm 1\kern-.17em l}}
\def\bkD{{\rm I\kern-.17em D}}
\def\bkR{{\rm I\kern-.17em R}}
\def\bkP{{\rm I\kern-.17em P}}
\def\Cov{{\rm Cov}}
\def\Cbb{{\mathbb{C}}}
\def\Fbb{{\mathbb{F}}}
\def\C{{\mathbb{C}}}
\def\Abb{{\mathbb{A}}}
\def\Mbb{{\mathbb{M}}}
\def\Bbb{{\mathbb{B}}}
\def\Vbb{{\mathbb{V}}}
\def\Mhatbb{{\mathbb{\widehat{M}}}}
\def\fnat{{\bf F}^{\textrm{nat}}}
\def\xref{x^{\textrm{ref}}}
\def\zref{z^{\textrm{ref}}}
\def\bkN{{\mathbb{N}}}
\def\bkZ{{\bf{Z}}}
\def\intt{{\textrm{int}}}
\def\fnatt{\widetilde{\bf F}^{\textrm{nat}}}
\def\gnat{{\bf G}^{\textrm{nat}}}
\def\jnat{{\bf J}^{\textrm{nat}}}
\def\bfone{{\bf 1}}
\def\bfzero{{\bf 0}}
\def\bfU{{\bf U}}
\def\bfC{{\bf C}}
\def\Fscrhat{\widehat{\Fscr}}
\def\bkE{{\rm I\kern-.17em E}}
\def\bk1{{\rm 1\kern-.17em l}}
\def\bkD{{\rm I\kern-.17em D}}
\def\bkR{{\rm I\kern-.17em R}}
\def\bkP{{\rm I\kern-.17em P}}



\makeatletter
\newcommand{\pushright}[1]{\ifmeasuring@#1\else\omit\hfill$\displaystyle#1$\fi\ignorespaces}
\newcommand{\pushleft}[1]{\ifmeasuring@#1\else\omit$\displaystyle#1$\hfill\fi\ignorespaces}
\makeatother


\newcommand{\scg}{shared-constraint game\@\xspace} %%% i.e.,
\newcommand{\scgs}{shared-constraint games\@\xspace}
\newcommand{\Scgs}{Shared-constraint games\@\xspace}
\def\bkZ{{\bf{Z}}}
\def\b12{(\beta_1,\beta_2)}
\def\Epec{{\mathscr{E}}}
\def\bfA{{\bf A}}


\newenvironment{proofarg}[1][]{\noindent\hspace{2em}{\itshape Proof #1: }}{\hspace*{\fill}~\qed\par\endtrivlist\unskip}
%\newenvironment{proof}[1][]{{\noindent \bf Proof #1: }}{\hfill \qed \vspace{6pt}\\ }
\newenvironment{example}{{\noindent \bf Example}}{\hfill $\square$\hspace{-4.5pt}\vspace{6pt}}
\newcounter{example}
\renewcommand{\theexample}{\thesection.\arabic{example}}
\newenvironment{examplec}[1][]{\refstepcounter{example}
\par\medskip \noindent%
   \textbf{Example~\theexample. #1} \rmfamily}{\hfill $\square$   \hspace{-4.5pt} \vspace{6pt}}


% \newenvironment{keywords}{{\noindent \bf Keywords:}}{}
% \newenvironment{amsclass}{{\noindent \bf AMS Subject Classification:}}{}
\def\eig{{\rm eig}}
\def\eigmax{\eig_{\max}}
\def\eigmin{\eig_{\min}}
\def\ebox{\fbox{\begin{minipage}{1pt}
\hspace{1pt}\vspace{1pt}
\end{minipage}}}


\newcounter{remark}
\renewcommand{\theremark}{\thesection.\arabic{remark}}
\newenvironment{remarkc}[1][]{\refstepcounter{remark}
\noindent{\itshape Remark~\theremark. #1} \rmfamily}{\hspace*{\fill}~$\square$\vspace{0pt}}
\newenvironment{remark}{{\noindent \it Remark: }}{\hfill $\square$}
\def\bfe{{\bf e}}
\def\t{^\top}
\def\Bscr{\mathscr{B}}
\def\Xscr{\mathcal{X}}
\def\Yscr{\mathcal{Y}}
\def\indep{\perp\!\!\!\!\perp}
\def\Ebb{\mathbb{E}}
\newlength{\noteWidth}
\setlength{\noteWidth}{.75in}
\long\def\notes#1{\ifinner
{\tiny #1}
\else
\marginpar{\parbox[t]{\noteWidth}{\raggedright\tiny #1}}
\fi\typeout{#1}}

 \def\notes#1{\typeout{read notes: #1}} %uncomment for final version

\def\mi{^{-i}}
%\input{/home/ankur/latexfiles/algorithm2e.sty}
%  \def\proof{\penalty 25\noindent{\bf Proof.}\nobreak\hskip 1em\ignorespaces}
\def\noproof{\hspace{1em}\blackslug}
\def\mm#1{\hbox{$#1$}}                 % math mode
\def\mtx#1#2{\renewcommand{\arraystretch}{1.2}%
    \left(\! \begin{array}{#1}#2\end{array}\! \right)}
\newcommand{\mQ}{\mathcal{Q}}

\newcommand{\I}[1]{\mathbb{I}_{\{#1\}}}
\def\Ubf{\mathbf{U}}
\def\Qbf{\mathbf{Q}}
\def\Mbf{\mathbf{M}}

%\newcommand{\nr}{\{1,\hdots,2^{nR}\}}
\newcommand{\MM}{\{1,\hdots,M\}}
\newcommand{\Qxs}{Q_{x|s}}
\newcommand{\Qsy}{Q_{\widehat{s}|y}}
\newcommand{\QXS}{Q_{X|S}}
\newcommand{\QSY}{Q_{\widehat{S}|Y}}

\newcommand{\wi}[1]{\widehat{#1}}
\newcommand{\mbb}[1]{\mathbb{#1}}
\newcommand{\mcal}[1]{\mathcal{#1}}
\newcommand{\epde}{$\epsilon$-$\delta\;$}
\DeclareMathOperator{\image}{Im}

\newcommand{\expec}[1]{\Ebb\Big[#1\Big]}
\newcommand{\absol}[1]{\left|#1\right|}
\newcommand{\curly}[1]{\left\{#1\right\}}
\newcommand{\round}[1]{\left(#1\right)}
\newcommand{\sqre}[1]{\left[#1\right]}
%\newcommand{\expec}[1]{\Ebb\Big[#1\Big]}
\newcommand{\Typeq}{T_{1/q}^q}
\newcommand{\TypeKq}{T_{1/q}^{Kq}}
\newcommand{\TypeK}{T_{\Yscr}^{K}}
\newcommand{\capacity}{\Xi(\mathscr{U})}
\newcommand{\Gs}{G_{\mathsf{s}}}
\newcommand{\Gsym}{G_{\mathsf{s}}^{\mathsf{Sym}}}
%\newcommand{\Gsn}{G_{\mathsf{s}}^n}
\newcommand{\Gc}{G_{\mathsf{c}}}
%\newcommand{\Gcn}{G_{\mathsf{c}}^n}
%\newcommand{\ut}{\mathsf{U}}
\newcommand{\ut}{\mathscr{U}}
\newcommand{\utsym}{\mathscr{U}^{\mathsf{Sym}}}
\def\utbar{\bar \ut} %\skew{3.8}
\def\uttilde{\overset{\sim}{ \ut}}
%\def\uttilde{\tilde \ut}
\newcommand{\best}{\mathscr{B}}%{\Bscr}
\newcommand{\pbest}{\mathscr{S}}%{\Sscr}

\newcommand{\Ihat}{\widehat{I}}


\newcommand{\ie}{i.e.\@\xspace} %%% i.e.,
\newcommand{\eg}{e.g.\@\xspace} %%% e.g.,
\newcommand{\etal}{et al.\@\xspace} %%% e.g., Gill \etal (1986)
\newcommand{\Matlab}{\textsc{Matlab}\xspace}
% Scientific Notation (E)
\newcommand{\e}[2]{{\small e}$\scriptstyle#1$#2}
% Miscellaneous items
%\renewcommand{\vec}[1]{\ensuremath{\mathbf{#1}}}  %requires amsmath package
\newcommand{\abs}[1]{\ensuremath{|#1|}}
\newcommand{\Real}{\ensuremath{\mathbb{R}}}
\newcommand{\inv}{^{-1}}
\newcommand{\dom}{\mathrm{dom}}
\newcommand{\minimize}[1]{\displaystyle\minim_{#1}}
\newcommand{\minim}{\mathop{\hbox{\rm min}}}
\newcommand{\maximize}[1]{\displaystyle\maxim_{#1}}
\newcommand{\maxim}{\mathop{\hbox{\rm max}}}
%\DeclareMathOperator*{\st}{subject\;to}
\newcommand{\twonorm}[1]{\norm{#1}_2}
\def\onenorm#1{\norm{#1}_1}
\def\infnorm#1{\norm{#1}_{\infty}}
\def\disp{\displaystyle}
\def\m{\phantom-}
\def\EM#1{{\em#1\/}}
\def\ur{\mbox{\bf u}}
\def\FR{I\!\!F}
\def\var{{\rm Var}}
\def\Complex{I\mskip -11.3mu C}
\def\sp{\hspace{10pt}}
\def\subject{\hbox{\rm s.t.}}
\def\eval#1#2{\left.#2\right|_{\alpha=#1}}
\def\fnat{{\bf F}^{\rm nat}}
\def\fnor{{\bf F}^{\textrm{nor}}}
\def\tmat#1#2{(\, #1 \ \ #2 \,)}
\def\tmatt#1#2#3{(\, #1 \ \  #2 \ \  #3\,)}
\def\tmattt#1#2#3#4{(\, #1 \ \  #2 \ \  #3 \ \  #4\,)}
\def\OPT{{\rm OPT}}
\def\FEA{{\rm FEA}}

\def\rarr{\rightarrow}
\def\larr{\leftarrow}
\def\asn{\buildrel{n}\over\rightarrow}
\def\xref{x^{\rm ref}}
\def\Hcal{\mathcal{H}}
\def\Gcal{\mathcal{G}}
\def\Cbb{\mathbb{C}}
\def\Ebb{\mathbb{E}}
\def\Cbb{{\mathbb{C}}}
\def\Pbb{{\mathbb{P}}}
\def\Nbb{{\mathbb{N}}}
\def\Qbb{{\mathbb{Q}}}
\def\Xbb{{\mathbb{X}}}
\def\Ybb{{\mathbb{Y}}}
\def\Zbb{{\mathbb{Z}}}
\def\Ibb{{\mathbb{I}}}
\def\as{{\rm a.s.}}
\def\reff{^{\rm ref}}
\def\Qcal{\mathcal{Q}}

\def\Gbf{{\bf G}}
\def\Cbf{{\bf C}}
\def\Bbf{{\bf B}}
\def\Dbf{{\bf D}}
\def\Wbf{{\bf W}}
\def\Mbf{{\bf M}}
\def\Nbf{{\bf N}}
\def\Ibf{{\bf I}}
\def\dbf{{\bf d}}
\def\fbf{{\bf f}}
\def\limn{\ds \lim_{n \rightarrow \infty}}
\def\limt{\ds \lim_{t \rightarrow \infty}}
\def\sumn{\ds \sum_{n=0}^{\infty}}
\def\ast{\buildrel{t}\over\rightarrow}
\def\asn{\buildrel{n}\over\rightarrow}
\def\inp{\buildrel{p.}\over\rightarrow}

\def\ind{{\rm ind}}
\def\ext{{\rm ext}}

\def\Gr{{\rm graph}}
\def\bfm{{\bf m}}
\def\mat#1#2{(\; #1 \quad #2 \;)}
\def\matt#1#2#3{(\; #1 \quad #2 \quad #3\;)}
\def\mattt#1#2#3#4{(\; #1 \quad #2 \quad #3 \quad #4\;)}
\def\blackslug{\hbox{\hskip 1pt \vrule width 4pt height 6pt depth 1.5pt
  \hskip 1pt}}
\def\boxit#1{\vbox{\hrule\hbox{\vrule\hskip 3pt
  \vbox{\vskip 3pt #1 \vskip 3pt}\hskip 3pt\vrule}\hrule}}
\def\cross{\scriptscriptstyle\times}
\def\cond{\hbox{\rm cond}}
\def\det{\mathop{{\rm det}}}
\def\diag{\mathop{\hbox{\rm diag}}}
\def\dim{\mathop{\hbox{\rm dim}}}
\def\exp{\mathop{\hbox{\rm exp}}}
\def\In{\mathop{\hbox{\it In}\,}}
\def\boundary{\mathop{\hbox{\rm bnd}}}
\def\interior{\mathop{\hbox{\rm int}}}
\def\Null{\mathop{{\rm null}}}
\def\rank{\mathop{\hbox{\rm rank}}}
\def\op{\mathop{\hbox{\rm op}}}
\def\Range{\mathop{{\rm range}}}
%\def\range{\mathop{\hbox{\rm range}}}
\def\range{{\rm range}}
\def\sign{\mathop{\hbox{\rm sign}}}
\def\sgn{\mathop{\hbox{\rm sgn}}}
\def\Span{\mbox{\rm span}}
\def\intt{{\rm int}}
\def\trace{\mathop{\hbox{\rm trace}}}
\def\drop{^{\null}}
%\def\etal{{\em et al.\ }}  %%% e.g., Gill \etal (1986)
\def\float{fl}
\def\grad{\nabla}
\def\Hess{\nabla^2}
\def\half  {{\textstyle{1\over 2}}}
\def\third {{\textstyle{1\over 3}}}
\def\fourth{{\textstyle{1\over 4}}}
\def\sixth{{\textstyle{1\over 6}}}
\def\Nth{{\textstyle{1\over N}}}
\def\nth{{\textstyle{1\over n}}}
\def\mth{{\textstyle{1\over m}}}
\def\invsq{^{-2}}
\def\limk{\lim_{k\to\infty}}
\def\limN{\lim_{N\to\infty}}
\def\conv{\textrm{conv}\:}
\def\cl{\mbox{cl}\:}
\def\mystrut{\vrule height10.5pt depth5.5pt width0pt}
\def\mod#1{|#1|}
\def\modd#1{\biggl|#1\biggr|}
\def\Mod#1{\left|#1\right|}
\def\normm#1{\biggl\|#1\biggr\|}
\def\norm#1{\|#1\|}

\def\spose#1{\hbox to 0pt{#1\hss}}
\def\sub#1{^{\null}_{#1}}
\def\text #1{\hbox{\quad#1\quad}}
\def\textt#1{\hbox{\qquad#1\qquad}}
\def\T{^T\!}
\def\ntext#1{\noalign{\vskip 2pt\hbox{#1}\vskip 2pt}}
\def\enddef{{$\null$}}
\def\xint{{x_{\rm int}}}
\def\vhat{{\hat v}}
\def\fhat{{\hat f}}
\def\xhat{{\hat x}}
\def\zhat{{\hat x}}
\def\phihat{{\widehat \phi}}
\def\xihat{{\widehat \xi}}
\def\betahat{{\widehat \beta}}
\def\psihat{{\widehat \psi}}
\def\thetahat{{\hat \theta}}
\def\Escr{\mathcal{E}}
\def \dtilde{\tilde{d}}
%%% Fixed-size glue, only for math mode

\def\fnatt{\widetilde{\bf F}^{\rm nat}}
\def\gnat{{\bf G}^{\rm nat}}
\def\jnat{{\bf J}^{\rm nat}}

\def\pthinsp{\mskip  2   mu}    %           thin space
\def\pmedsp {\mskip  2.75mu}    %         medium space
\def\pthiksp{\mskip  3.5 mu}    %          thick space
\def\nthinsp{\mskip -2   mu}
\def\nmedsp {\mskip -2.75mu}
\def\nthiksp{\mskip -3.5 mu}

%%% Specially lowered subscripts (see \sub above)

\def\Asubk{A\sub{\nthinsp k}}
\def\Bsubk{B\sub{\nthinsp k}}
\def\Fsubk{F\sub{\nmedsp k}}
\def\Gsubk{G\sub{\nthinsp k}}
\def\Hsubk{H\sub{\nthinsp k}}
\def\HsubB{H\sub{\nthinsp \scriptscriptstyle{B}}}
\def\HsubS{H\sub{\nthinsp \scriptscriptstyle{S}}}
\def\HsubBS{H\sub{\nthinsp \scriptscriptstyle{BS}}}
\def\Hz{H\sub{\nthinsp \scriptscriptstyle{Z}}}
\def\Jsubk{J\sub{\nmedsp k}}
\def\Psubk{P\sub{\nthiksp k}}
\def\Qsubk{Q\sub{\nthinsp k}}
\def\Vsubk{V\sub{\nmedsp k}}
\def\Vsubone{V\sub{\nmedsp 1}}
\def\Vsubtwo{V\sub{\nmedsp 2}}
\def\Ysubk{Y\sub{\nmedsp k}}
\def\Wsubk{W\sub{\nthiksp k}}
\def\Zsubk{Z\sub{\nthinsp k}}
\def\Zsubi{Z\sub{\nthinsp i}}

%%% Subscripts

\def\subfr{_{\scriptscriptstyle{\it FR}}}
\def\subfx{_{\scriptscriptstyle{\it FX}}}
\def\fr{_{\scriptscriptstyle{\it FR}}}
\def\fx{_{\scriptscriptstyle{\it FX}}}
\def\submax{_{\max}}
\def\submin{_{\min}}
\def\subminus{_{\scriptscriptstyle -}}
\def\subplus {_{\scriptscriptstyle +}}

\def\Ass{_{\scriptscriptstyle A}}
\def\Bss{_{\scriptscriptstyle B}}
\def\BSss{_{\scriptscriptstyle BS}}
\def\Css{_{\scriptscriptstyle C}}
\def\Dss{_{\scriptscriptstyle D}}
\def\Ess{_{\scriptscriptstyle E}}
\def\Fss{_{\scriptscriptstyle F}}
\def\Gss{_{\scriptscriptstyle G}}
\def\Hss{_{\scriptscriptstyle H}}
\def\Iss{_{\scriptscriptstyle I}}
\def\w{_w}
\def\k{_k}
\def\kp#1{_{k+#1}}
\def\km#1{_{k-#1}}
\def\j{_j}
\def\jp#1{_{j+#1}}
\def\jm#1{_{j-#1}}
\def\K{_{\scriptscriptstyle K}}
\def\L{_{\scriptscriptstyle L}}
\def\M{_{\scriptscriptstyle M}}
\def\N{_{\scriptscriptstyle N}}
\def\subb{\B}
\def\subc{\C}
\def\subm{\M}
\def\subf{\F}
\def\subn{\N}
\def\subh{\H}
\def\subk{\K}
\def\subt{_{\scriptscriptstyle T}}
\def\subr{\R}
\def\subz{\Z}
\def\O{_{\scriptscriptstyle O}}
\def\Q{_{\scriptscriptstyle Q}}
\def\P{_{\scriptscriptstyle P}}
\def\R{_{\scriptscriptstyle R}}
\def\S{_{\scriptscriptstyle S}}
\def\U{_{\scriptscriptstyle U}}
\def\V{_{\scriptscriptstyle V}}
\def\W{_{\scriptscriptstyle W}}
\def\Y{_{\scriptscriptstyle Y}}
\def\Z{_{\scriptscriptstyle Z}}

%%% Superscript stars

\def\superstar{^{\raise 0.5pt\hbox{$\nthinsp *$}}}
\def\SUPERSTAR{^{\raise 0.5pt\hbox{$*$}}}
\def\shrink{\mskip -7mu}  %%% Less space after \xstar if followed by , or .

\def\alphastar{\alpha\superstar}
\def\lamstar  {\lambda\superstar}
\def\lamstarT {\lambda^{\raise 0.5pt\hbox{$\nthinsp *$}T}}
\def\lambdastar{\lamstar}
\def\nustar{\nu\superstar}
\def\mustar{\mu\superstar}
\def\pistar{\pi\superstar}
\def\s{^*}

\def\cstar{c\superstar}
\def\dstar{d\SUPERSTAR}
\def\fstar{f\SUPERSTAR}
\def\gstar{g\superstar}
\def\hstar{h\superstar}
\def\mstar{m\superstar}
\def\pstar{p\superstar}
\def\sstar{s\superstar}
\def\ustar{u\superstar}
\def\vstar{v\superstar}
\def\wstar{w\superstar}
\def\xstar{x\superstar}
\def\ystar{y\superstar}
\def\zstar{z\superstar}

\def\Astar{A\SUPERSTAR}
\def\Bstar{B\SUPERSTAR}
\def\Cstar{C\SUPERSTAR}
\def\Fstar{F\SUPERSTAR}
\def\Gstar{G\SUPERSTAR}
\def\Hstar{H\SUPERSTAR}
\def\Jstar{J\SUPERSTAR}
\def\Kstar{K\SUPERSTAR}
\def\Mstar{M\SUPERSTAR}
\def\Ustar{U\SUPERSTAR}
\def\Wstar{W\SUPERSTAR}
\def\Xstar{X\SUPERSTAR}
\def\Zstar{Z\SUPERSTAR}

%%% bars, hats, tildes  (\skew4 is the default)
\def\del{\partial}
\def\psibar{\bar \psi}
\def\alphabar{\bar \alpha}
\def\alphahat{\widehat \alpha}
\def\alphatilde{\skew3\tilde \alpha}
\def\betabar{\skew{2.8}\bar\beta}
\def\betahat{\skew{2.8}\hat\beta}
\def\betatilde{\skew{2.8}\tilde\beta}
\def\ftilde{\skew{2.8}\tilde f}
\def\delbar{\bar\delta}
\def\deltilde{\skew5\tilde \delta}
\def\deltabar{\delbar}
\def\deltatilde{\deltilde}
\def\lambar{\bar\lambda}
\def\etabar{\bar\eta}
\def\gammabar{\bar\gamma}
\def\lamhat{\skew{2.8}\hat \lambda}
\def\lambdabar{\lambar}
\def\lambdahat{\lamhat}
\def\mubar{\skew3\bar \mu}
\def\muhat{\skew3\hat \mu}
\def\mutilde{\skew3\tilde\mu}
\def\nubar{\skew3\bar\nu}
\def\nuhat{\skew3\hat\nu}
\def\nutilde{\skew3\tilde\nu}
\def\Mtilde{\skew3\tilde M}
\def\omegabar{\bar\omega}
\def\pibar{\bar\pi}
\def\pihat{\skew1\widehat \pi}
\def\sigmabar{\bar\sigma}
\def\sigmahat{\hat\sigma}
\def\rhobar{\bar\rho}
\def\rhohat{\widehat\rho}
\def\taubar{\bar\tau}
\def\tautilde{\tilde\tau}
\def\tauhat{\hat\tau}
\def\thetabar{\bar\theta}
\def\xibar{\skew4\bar\xi}
\def\UBbar{\skew4\overline {UB}}
\def\LBbar{\skew4\overline {LB}}
\def\xdot{\dot{x}}

%%%%%%%%%%%%%%%%%%%%%%%%%%%%%%%%%%%%%
%  Math italic and calligraphic fonts
%%%%%%%%%%%%%%%%%%%%%%%%%%%%%%%%%%%%%
\def\Deltait{{\mit \Delta}}
\def\Deltabar{{\bar \Delta}}
\def\Gammait{{\mit \Gamma}}
\def\Lambdait{{\mit \Lambda}}
\def\Sigmait{{\mit \Sigma}}
\def\Lambarit{\skew5\bar{\mit \Lambda}}
\def\Omegait{{\mit \Omega}}
\def\Omegaitbar{\skew5\bar{\mit \Omega}}
\def\Thetait{{\mit \Theta}}
\def\Piitbar{\skew5\bar{\mit \Pi}}
\def\Piit{{\mit \Pi}}
\def\Phiit{{\mit \Phi}}
\def\Ascr{{\cal A}}
\def\Bscr{{\cal B}}
\def\Fscr{{\cal F}}
\def\Dscr{{\cal D}}
\def\Iscr{{\cal I}}
\def\Jscr{{\cal J}}
\def\Hscr{{\cal H}}
\def\Lscr{{\cal L}}
\def\Mscr{{\cal M}}
\def\Tscr{{\cal T}}
\def\lscr{\ell}
\def\lscrbar{\ellbar}
\def\Oscr{{\cal O}}
\def\Pscr{{\cal P}}
\def\Qscr{{\cal Q}}
\def\Sscr{{\cal S}}
\def\Sscrhat{\widehat{\mathcal{S}}}
\def\Uscr{{\cal U}}
\def\Vscr{{\cal V}}
\def\Wscr{{\cal W}}
\def\Mscr{{\cal M}}
\def\Nscr{{\cal N}}
\def\Rscr{{\cal R}}
\def\Gscr{{\cal G}}
\def\Cscr{{\cal C}}
\def\Zscr{{\cal Z}}
\def\Xscr{{\cal X}}
\def\Yscr{{\cal Y}}
\def\Kscr{{\cal K}}
\def\Mhatscr{\cal{\widehat{M}}}
\def\abar{\skew3\bar a}
\def\ahat{\skew2\widehat a}
\def\atilde{\skew2\widetilde a}
\def\Abar{\skew7\bar A}
\def\Ahat{\widehat A}
\def\Atilde{\widetilde A}
\def\bbar{\skew3\bar b}
\def\bhat{\skew2\widehat b}
\def\btilde{\skew2\widetilde b}
\def\Bbar{\bar B}
\def\Bhat{\widehat B}
\def\cbar{\skew5\bar c}
\def\chat{\skew3\widehat c}
\def\ctilde{\widetilde c}
\def\Cbar{\bar C}
\def\Chat{\widehat C}
\def\Ctilde{\widetilde C}
\def\dbar{\bar d}
\def\dhat{\widehat d}
\def\dtilde{\widetilde d}
\def\Dbar{\bar D}
\def\Dhat{\widehat D}
\def\Dtilde{\widetilde D}
\def\ehat{\skew3\widehat e}
\def\ebar{\bar e}
\def\Ebar{\bar E}
\def\Ehat{\widehat E}
\def\fbar{\bar f}
\def\fhat{\widehat f}
\def\ftilde{\widetilde f}
\def\Fbar{\bar F}
\def\Fhat{\widehat F}
\def\gbar{\skew{4.3}\bar g}
\def\ghat{\skew{4.3}\widehat g}
\def\gtilde{\skew{4.5}\widetilde g}
\def\Gbar{\bar G}
\def\Ghat{\widehat G}
%\def\hbar{\skew{4.2}\bar h}
\def\hhat{\skew2\widehat h}
\def\htilde{\skew3\widetilde h}
\def\Hbar{\skew5\bar H}
\def\Hhat{\widehat H}
\def\Htilde{\widetilde H}
\def\Ibar{\skew5\bar I}
\def\Itilde{\widetilde I}
\def\Jbar{\skew6\bar J}
\def\Jhat{\widehat J}
\def\Jtilde{\widetilde J}
\def\kbar{\skew{4.4}\bar k}
\def\Khat{\widehat K}
\def\Kbar{\skew{4.4}\bar K}
\def\Ktilde{\widetilde K}
\def\ellbar{\bar \ell}
\def\lhat{\skew2\widehat l}
\def\lbar{\skew2\bar l}
\def\Lbar{\skew{4.3}\bar L}
\def\Lhat{\widehat L}
\def\Ltilde{\widetilde L}
\def\mbar{\skew2\bar m}
\def\mhat{\widehat m}
\def\Mbar{\skew{4.4}\bar M}
\def\Mhat{\widehat M}
\def\Mtilde{\widetilde M}
\def\Nbar{\skew{4.4}\bar N}
\def\Ntilde{\widetilde N}
\def\nbar{\skew2\bar n}
\def\pbar{\skew2\bar p}
%\def\phat{\skew2\widehat p}
\def\phat{\skew2\widehat p}
\def\ptilde{\skew2\widetilde p}
\def\Pbar{\skew5\bar P}
\def\phibar{\skew5\bar \phi}
\def\Phat{\widehat P}
\def\Ptilde{\skew5\widetilde P}
\def\qbar{\bar q}
\def\qhat{\skew2\widehat q}
\def\qtilde{\widetilde q}
\def\Qbar{\bar Q}
\def\Qhat{\widehat Q}
\def\Qtilde{\widetilde Q}
\def\rbar{\skew3\bar r}
\def\rhat{\skew3\widehat r}
\def\rtilde{\skew3\widetilde r}
\def\Rbar{\skew5\bar R}
\def\Rhat{\widehat R}
\def\Rtilde{\widetilde R}
\def\sbar{\bar s}
\def\shat{\widehat s}
\def\stilde{\widetilde s}
\def\Shat{\widehat S}
\def\Sbar{\skew2\bar S}
\def\tbar{\bar t}
\def\ttilde{\widetilde t}
\def\that{\widehat t}
\def\th{^{\rm th}}
\def\Tbar{\bar T}
\def\That{\widehat T}
\def\Ttilde{\widetilde T}
\def\ubar{\skew3\bar u}
\def\uhat{\skew3\widehat u}
\def\utilde{\skew3\widetilde u}
\def\Ubar{\skew2\bar U}
\def\Uhat{\widehat U}
\def\Utilde{\widetilde U}
\def\vbar{\skew3\bar v}
\def\vhat{\skew3\widehat v}
\def\vtilde{\skew3\widetilde v}
\def\Vbar{\skew2\bar V}
\def\Vhat{\widehat V}
\def\Vtilde{\widetilde V}
\def\UBtilde{\widetilde {UB}}
\def\Utilde{\widetilde {U}}
\def\LBtilde{\widetilde {LB}}
\def\Ltilde{\widetilde {L}}
\def\wbar{\skew3\bar w}
\def\what{\skew3\widehat w}
\def\wtilde{\skew3\widetilde w}
\def\Wbar{\skew3\bar W}
\def\What{\widehat W}
\def\Wtilde{\widetilde W}
\def\xbar{\skew{2.8}\bar x}
\def\xhat{\skew{2.8}\widehat x}
\def\xtilde{\skew3\widetilde x}
\def\Xhat{\widehat X}
\def\Xbar{\bar X}
\def\ybar{\skew3\bar y}
\def\yhat{\skew3\widehat y}
\def\ytilde{\skew3\widetilde y}
\def\Ybar{\skew2\bar Y}
\def\Yhat{\widehat Y}
\def\zbar{\skew{2.8}\bar z}
\def\zhat{\skew{2.8}\widehat z}
\def\ztilde{\skew{2.8}\widetilde z}
\def\Zbar{\skew5\bar Z}
\def\Zhat{\widehat Z}
\def\shatbar{\skew3 \bar \shat}
\def\Ztilde{\widetilde Z}
\def\MINOS{{\small MINOS}}
\def\NPSOL{{\small NPSOL}}
\def\QPSOL{{\small QPSOL}}
\def\LUSOL{{\small LUSOL}}
\def\LSSOL{{\small LSSOL}}

\def\Mbf{{\bf M}}
\def\Nbf{{\bf N}}
\def\Dbf{{\bf D}}
\def\Nbf{{\bf N}}
\def\Pbf{{\bf P}}
\def\vbf{{\bf v}}
\def\xbf{{\bf x}}
\def\shatbf{{\bf \shat}}
\def\ybf{{\bf y}}
\def\rbf{{\bf r}}
\def\sbf{{\bf s}}
\def \mbf{{\bf m}}
\def \kbf{{\bf k}}
\def\supp{{\rm supp}}
\def\aur{\;\textrm{and}\;}
\def\where{\;\textrm{where}\;}
\def\eef{\;\textrm{if}\;}
\def\ow{\;\textrm{otherwise}\;}
\def\del{\partial}
\def\non{\nonumber}
\def\SOL{{\rm SOL}}
\def\VI{{\rm VI}}
\def\AVI{{\rm AVI}}

\def\VIs{{\rm VIs}}
\def\QVI{{\rm QVI}}
\def\QVIs{{\rm QVIs}}
\def\LCP{{\rm LCP}}
\def\GNE{{\rm GNE}}
\def\VE{{\rm VE}}
\def\GNEs{{\rm GNEs}}
\def\VEs{{\rm VEs}}
\def\SYS{{\rm (SYS)}}
\def\bfI{{\bf I}}
\let\forallnew\forall
\renewcommand{\forall}{\forallnew\ }
\let\forall\forallnew
\newcommand{\mlmfg}{multi-leader multi-follower game\@\xspace}
\newcommand{\mlmfgs}{multi-leader multi-follower games\@\xspace}
\def\ds{\displaystyle}
		\def\bkE{{\rm I\kern-.17em E}}
		\def\bk1{{\rm 1\kern-.17em l}}
		\def\bkD{{\rm I\kern-.17em D}}
		\def\bkR{{\rm I\kern-.17em R}}
		\def\bkP{{\rm I\kern-.17em P}}
		\def\bkY{{\bf \kern-.17em Y}}
		\def\bkZ{{\bf \kern-.17em Z}}
		\def\bkC{{\bf  \kern-.17em C}}
%		\usepackage[ruled,vlined,boxed]{algorithm2e}
		%\setlength{\baselineskip}{24pt}
		%\setlength{\topmargin}     { 0.0in}    % One inch top margins
		%\setlength{\headsep}       {24.0pt}    % with Header
		%\setlength{\textheight}    { 8.5in}
		%\setlength{\textwidth}     { 6.5in}    % One inch side margins
		%\setlength{\oddsidemargin} { 0.0in}
		%\setlength{\evensidemargin}{ 0.0in}
		%\setlength{\marginparsep}  { 0.0pt}    % No Margin notes
		%\setlength{\marginparwidth}{ 0.0pt}
		%\pagestyle{headings}

%%%%%%%%%%%%%%%%%%%%%
%Colours
%%%%%%%%%%%%%%%%%%%%%
%\definecolor{violet}{rgb}{.5,0.0,.80}
%\def\violet{{ \color{violet} }}
%\definecolor{dgreen}{rgb}{0,0.5,0.0}


\def\shiftpar#1{\noindent \hangindent#1}
\newenvironment{indentpar}[1]%
{\begin{list}{}%
         {\setlength{\leftmargin}{#1}}%
         \item[]%
}
{\end{list}}
\def\Psf{\mathsf{P}}
\def\Tsf{\mathsf{T}}
\def\Esf{\mathsf{E}}
\def\Dsf{\mathsf{D}}
\def\Wsf{\mathsf{W}}
\def\Bsf{\mathsf{B}}
\def\Rsf{\mathsf{R}}
		% -- Environments --
		\def\bsp{\begin{split}}
		\def\beq{\begin{eqnarray}}
		\def\bal{\begin{align*}}
		\def\bc{\begin{center}}
		\def\be{\begin{enumerate}}
		\def\bi{\begin{itemize}}
		\def\bs{\begin{small}}
		\def\bS{\begin{slide}}
		\def\ec{\end{center}}
		\def\ee{\end{enumerate}}
		\def\ei{\end{itemize}}
		\def\es{\end{small}}
		\def\eS{\end{slide}}
		\def\eeq{\end{eqnarray}}
		\def\eal{\end{align*}}
		\def\esp{\end{split}}
		\def\qed{ \vrule height7.5pt width7.5pt depth0pt}  %width4.17pt depth0pt}
		\def\problemlarge#1#2#3#4{\fbox
		 {\begin{tabular*}{.88\textwidth}
			{@{}l@{\extracolsep{\fill}}l@{\extracolsep{6pt}}l@{\extracolsep{\fill}}c@{}}
				#1 & $\minimize{#2}$ & $#3$ & $ $ \\[5pt]
					 & $\subject\ $    & $#4$ & $ $
			\end{tabular*}}
			}
		\def\problemsuperlarge#1#2#3#4{\fbox
           {\tiny      {\begin{tabular*}{.95\textwidth}
                        {@{}l@{\extracolsep{\fill}}l@{\extracolsep{6pt}}l@{\extracolsep{\fill}}c@{}}
                              #1 & $\minimize{#2}$ &  $#3$ & $ $ \\[12pt]
                                         & $\subject\ $    & $#4$ & $ $
                        \end{tabular*}} }
                        }
	\def\maxproblemlarge#1#2#3#4{\fbox
		 {\begin{tabular*}{1.05\textwidth}
			{@{}l@{\extracolsep{\fill}}l@{\extracolsep{6pt}}l@{\extracolsep{\fill}}c@{}}
				#1 & $\maximize{#2}$ & $#3$ & $ $ \\[5pt]
					 & $\subject\ $    & $#4$ & $ $
			\end{tabular*}}
			}
	\def\maxproblemsmall#1#2#3#4{\fbox
		 {\begin{tabular*}{0.47\textwidth}
			{@{}l@{\extracolsep{\fill}}l@{\extracolsep{6pt}}l@{\extracolsep{\fill}}c@{}}
				#1 & $\maximize{#2}$ & $#3$ & $ $ \\[5pt]
					 & $\subject\ $    & $#4$ & $ $
			\end{tabular*}}
			}



\newcommand{\boxedeqnnew}[2]{\vspace{12pt}
 \noindent\framebox[\textwidth]{\begin{tabular*}{\textwidth}{p{.8\textwidth}p{.15\textwidth}}
                             \hspace{.05 \textwidth}    #1 & \raggedleft #2
                                \end{tabular*}
}\vspace{12pt}
 }

\newcommand{\boxedeqn}[3]{\vspace{-.2in}
\begin{center}
  \fbox{%
      \begin{minipage}{#2 \textwidth}%
      \begin{equation}#1 #3 \end{equation}%
      \end{minipage}%
    }%
    \end{center}
}
\newcommand{\boxedeqnsmall}[6]{\vspace{12pt}
 \noindent\framebox[#3 \textwidth]{\begin{tabular*}{#3 \textwidth}{p{#5 \textwidth }p{#6 \textwidth }}
                               \hspace{#4 \textwidth}  #1 & \raggedleft #2
                                \end{tabular*}
}\vspace{12pt}
 }


\newcommand{\anybox}[2]
{\fbox{
\begin{minipage}
\begin{tabular*}{#2 \texwidth}{l}
#1
\end{tabular*}
\end{minipage}
}}

		\def\problemsmall#1#2#3#4{\fbox
		 {\begin{tabular*}{0.47\textwidth}
			{@{}l@{\extracolsep{\fill}}l@{\extracolsep{6pt}}l@{\extracolsep{\fill}}c@{}}
				#1 & $\minimize{#2}$ & $#3$ & $ $ \\[5pt]
					  $\sub\ $ &    & $#4$ & $ $
			\end{tabular*}}
			}
\def\problemsmallsingcol#1#2#3#4{\fbox
		 {\begin{tabular*}{0.95\columnwidth}
			{@{}l@{\extracolsep{\fill}}l@{\extracolsep{6pt}}l@{\extracolsep{\fill}}c@{}}
				#1 & $\minimize{#2}$ & $#3$ & $ $ \\[5pt]
					  $\sub\ $ &    & $#4$ & $ $
			\end{tabular*}}
			}

                        \def\problemsmallsingcolwb#1#2#3#4{
		 \begin{tabular*}{0.95\columnwidth}
			{@{}l@{\extracolsep{\fill}}l@{\extracolsep{6pt}}l@{\extracolsep{\fill}}c@{}}
				#1 & $\minimize{#2}$ & $#3$ & $ $ \\[5pt]
					  $\sub\ $ &    & $#4$ & $ $
			\end{tabular*}
			}

		\def\problem#1#2#3#4{\fbox
		 {\begin{tabular*}{0.80\textwidth}
			{@{}l@{\extracolsep{\fill}}l@{\extracolsep{6pt}}l@{\extracolsep{\fill}}c@{}}
				#1 & $\minimize{#2}$ & $#3$ & $ $ \\[5pt]
					 & $\subject\ $    & $#4$ & $ $
			\end{tabular*}}
			}

		\def\minmaxproblem#1#2#3#4#5{\fbox
		 {\begin{tabular*}{0.80\textwidth}
			{@{}l@{\extracolsep{\fill}}l@{\extracolsep{6pt}}l@{\extracolsep{\fill}}c@{}}
				#1 & $\minimize{#2}$ $ \maximize{#3}$& $#4$ & $ $ \\[5pt]
					 & $\subject\ $    & $#5$ & $ $
			\end{tabular*}}
			}
		\def\minmaxproblemsmall#1#2#3#4#5{\fbox
		 {\begin{tabular*}{0.47\textwidth}
			{@{}l@{\extracolsep{\fill}}l@{\extracolsep{6pt}}l@{\extracolsep{\fill}}c@{}}
				#1 & $\minimize{#2}$ $ \maximize{#3}$& $#4$ & $ $ \\[5pt]
					 & $\subject\ $    & $#5$ & $ $
			\end{tabular*}}
			}
		\def\minmaxproblemcompactsmall#1#2#3#4#5{\fbox
		 {\begin{tabular*}{0.47\textwidth}
			{@{}l@{\extracolsep{\fill}}l@{\extracolsep{6pt}}l@{\extracolsep{\fill}}c@{}}
				#1 & $\minimize{#2}$ $ \maximize{#3}$& $#4$ & 			 $\ \subject$     $#5$
			\end{tabular*}}
			}


\def\maxproblem#1#2#3#4{\fbox
		 {\begin{tabular*}{0.80\textwidth}
			{@{}l@{\extracolsep{\fill}}l@{\extracolsep{6pt}}l@{\extracolsep{\fill}}c@{}}
				#1 & $\maximize{#2}$ & $#3$ & $ $ \\[5pt]
					 & $\subject\ $    & $#4$ & $ $
			\end{tabular*}}
			}

\def\maxxproblem#1#2#3#4{\fbox
		 {\begin{tabular*}{1.1\textwidth}
			{@{}l@{\extracolsep{\fill}}l@{\extracolsep{6pt}}l@{\extracolsep{\fill}}c@{}}
				#1 & $\maximize{#2}$ & $#3$ & $ $ \\[5pt]
					 & $\subject\ $    & $#4$ & $ $
			\end{tabular*}}
			}

\def\maxproblemsmallwb#1#2#3#4{
		 \begin{tabular*}{0.47\textwidth}
			{@{}l@{\extracolsep{\fill}}l@{\extracolsep{-4pt}}l@{\extracolsep{\fill}}c@{}}
				#1 &  & $\maximize{#2}$  $#3$ & $ $ \\[4pt]
					  $\sub \ $  &   & $#4$ &  $ $
			\end{tabular*}
			}


\def\maximproblemsmalla#1#2#3#4{\fbox
		 {\begin{tabular*}{0.49\textwidth}
			{@{}l@{\extracolsep{\fill}}l@{\extracolsep{1pt}}l@{\extracolsep{\fill}}c@{}}
				#1 &  & $\maximum{#2}$  $#3$ & $ $ \\[4pt]
					  $\sub \ $  &   & $#4$ &  $ $
			\end{tabular*}}
			}

\def\problemsmallawb#1#2#3#4{
		 \begin{tabular*}{0.47\textwidth}
			{@{}l@{\extracolsep{\fill}}l@{\extracolsep{-4pt}}l@{\extracolsep{\fill}}c@{}}
				#1 &  & $\minimize{#2}$  $#3$ & $ $ \\[4pt]
					  $\sub \ $  &   & $#4$ &  $ $
			\end{tabular*}
                      }
\def\problemsmallabb#1#2#3#4{\fbox{
		 \begin{tabular*}{0.56\textwidth}
			{@{}l@{\extracolsep{\fill}}l@{\extracolsep{-4pt}}l@{\extracolsep{\fill}}c@{}}
				#1 &  & \quad $= $ $\minimize{#2}$  $#3$&   \\[4pt]
                        $\sub \ $  &   & $#4$ &  $ $
			\end{tabular*}}
			}
\def\maxproblemsmallbb#1#2#3#4{\fbox{
		 \begin{tabular*}{0.56\textwidth}
			{@{}l@{\extracolsep{\fill}}l@{\extracolsep{-4pt}}l@{\extracolsep{\fill}}c@{}}
				#1 &  & \quad $ = $ $\maximize{#2}$  $#3$ & $ $ \\[4pt]
					  $\sub \ $  &   & $#4$ &  $ $
			\end{tabular*}
			}}

\def\unconmaxproblemsmall#1#2#3#4{\fbox
		 {\begin{tabular*}{0.47\textwidth}
			{@{}l@{\extracolsep{\fill}}l@{\extracolsep{6pt}}l@{\extracolsep{\fill}}c@{}}
				#1 & $\maximize{#2}$ & $#3$ & $ $ \\[5pt]
			\end{tabular*}}
			}

\def\unconmaxproblem#1#2#3#4{\fbox
		 {\begin{tabular*}{0.95\textwidth}
			{@{}l@{\extracolsep{\fill}}l@{\extracolsep{6pt}}l@{\extracolsep{\fill}}c@{}}
				#1 & $\maximize{#2}$ & $#3$ & $ $ \\[5pt]
			\end{tabular*}}
			}
	\def\problemcompactsmall#1#2#3#4{\fbox
		 {\begin{tabular*}{0.47\textwidth}
			{@{}l@{\extracolsep{\fill}}l@{\extracolsep{6pt}}l@{\extracolsep{\fill}}c@{}}
				#1 &$\minimize{#2}$\ $#3$ &{\rm s.t.} $#4 $
			\end{tabular*}}}
	\def\maxproblemcompactsmall#1#2#3#4{\fbox
		 {\begin{tabular*}{0.47\textwidth}
			{@{}l@{\extracolsep{\fill}}l@{\extracolsep{6pt}}l@{\extracolsep{\fill}}c@{}}
				#1 &$\maximize{#2}$\ $#3$ &{\rm s.t.}\ $#4 $
			\end{tabular*}}}


\def\unconproblem#1#2#3#4{\fbox
		 {\begin{tabular*}{0.95\textwidth}
			{@{}l@{\extracolsep{\fill}}l@{\extracolsep{6pt}}l@{\extracolsep{\fill}}c@{}}
				#1 & $\minimize{#2}$ & $#3$ & $ $ \\[5pt]
			\end{tabular*}}
			}
	\def\cpproblem#1#2#3#4{\fbox
		 {\begin{tabular*}{0.95\textwidth}
			{@{}l@{\extracolsep{\fill}}l@{\extracolsep{6pt}}l@{\extracolsep{\fill}}c@{}}
				#1 & & $#4 $
			\end{tabular*}}}
				\def\viproblem#1#2#3{\small\fbox
					 {\begin{tabular*}{0.47\textwidth}
						{@{}l@{\extracolsep{\fill}}l@{\extracolsep{2pt}}l@{\extracolsep{\fill}}c@{}}
							#1 &#2 & $#3 $
						\end{tabular*}}}

	\def\cp2problem#1#2#3#4{\fbox
		 {\begin{tabular*}{0.9\textwidth}
			{@{}l@{\extracolsep{\fill}}l@{\extracolsep{6pt}}l@{\extracolsep{\fill}}c@{}}
				#1 & & $#4 $
			\end{tabular*}}}
\def\z{\phantom 0}
\newcommand{\pmat}[1]{\begin{pmatrix} #1 \end{pmatrix}}
		\newcommand{\botline}{\vspace*{-1ex}\hrulefill}
		\newcommand{\clin}[2]{c\L(#1,#2)}
		\newcommand{\dlin}[2]{d\L(#1,#2)}
		\newcommand{\Fbox}[1]{\fbox{\quad#1\quad}}
		\newcommand{\mcol}[3]{\multicolumn{#1}{#2}{#3}}
		\newcommand{\assign}{\ensuremath{\leftarrow}}
		\newcommand{\ssub}[1]{\mbox{\scriptsize #1}}
		\def\bkE{{\rm I\kern-.17em E}}
		\def\bk1{{\rm 1\kern-.17em l}}
		\def\bkD{{\rm I\kern-.17em D}}
		\def\bkR{{\rm I\kern-.17em R}}
		\def\bkP{{\rm I\kern-.17em P}}

		\def\bkZ{{\bf{Z}}}
		\def\bkI{{\bf{I}}}

\def\boldp{\boldsymbol{p}}
\def\boldy{\boldsymbol{y}}


\newcommand {\beeq}[1]{\begin{equation}\label{#1}}
\newcommand {\eeeq}{\end{equation}}
\newcommand {\bea}{\begin{eqnarray}}
\newcommand {\eea}{\end{eqnarray}}
\newcommand{\mb}[1]{\mbox{\boldmath $#1$}}
\newcommand{\mbt}[1]{\mbox{\boldmath $\tilde{#1}$}}
\newcommand{\mbs}[1]{{\mbox{\boldmath \scriptsize{$#1$}}}}
\newcommand{\mbst}[1]{{\mbox{\boldmath \scriptsize{$\tilde{#1}$}}}}
\newcommand{\mbss}[1]{{\mbox{\boldmath \tiny{$#1$}}}}
\newcommand{\dpv}{\displaystyle \vspace{3pt}}
\def\texitem#1{\par\smallskip\noindent\hangindent 25pt
               \hbox to 25pt {\hss #1 ~}\ignorespaces}
\def\smskip{\par\vskip 7pt}
%\def\example{{\bf Example :\ }}
\def\st{\mbox{subject to}}

%%% Local Variables:
%%% mode: latex
%%% TeX-master: "report"
%%% End:

\def\bsp{\begin{split}}
		\def\beq{\begin{eqnarray}}
		\def\bal{\begin{align*}}
		\def\bc{\begin{center}}
		\def\be{\begin{enumerate}}
		\def\bi{\begin{itemize}}
		\def\bs{\begin{small}}
		\def\bS{\begin{slide}}
		\def\ec{\end{center}}
		\def\ee{\end{enumerate}}
		\def\ei{\end{itemize}}
		\def\es{\end{small}}
		\def\eS{\end{slide}}
		\def\eeq{\end{eqnarray}}
		\def\eal{\end{align*}}
		\def\esp{\end{split}}
		\def\qed{ \vrule height7.5pt width7.5pt depth0pt}  %width4.17pt depth0pt}
		\def\problemlarge#1#2#3#4{\fbox
		 {\begin{tabular*}{0.47\textwidth}
			{@{}l@{\extracolsep{\fill}}l@{\extracolsep{3pt}}l@{\extracolsep{\fill}}c@{}}
				#1 & $\minimize{#2}$ & $#3$ & $ $ \\[2pt]
					$\subject\ $  &    & $#4$ & $ $
			\end{tabular*}}
			}
		\def\problemsuperlarge#1#2#3#4{\fbox
           {\tiny      {\begin{tabular*}{0.485\textwidth}
                        {@{}l@{\extracolsep{\fill}}l@{\extracolsep{8pt}}l@{\extracolsep{\fill}}c@{}}
                              #1 & $\minimize{#2}$ &  $#3$ & $ $ \\[14pt]
                                         & $\subject\ $    & $#4$ & $ $
                        \end{tabular*}} }
                        }
                        \newenvironment{examplee}{{ \emph{Example:} }}{\hfill $\Box$ \vspace{6pt}}

                        \def\unconproblem#1#2#3{\fbox
		 {\begin{tabular*}{0.47\textwidth}
			{@{}l@{\extracolsep{\fill}}l@{\extracolsep{6pt}}l@{\extracolsep{\fill}}c@{}}
				#1 & $\minimize{#2}$ & $#3$ & $ $ \\[5pt]
			\end{tabular*}}
			}
                        %%%%%%%%%%%%%%%%%%%%%%%%%%%%%%%%%%%%%
%        macros.tex                  %
%%%%%%%%%%%%%%%%%%%%%%%%%%%%%%%%%%%%%%

\usepackage{amsmath, amssymb, xspace}
\usepackage{epsfig}
\usepackage{longtable}
\usepackage{color}
\usepackage{mathrsfs}
\usepackage{subfig}
%\newtheorem{theorem}{Theorem}[section]
%\newtheorem{assumption}{Assumption}[section]
%\newtheorem{lemma}{Lemma}[section]
%\newtheorem{conjecture}{Conjecture}[section]
%\newtheorem{claim}[theorem]{Claim}
%\newtheorem{proposition}[theorem]{Proposition}
%\newtheorem{prop}[theorem]{Proposition}
%\newtheorem{corollary}[theorem]{Corollary}
%\newtheorem{remark}[theorem]{Remark}
%\newtheorem{cor}[theorem]{Corollary}
%\newtheorem{result}[theorem]{Result}
%\newtheorem{definition}{Definition}[section]
%\newenvironment{proof}[1][]{{\noindent \emph {Proof} #1: }}{\hfill \qed \vspace{3pt}\\ }
\def\ext{{\rm ext}}
%\newenvironment{examplee}{{\noindent \bf Example}}{\hfill $\square$\hspace{-4.5pt}\vspace{6pt}}
%\newcommand{\ie}{i.e.\@\xspace}
\def\Cscr{{\cal C}}
\def\Nscr{{\cal N}}
\def\subject{\hbox{\rm subject to}}
\def\sub{\hbox{\rm s.t}}
\def\such{\hbox{\rm such that}}
\def\conv{{\rm conv}}
%\newcommand{\minim}{\mathop{\hbox{\rm min}}}
%\newcommand{\maximize}[1]{\displaystyle\maxim_{#1}}
%\newcommand{\maxim}{\mathop{\hbox{\rm maximize}}}
\newcommand{\maximum}{\mathop{\hbox{\rm max}}}
%\newcommand{\minimize}[1]{\displaystyle\minim_{#1}}
\def\unconproblem#1#2#3{\fbox
		 {\begin{tabular*}{0.47\textwidth}
			{@{}l@{\extracolsep{\fill}}l@{\extracolsep{6pt}}l@{\extracolsep{\fill}}c@{}}
				#1 & $\minimize{#2}$ & $#3$ & $ $ \\[5pt]
			\end{tabular*}}
			}
%			\def\problemsmall#1#2#3#4{\fbox
%		 {\begin{tabular*}{0.8 \textwidth}
%			{@{}l@{\extracolsep{\fill}}l@{\extracolsep{-3pt}}l@{\extracolsep{\fill}}c@{}}
%				#1 & & $\minimize{#2}$ $#3$ & $ $ \\[2pt]
%				$\sub\ $	 &$ $     & $#4$ & $ $
%			\end{tabular*}}
%			}
			\def\problemsmalla#1#2#3#4{\fbox
		 {\begin{tabular*}{0.47\textwidth}
			{@{}l@{\extracolsep{\fill}}l@{\extracolsep{-4pt}}l@{\extracolsep{\fill}}c@{}}
				#1 &  & $\minimize{#2}$  $#3$ & $ $ \\[4pt]
					  $\sub \ $  &   & $#4$ &  $ $
			\end{tabular*}}
			}
			\def\singlecolproblemsmalla#1#2#3#4{\fbox
		 {\begin{tabular*}{0.85\textwidth}
			{@{}l@{\extracolsep{\fill}}l@{\extracolsep{-4pt}}l@{\extracolsep{\fill}}c@{}}
				#1 &  & $\minimize{#2}$  $#3$ & $ $ \\[4pt]
					  $\sub \ $  &   & $#4$ &  $ $
			\end{tabular*}}
			}
			\def\maximproblemsmalla#1#2#3#4{\fbox
		 {\begin{tabular*}{0.49\textwidth}
			{@{}l@{\extracolsep{\fill}}l@{\extracolsep{-4pt}}l@{\extracolsep{\fill}}c@{}}
				#1 &  & $\maximum{#2}$  $#3$ & $ $ \\[4pt]
					  $\sub \ $  &   & $#4$ &  $ $
			\end{tabular*}}
			}
%\def\maximproblemsmalla#1#2#3#4{\fbox
%		 {\begin{tabular*}{0.49\textwidth}
%			{@{}l@{\extracolsep{\fill}}l@{\extracolsep{-4pt}}l@{\extracolsep{\fill}}c@{}}
%				#1 &  & $\maximum{#2}$  $#3$ & $ $ \\[4pt]
%					  $\sub \ $  &   & $#4$ &  $ $
%			\end{tabular*}}
%			}
%			\def\maxproblemlarge#1#2#3#4{\fbox
%		 {\begin{tabular*}{1.0\textwidth}
%			{@{}l@{\extracolsep{\fill}}l@{\extracolsep{-2pt}}l@{\extracolsep{\fill}}c@{}}
%				#1 & $\max{#2}$ & $#3$ & $ $ \\[3pt]
%					 & $\sub $    & $#4$ & $ $
%			\end{tabular*}}
%			}
	\def\maxproblemsmall#1#2#3#4{\fbox
		 {\begin{tabular*}{0.47\textwidth}
			{@{}l@{\extracolsep{\fill}}l@{\extracolsep{-4pt}}l@{\extracolsep{\fill}}c@{}}
				#1 &  & $\maximize {#2}$ $#3$ & $ $ \\[4pt]
					 $\sub \ $  &   & $#4$ & $ $
			\end{tabular*}}
                    }
	\def\maxproblemsmallsingcol#1#2#3#4{\fbox
		 {\begin{tabular*}{0.95\columnwidth}
			{@{}l@{\extracolsep{\fill}}l@{\extracolsep{-4pt}}l@{\extracolsep{\fill}}c@{}}
				#1 &  & $\maximize {#2}$ $#3$ & $ $ \\[4pt]
					 $\sub \ $  &   & $#4$ & $ $
			\end{tabular*}}
                    }
                    	\def\maxproblemsmallsingcolwb#1#2#3#4{
		 \begin{tabular*}{0.95\columnwidth}
			{@{}l@{\extracolsep{\fill}}l@{\extracolsep{-4pt}}l@{\extracolsep{\fill}}c@{}}
				#1 &  & $\maximize {#2}$ $#3$ & $ $ \\[4pt]
					 $\sub \ $  &   & $#4$ & $ $
			\end{tabular*}
			}
\def\singlecolmaxproblemsmall#1#2#3#4{\fbox
		 {\begin{tabular*}{0.85\textwidth}
			{@{}l@{\extracolsep{\fill}}l@{\extracolsep{-4pt}}l@{\extracolsep{\fill}}c@{}}
				#1 &  & $\maximize {#2}$ $#3$ & $ $ \\[4pt]
					 $\sub \ $  &   & $#4$ & $ $
			\end{tabular*}}
			}
%%%%%%%%%%%%%%%%%%%%%%%%%%%%%
% 	bb, bf and bk
%%%%%%%%%%%%%%%%%%%%%%%%%%%%%



% Add more packages as needed

% Theorem-like environments
\usepackage{amsthm} % For theorem-like environments

\theoremstyle{plain}
\newtheorem{thm}{Theorem}[section]
\newtheorem{lem}[thm]{Lemma}
\newtheorem{rem}{Remark}[section]

%\newtheorem{prop}{Proposition}
%newtheorem{cor}{Corollary}

\theoremstyle{definition}
\newtheorem{defn}{Definition}[section]

%\newtheorem{eg}{Example}[section]

\newtheorem{exg}{Example}
\newtheorem{case}{Case}
\newtheorem{casecase}{Scenario}
\newtheorem{claims}{Claim}
\renewcommand{\theclaims}{\arabic{claims}}
\setcounter{claims}{0}
\newtheorem{step}{Step}
\newtheorem{subcase}{Subcase}
\newtheorem{subcases}{Subcase}
\newtheorem*{notn}{Notation}
\newtheorem*{locnotn}{Local Notation}
\newtheorem{apdxproof}{Proof}[section]  % Numbered by section, e.g., Proof 3.1
% Redefine it to add QED box
\let\oldapdxproof\apdxproof
\let\endoldapdxproof\endapdxproof
\RenewEnviron{apdxproof}{
  \begin{oldapdxproof}
    \BODY
  \qed  % Adds the end-of-proof box
  \end{oldapdxproof}
}

\AtBeginEnvironment{apdxproof}{\setcounter{case}{0}}

% \makeatletter
% \pretocmd{\proof}{\setcounter{case}{0}}{}{}
% \makeatother

% Define more theorem environments as needed


% Abbreviations
%\newcommand{\etal}{et al.}
%\newcommand{\ie}{i.e.,}
\newcommand{\wrt}{w.r.t.}
\newcommand{\etc}{etc.}
\newcommand{\cf}{cf.}
\newcommand{\iid}{i.i.d.}
% Add more abbreviation as needed


% Math symbols and operators
\newcommand{\vect}[1]{\mathbf{#1}} % Vector notation (bold)
\newcommand{\mat}[1]{\mathbf{#1}}  % Matrix notation (bold)
\newcommand{\R}{\mathbb{R}} % Real numbers
\newcommand{\N}{\mathbb{N}} % Natural numbers
\newcommand{\C}{\mathbb{C}} % Complex numbers
\newcommand{\Z}{\mathbb{Z}} % Integers
\newcommand{\Q}{\mathbb{Q}} % Rational numbers
\newcommand{\eps}{\epsilon} % epsilon
\newcommand{\dlt}{\delta} % delta
\newcommand{\argmax}{\operatornamewithlimits{argmax}} % argmax operator
\newcommand{\argmin}{\operatornamewithlimits{argmin}} % argmin operator
% Add more operator and symbols as needed


% Probability theory
\newcommand{\prob}{\mathbb{P}} % Probability notation
\newcommand{\E}{\mathbb{E}} % Expectation operator
\newcommand{\event}{\mathcal{E}} % Event notation
% Add more from probability theory as needed


% Learning theory
\newcommand{\loss}{\mathcal{L}} % Loss function
\newcommand{\data}{\mathcal{D}} % Data set
\newcommand{\hypothesis}{\mathcal{F}} % Hypothesis class
% Add more from learning theory as needed


% Game theory
\newcommand{\neqb}{\text{Nash Equilibrium}} % Nash Equilibrium
\newcommand{\seqb}{\text{Stackelberg Equilibrium}} % Stackelberg Equilibrium
% Add more from game theory as needed


% Scripts
\newcommand{\sr}[1]{\mathscr{#1}} % Script letters
\newcommand{\cg}[1]{\mathcal{#1}} % Calligraphic letterss
\newcommand{\bb}[1]{\mathbb{#1}} % Blackboard bold letters
\newcommand{\fk}[1]{\mathfrak{#1}} % Fraktur letters
\newcommand{\mrm}[1]{\mathrm{#1}} % Roman letters
\newcommand{\msf}[1]{\mathsf{#1}} % Serif


% Add more scripts as needed

% Topic specific packages and shortcuts


\newcommand{\h}{\mathbf{h}}
\newcommand{\xm}{\mathbf{i_{\epsilon}}}
\newcommand{\zm}{\mathbf{k_{\epsilon}}}
\newcommand{\xp}{\mathbf{j_{\epsilon}}}
\newcommand{\zp}{\mathbf{l_{\epsilon}}}
\newcommand{\zb}{\mathbf{z}}
\newcommand{\xxm}{\mathbf{i}}
\newcommand{\xxp}{\mathbf{j}}
\newcommand{\um}{\mathbf{u_-}}
\newcommand{\up}{\mathbf{u_+}}
\newcommand{\uu}{\mathbf{u}}
\newcommand{\vv}{\mathbf{v}}
\newcommand{\rt}{\tilde{r}}
\newcommand{\xt}{\tilde{x}}
\newcommand{\ph}{\hat{p}}
\newcommand{\stt}{\tilde{s}}
\newcommand{\srt}{s_{\tilde{r}}}
\newcommand{\srr}{s_r}
\newcommand{\pt}{\tilde{p}}
\newcommand{\hht}{\tilde{h}}
\newcommand{\hh}{\hat{h}}
\newcommand{\bbb}{\mathbf{b}}
\newcommand{\qt}{\tilde{q}}
\newcommand{\minus}{\mathbf{-}}
\newcommand{\plus}{\mathbf{+}}
\newcommand{\rs}{r^{\star}}
\newcommand{\srs}{s_{r^{\star}}}
\newcommand{\ib}{\mathbf{h^{-}}}
\newcommand{\jb}{\mathbf{h^{+}}}
\newcommand{\kb}{\mathbf{k}}
\newcommand{\lb}{\mathbf{g}}
\newcommand{\ibb}{\mathbf{i}}
\newcommand{\ab}{\mathbf{a}}
\newcommand{\jbb}{\mathbf{j}}
\newcommand{\kbb}{\mathbf{k}}
\newcommand{\lbb}{\mathbf{l}}
\newcommand{\hon}{\mathbf{h_1}^{(-)}}
\newcommand{\hop}{\mathbf{h_1}^{(+)}}
\newcommand{\htn}{\mathbf{h_2}^{(-)}}
\newcommand{\htp}{\mathbf{h_2}^{(+)}}
\newcommand{\rcst}{\tilde{r}\circ s_{\tilde{r}}}
\newcommand{\tlbd}{\tilde{\mathbf{g}'}}
\newcommand{\tlb}{\tilde{\mathbf{g}}}
\newcommand{\qh}{\hat{q}}
\newcommand{\qtt}{\tilde{\tilde{q}}}
\newcommand{\rh}{\hat{r}}
\newcommand{\bqo}{\bar{q}_1}
\newcommand{\bqt}{\bar{q}_2}
\newcommand{\ibt}{\widehat{\mathbf{h^{(-)}}}}
\newcommand{\jbt}{\widehat{\mathbf{h^{(+)}}}}
\newcommand{\hds}{\hat{\mathfrak{D}}(S)}
\newcommand{\senpayoff}{\mathbb{U}}

\newcommand{\ibeps}{\mathbf{h}_{\eps}^-}
\newcommand{\jbeps}{\mathbf{h}_{\eps}^+}

\title{Strategic Classification for Non-uniform Preferences using Penalty Labels and Randomisation
}

\author{
    \IEEEauthorblockN{Manish Kumar Singh\IEEEauthorrefmark{1}, 
                      Ankur A. Kulkarni\IEEEauthorrefmark{1}}
    \IEEEauthorblockA{\\ \IEEEauthorrefmark{1}Center for Systems and Control Engineering, \\
                      Indian Institute of Technology Bombay, \\
                      Mumbai, Maharashtra 400076, India \\
                      Email: {manish@sc.iitb.ac.in, kulkarni.ankur@iitb.ac.in}}
}

%%%%%%%%%%%%%%%%%%%%%%%%%%%%%%%%%%%%%
%        macros.tex                  %
%%%%%%%%%%%%%%%%%%%%%%%%%%%%%%%%%%%%%%

%From PhD thesis
\usepackage{amsmath, amssymb, bbm, xspace}
\usepackage{epsfig}
\usepackage{longtable}
\usepackage{color}
\usepackage{mathrsfs}
\usepackage{comment}
\usepackage{ifthen}
\newboolean{showcomments}
\setboolean{showcomments}{true}
%\newcommand{\aak}[1]{  \ifthenelse{\boolean{showcomments}}
%{ \textcolor{blue}{(AAK says:  #1)}} {}  }
\def\sthat{\ {\rm such\ that}\ }
%\usepackage{fancyhdr}
% \usepackage[T1]{fontenc}
% \usepackage{arev}
%\usepackage{helvet}
\usepackage{courier}
%\usepackage{courier}
% \usepackage[T1]{fontenc}
% \usepackage[sc]{mathpazo}
% \linespread{1.05}         % Palatino needs more leading (space between lines)
\usepackage{xcolor}
\usepackage{todonotes}
\usepackage[framemethod=tikz]{mdframed}
\usepackage{lineno}
\newmdenv[leftmargin=\dimexpr-0.4em, innerleftmargin=0.5em,
rightmargin=\dimexpr-0.4em, innerrightmargin=0.5em,
linewidth=2pt,linecolor=red, topline=false, bottomline=false,
innertopmargin=0pt,innerbottommargin=0pt,skipbelow=0pt,skipabove=0pt,%
]{notex}

\newenvironment{note}%
{\vskip\dimexpr\dp\strutbox-\prevdepth\relax\notex\strut\ignorespaces}%
{\xdef\notetpd{\the\prevdepth}\endnotex\vskip-\notetpd\relax}

\let\oldtodo\todo

\makeatletter%
\DeclareDocumentCommand{\todo}{ O{} +g +d<> }{%
		\setlength{\marginparwidth}{1.5cm}%
	%
	\IfNoValueTF{#2}{\relax}{%
		\oldtodo[caption={#2},size=\scriptsize,#1]{\renewcommand{\baselinestretch}{0.8}\selectfont\sffamily#2\par}%
	}%
	\IfNoValueTF{#3}{\relax}{%
		\IfNoValueTF{#2}{% when parmark but no todo
			\begin{note}%
				\begin{internallinenumbers}%
					\indent%
					#3%
				\end{internallinenumbers}%
			\end{note}%
		}{% when parmark and todo
			\vspace{-0\baselineskip}%
			\begin{note}%
				\begin{internallinenumbers}%
					\indent%
					#3%
				\end{internallinenumbers}%
			\end{note}%
		}%
	}%
}%
\makeatother
\usepackage{soul}

\newcommand{\hlc}[2][yellow]{{%
		\colorlet{foo}{#1}%
		\sethlcolor{foo}\hl{#2}}%
}

\newtheorem{conjecture}{Conjecture}[section]

\newcommand{\removetodo}[2]{\todo[color=pink]{\textbf{delete:} ``#1'' #2}\hlc[pink]{#1}}
\newcommand{\inserttodo}[1]{\todo[color=green!40]{\textbf{insert:} #1}}


\newcommand{\hltodoy}[2]{\todo[color=yellow!40]{#2}\hl{#1} }

\newcommand{\hltodoc}[3]{\todo[color=#3!40]{#2}\hlc[#3]{#1} }

\newcommand{\hltodo}[2]{\todo[color=orange!40]{#2}\hlc[orange!40]{#1} }
\newcommand{\replacetodo}[2]{\todo[color=pink!40]{\textbf{replace with:}``#2'' }\hl{#1} }

\usepackage{marginnote}
\newcommand{\todol}[1]{{%
		\let\marginpar\marginnote
		\reversemarginpar
		\renewcommand{\baselinestretch}{0.8}%
		\todo{#1}}}

\newcommand{\inserttodol}[1]{{%
		\let\marginpar\marginnote
		\reversemarginpar
		\renewcommand{\baselinestretch}{0.8}%
		\inserttodo{#1}}}

\newcommand{\removetodol}[2]{{%
		\let\marginpar\marginnote
		\reversemarginpar
		\renewcommand{\baselinestretch}{0.8}%
		\removetodo{#1}{#2}}}

\newcommand{\hltodol}[2]{{%
		\let\marginpar\marginnote
		\reversemarginpar
		\renewcommand{\baselinestretch}{0.8}%
		\hltodo{#1}{#2}}}

\newcommand{\replacetodol}[2]{{%
		\let\marginpar\marginnote
		\reversemarginpar
		\renewcommand{\baselinestretch}{0.8}%
		\replacetodo{#1}{#2}}}

\newcommand{\hltodoyl}[2]{{%
		\let\marginpar\marginnote
		\reversemarginpar
		\renewcommand{\baselinestretch}{0.8}%
		\hltodoy{#1}{#2}}}

\newcommand{\hltodocl}[3]{{		\let\marginpar\marginnote
		\reversemarginpar
		\renewcommand{\baselinestretch}{0.8}%
		\hltodoc{#1}{#2}{#3}}}

\newcommand{\tododone}[1]{\todo[nolist,color=gray!20,bordercolor=gray!60]{#1}}
% \usepackage{titlesec}
% \titlelabel{\qed \thetitle \quad}
% \titleformat*{\section}{\itshape}

% This abstract environment is a little narrower
% than the default.  Note the capital A...
% \newenvironment{Abstract}
%   {\begin{center}
%     \textbf{Abstract} \\
%     \medskip
%   \begin{minipage}{4.8in}}
%   {\end{minipage}
% \end{center}}



\newtheorem{theorem}{Theorem}[section]
\newtheorem{assumption}{Assumption}[section]
\newtheorem{lemma}[theorem]{Lemma}
\newtheorem{claim}[theorem]{Claim}
\newtheorem{proposition}[theorem]{Proposition}
\newtheorem{prop}[theorem]{Proposition}
\newtheorem{corollary}[theorem]{Corollary}
\newtheorem{cor}[theorem]{Corollary}
\newtheorem{result}[theorem]{Result}
\newtheorem{definition}{Definition}[section]

\newcommand{\comb}[2]{\prescript{#1}{}C_{#2}}

%\newcounter{theorem}[chapter]
% \renewcommand{\thetheorem}{\thechapter.\arabic{theorem}}
%
% \renewenvironment{theorem}{\refstepcounter{theorem} {\ \\[2ex]}{\bf \large \textsf{Theorem \thetheorem}} \newline \it}{}

\def\bkE{{\rm I\kern-.17em E}}
\def\bk1{{\rm 1\kern-.17em l}}
\def\bkD{{\rm I\kern-.17em D}}
\def\bkR{{\rm I\kern-.17em R}}
\def\bkP{{\rm I\kern-.17em P}}
\def\Cov{{\rm Cov}}
\def\Cbb{{\mathbb{C}}}
\def\Fbb{{\mathbb{F}}}
\def\C{{\mathbb{C}}}
\def\Abb{{\mathbb{A}}}
\def\Mbb{{\mathbb{M}}}
\def\Bbb{{\mathbb{B}}}
\def\Vbb{{\mathbb{V}}}
\def\Mhatbb{{\mathbb{\widehat{M}}}}
\def\fnat{{\bf F}^{\textrm{nat}}}
\def\xref{x^{\textrm{ref}}}
\def\zref{z^{\textrm{ref}}}
\def\bkN{{\mathbb{N}}}
\def\bkZ{{\bf{Z}}}
\def\intt{{\textrm{int}}}
\def\fnatt{\widetilde{\bf F}^{\textrm{nat}}}
\def\gnat{{\bf G}^{\textrm{nat}}}
\def\jnat{{\bf J}^{\textrm{nat}}}
\def\bfone{{\bf 1}}
\def\bfzero{{\bf 0}}
\def\bfU{{\bf U}}
\def\bfC{{\bf C}}
\def\Fscrhat{\widehat{\Fscr}}
\def\bkE{{\rm I\kern-.17em E}}
\def\bk1{{\rm 1\kern-.17em l}}
\def\bkD{{\rm I\kern-.17em D}}
\def\bkR{{\rm I\kern-.17em R}}
\def\bkP{{\rm I\kern-.17em P}}



\makeatletter
\newcommand{\pushright}[1]{\ifmeasuring@#1\else\omit\hfill$\displaystyle#1$\fi\ignorespaces}
\newcommand{\pushleft}[1]{\ifmeasuring@#1\else\omit$\displaystyle#1$\hfill\fi\ignorespaces}
\makeatother


\newcommand{\scg}{shared-constraint game\@\xspace} %%% i.e.,
\newcommand{\scgs}{shared-constraint games\@\xspace}
\newcommand{\Scgs}{Shared-constraint games\@\xspace}
\def\bkZ{{\bf{Z}}}
\def\b12{(\beta_1,\beta_2)}
\def\Epec{{\mathscr{E}}}
\def\bfA{{\bf A}}


\newenvironment{proofarg}[1][]{\noindent\hspace{2em}{\itshape Proof #1: }}{\hspace*{\fill}~\qed\par\endtrivlist\unskip}
%\newenvironment{proof}[1][]{{\noindent \bf Proof #1: }}{\hfill \qed \vspace{6pt}\\ }
\newenvironment{example}{{\noindent \bf Example}}{\hfill $\square$\hspace{-4.5pt}\vspace{6pt}}
\newcounter{example}
\renewcommand{\theexample}{\thesection.\arabic{example}}
\newenvironment{examplec}[1][]{\refstepcounter{example}
\par\medskip \noindent%
   \textbf{Example~\theexample. #1} \rmfamily}{\hfill $\square$   \hspace{-4.5pt} \vspace{6pt}}


% \newenvironment{keywords}{{\noindent \bf Keywords:}}{}
% \newenvironment{amsclass}{{\noindent \bf AMS Subject Classification:}}{}
\def\eig{{\rm eig}}
\def\eigmax{\eig_{\max}}
\def\eigmin{\eig_{\min}}
\def\ebox{\fbox{\begin{minipage}{1pt}
\hspace{1pt}\vspace{1pt}
\end{minipage}}}


\newcounter{remark}
\renewcommand{\theremark}{\thesection.\arabic{remark}}
\newenvironment{remarkc}[1][]{\refstepcounter{remark}
\noindent{\itshape Remark~\theremark. #1} \rmfamily}{\hspace*{\fill}~$\square$\vspace{0pt}}
\newenvironment{remark}{{\noindent \it Remark: }}{\hfill $\square$}
\def\bfe{{\bf e}}
\def\t{^\top}
\def\Bscr{\mathscr{B}}
\def\Xscr{\mathcal{X}}
\def\Yscr{\mathcal{Y}}
\def\indep{\perp\!\!\!\!\perp}
\def\Ebb{\mathbb{E}}
\newlength{\noteWidth}
\setlength{\noteWidth}{.75in}
\long\def\notes#1{\ifinner
{\tiny #1}
\else
\marginpar{\parbox[t]{\noteWidth}{\raggedright\tiny #1}}
\fi\typeout{#1}}

 \def\notes#1{\typeout{read notes: #1}} %uncomment for final version

\def\mi{^{-i}}
%\input{/home/ankur/latexfiles/algorithm2e.sty}
%  \def\proof{\penalty 25\noindent{\bf Proof.}\nobreak\hskip 1em\ignorespaces}
\def\noproof{\hspace{1em}\blackslug}
\def\mm#1{\hbox{$#1$}}                 % math mode
\def\mtx#1#2{\renewcommand{\arraystretch}{1.2}%
    \left(\! \begin{array}{#1}#2\end{array}\! \right)}
\newcommand{\mQ}{\mathcal{Q}}

\newcommand{\I}[1]{\mathbb{I}_{\{#1\}}}
\def\Ubf{\mathbf{U}}
\def\Qbf{\mathbf{Q}}
\def\Mbf{\mathbf{M}}

%\newcommand{\nr}{\{1,\hdots,2^{nR}\}}
\newcommand{\MM}{\{1,\hdots,M\}}
\newcommand{\Qxs}{Q_{x|s}}
\newcommand{\Qsy}{Q_{\widehat{s}|y}}
\newcommand{\QXS}{Q_{X|S}}
\newcommand{\QSY}{Q_{\widehat{S}|Y}}

\newcommand{\wi}[1]{\widehat{#1}}
\newcommand{\mbb}[1]{\mathbb{#1}}
\newcommand{\mcal}[1]{\mathcal{#1}}
\newcommand{\epde}{$\epsilon$-$\delta\;$}
\DeclareMathOperator{\image}{Im}

\newcommand{\expec}[1]{\Ebb\Big[#1\Big]}
\newcommand{\absol}[1]{\left|#1\right|}
\newcommand{\curly}[1]{\left\{#1\right\}}
\newcommand{\round}[1]{\left(#1\right)}
\newcommand{\sqre}[1]{\left[#1\right]}
%\newcommand{\expec}[1]{\Ebb\Big[#1\Big]}
\newcommand{\Typeq}{T_{1/q}^q}
\newcommand{\TypeKq}{T_{1/q}^{Kq}}
\newcommand{\TypeK}{T_{\Yscr}^{K}}
\newcommand{\capacity}{\Xi(\mathscr{U})}
\newcommand{\Gs}{G_{\mathsf{s}}}
\newcommand{\Gsym}{G_{\mathsf{s}}^{\mathsf{Sym}}}
%\newcommand{\Gsn}{G_{\mathsf{s}}^n}
\newcommand{\Gc}{G_{\mathsf{c}}}
%\newcommand{\Gcn}{G_{\mathsf{c}}^n}
%\newcommand{\ut}{\mathsf{U}}
\newcommand{\ut}{\mathscr{U}}
\newcommand{\utsym}{\mathscr{U}^{\mathsf{Sym}}}
\def\utbar{\bar \ut} %\skew{3.8}
\def\uttilde{\overset{\sim}{ \ut}}
%\def\uttilde{\tilde \ut}
\newcommand{\best}{\mathscr{B}}%{\Bscr}
\newcommand{\pbest}{\mathscr{S}}%{\Sscr}

\newcommand{\Ihat}{\widehat{I}}


\newcommand{\ie}{i.e.\@\xspace} %%% i.e.,
\newcommand{\eg}{e.g.\@\xspace} %%% e.g.,
\newcommand{\etal}{et al.\@\xspace} %%% e.g., Gill \etal (1986)
\newcommand{\Matlab}{\textsc{Matlab}\xspace}
% Scientific Notation (E)
\newcommand{\e}[2]{{\small e}$\scriptstyle#1$#2}
% Miscellaneous items
%\renewcommand{\vec}[1]{\ensuremath{\mathbf{#1}}}  %requires amsmath package
\newcommand{\abs}[1]{\ensuremath{|#1|}}
\newcommand{\Real}{\ensuremath{\mathbb{R}}}
\newcommand{\inv}{^{-1}}
\newcommand{\dom}{\mathrm{dom}}
\newcommand{\minimize}[1]{\displaystyle\minim_{#1}}
\newcommand{\minim}{\mathop{\hbox{\rm min}}}
\newcommand{\maximize}[1]{\displaystyle\maxim_{#1}}
\newcommand{\maxim}{\mathop{\hbox{\rm max}}}
%\DeclareMathOperator*{\st}{subject\;to}
\newcommand{\twonorm}[1]{\norm{#1}_2}
\def\onenorm#1{\norm{#1}_1}
\def\infnorm#1{\norm{#1}_{\infty}}
\def\disp{\displaystyle}
\def\m{\phantom-}
\def\EM#1{{\em#1\/}}
\def\ur{\mbox{\bf u}}
\def\FR{I\!\!F}
\def\var{{\rm Var}}
\def\Complex{I\mskip -11.3mu C}
\def\sp{\hspace{10pt}}
\def\subject{\hbox{\rm s.t.}}
\def\eval#1#2{\left.#2\right|_{\alpha=#1}}
\def\fnat{{\bf F}^{\rm nat}}
\def\fnor{{\bf F}^{\textrm{nor}}}
\def\tmat#1#2{(\, #1 \ \ #2 \,)}
\def\tmatt#1#2#3{(\, #1 \ \  #2 \ \  #3\,)}
\def\tmattt#1#2#3#4{(\, #1 \ \  #2 \ \  #3 \ \  #4\,)}
\def\OPT{{\rm OPT}}
\def\FEA{{\rm FEA}}

\def\rarr{\rightarrow}
\def\larr{\leftarrow}
\def\asn{\buildrel{n}\over\rightarrow}
\def\xref{x^{\rm ref}}
\def\Hcal{\mathcal{H}}
\def\Gcal{\mathcal{G}}
\def\Cbb{\mathbb{C}}
\def\Ebb{\mathbb{E}}
\def\Cbb{{\mathbb{C}}}
\def\Pbb{{\mathbb{P}}}
\def\Nbb{{\mathbb{N}}}
\def\Qbb{{\mathbb{Q}}}
\def\Xbb{{\mathbb{X}}}
\def\Ybb{{\mathbb{Y}}}
\def\Zbb{{\mathbb{Z}}}
\def\Ibb{{\mathbb{I}}}
\def\as{{\rm a.s.}}
\def\reff{^{\rm ref}}
\def\Qcal{\mathcal{Q}}

\def\Gbf{{\bf G}}
\def\Cbf{{\bf C}}
\def\Bbf{{\bf B}}
\def\Dbf{{\bf D}}
\def\Wbf{{\bf W}}
\def\Mbf{{\bf M}}
\def\Nbf{{\bf N}}
\def\Ibf{{\bf I}}
\def\dbf{{\bf d}}
\def\fbf{{\bf f}}
\def\limn{\ds \lim_{n \rightarrow \infty}}
\def\limt{\ds \lim_{t \rightarrow \infty}}
\def\sumn{\ds \sum_{n=0}^{\infty}}
\def\ast{\buildrel{t}\over\rightarrow}
\def\asn{\buildrel{n}\over\rightarrow}
\def\inp{\buildrel{p.}\over\rightarrow}

\def\ind{{\rm ind}}
\def\ext{{\rm ext}}

\def\Gr{{\rm graph}}
\def\bfm{{\bf m}}
\def\mat#1#2{(\; #1 \quad #2 \;)}
\def\matt#1#2#3{(\; #1 \quad #2 \quad #3\;)}
\def\mattt#1#2#3#4{(\; #1 \quad #2 \quad #3 \quad #4\;)}
\def\blackslug{\hbox{\hskip 1pt \vrule width 4pt height 6pt depth 1.5pt
  \hskip 1pt}}
\def\boxit#1{\vbox{\hrule\hbox{\vrule\hskip 3pt
  \vbox{\vskip 3pt #1 \vskip 3pt}\hskip 3pt\vrule}\hrule}}
\def\cross{\scriptscriptstyle\times}
\def\cond{\hbox{\rm cond}}
\def\det{\mathop{{\rm det}}}
\def\diag{\mathop{\hbox{\rm diag}}}
\def\dim{\mathop{\hbox{\rm dim}}}
\def\exp{\mathop{\hbox{\rm exp}}}
\def\In{\mathop{\hbox{\it In}\,}}
\def\boundary{\mathop{\hbox{\rm bnd}}}
\def\interior{\mathop{\hbox{\rm int}}}
\def\Null{\mathop{{\rm null}}}
\def\rank{\mathop{\hbox{\rm rank}}}
\def\op{\mathop{\hbox{\rm op}}}
\def\Range{\mathop{{\rm range}}}
%\def\range{\mathop{\hbox{\rm range}}}
\def\range{{\rm range}}
\def\sign{\mathop{\hbox{\rm sign}}}
\def\sgn{\mathop{\hbox{\rm sgn}}}
\def\Span{\mbox{\rm span}}
\def\intt{{\rm int}}
\def\trace{\mathop{\hbox{\rm trace}}}
\def\drop{^{\null}}
%\def\etal{{\em et al.\ }}  %%% e.g., Gill \etal (1986)
\def\float{fl}
\def\grad{\nabla}
\def\Hess{\nabla^2}
\def\half  {{\textstyle{1\over 2}}}
\def\third {{\textstyle{1\over 3}}}
\def\fourth{{\textstyle{1\over 4}}}
\def\sixth{{\textstyle{1\over 6}}}
\def\Nth{{\textstyle{1\over N}}}
\def\nth{{\textstyle{1\over n}}}
\def\mth{{\textstyle{1\over m}}}
\def\invsq{^{-2}}
\def\limk{\lim_{k\to\infty}}
\def\limN{\lim_{N\to\infty}}
\def\conv{\textrm{conv}\:}
\def\cl{\mbox{cl}\:}
\def\mystrut{\vrule height10.5pt depth5.5pt width0pt}
\def\mod#1{|#1|}
\def\modd#1{\biggl|#1\biggr|}
\def\Mod#1{\left|#1\right|}
\def\normm#1{\biggl\|#1\biggr\|}
\def\norm#1{\|#1\|}

\def\spose#1{\hbox to 0pt{#1\hss}}
\def\sub#1{^{\null}_{#1}}
\def\text #1{\hbox{\quad#1\quad}}
\def\textt#1{\hbox{\qquad#1\qquad}}
\def\T{^T\!}
\def\ntext#1{\noalign{\vskip 2pt\hbox{#1}\vskip 2pt}}
\def\enddef{{$\null$}}
\def\xint{{x_{\rm int}}}
\def\vhat{{\hat v}}
\def\fhat{{\hat f}}
\def\xhat{{\hat x}}
\def\zhat{{\hat x}}
\def\phihat{{\widehat \phi}}
\def\xihat{{\widehat \xi}}
\def\betahat{{\widehat \beta}}
\def\psihat{{\widehat \psi}}
\def\thetahat{{\hat \theta}}
\def\Escr{\mathcal{E}}
\def \dtilde{\tilde{d}}
%%% Fixed-size glue, only for math mode

\def\fnatt{\widetilde{\bf F}^{\rm nat}}
\def\gnat{{\bf G}^{\rm nat}}
\def\jnat{{\bf J}^{\rm nat}}

\def\pthinsp{\mskip  2   mu}    %           thin space
\def\pmedsp {\mskip  2.75mu}    %         medium space
\def\pthiksp{\mskip  3.5 mu}    %          thick space
\def\nthinsp{\mskip -2   mu}
\def\nmedsp {\mskip -2.75mu}
\def\nthiksp{\mskip -3.5 mu}

%%% Specially lowered subscripts (see \sub above)

\def\Asubk{A\sub{\nthinsp k}}
\def\Bsubk{B\sub{\nthinsp k}}
\def\Fsubk{F\sub{\nmedsp k}}
\def\Gsubk{G\sub{\nthinsp k}}
\def\Hsubk{H\sub{\nthinsp k}}
\def\HsubB{H\sub{\nthinsp \scriptscriptstyle{B}}}
\def\HsubS{H\sub{\nthinsp \scriptscriptstyle{S}}}
\def\HsubBS{H\sub{\nthinsp \scriptscriptstyle{BS}}}
\def\Hz{H\sub{\nthinsp \scriptscriptstyle{Z}}}
\def\Jsubk{J\sub{\nmedsp k}}
\def\Psubk{P\sub{\nthiksp k}}
\def\Qsubk{Q\sub{\nthinsp k}}
\def\Vsubk{V\sub{\nmedsp k}}
\def\Vsubone{V\sub{\nmedsp 1}}
\def\Vsubtwo{V\sub{\nmedsp 2}}
\def\Ysubk{Y\sub{\nmedsp k}}
\def\Wsubk{W\sub{\nthiksp k}}
\def\Zsubk{Z\sub{\nthinsp k}}
\def\Zsubi{Z\sub{\nthinsp i}}

%%% Subscripts

\def\subfr{_{\scriptscriptstyle{\it FR}}}
\def\subfx{_{\scriptscriptstyle{\it FX}}}
\def\fr{_{\scriptscriptstyle{\it FR}}}
\def\fx{_{\scriptscriptstyle{\it FX}}}
\def\submax{_{\max}}
\def\submin{_{\min}}
\def\subminus{_{\scriptscriptstyle -}}
\def\subplus {_{\scriptscriptstyle +}}

\def\Ass{_{\scriptscriptstyle A}}
\def\Bss{_{\scriptscriptstyle B}}
\def\BSss{_{\scriptscriptstyle BS}}
\def\Css{_{\scriptscriptstyle C}}
\def\Dss{_{\scriptscriptstyle D}}
\def\Ess{_{\scriptscriptstyle E}}
\def\Fss{_{\scriptscriptstyle F}}
\def\Gss{_{\scriptscriptstyle G}}
\def\Hss{_{\scriptscriptstyle H}}
\def\Iss{_{\scriptscriptstyle I}}
\def\w{_w}
\def\k{_k}
\def\kp#1{_{k+#1}}
\def\km#1{_{k-#1}}
\def\j{_j}
\def\jp#1{_{j+#1}}
\def\jm#1{_{j-#1}}
\def\K{_{\scriptscriptstyle K}}
\def\L{_{\scriptscriptstyle L}}
\def\M{_{\scriptscriptstyle M}}
\def\N{_{\scriptscriptstyle N}}
\def\subb{\B}
\def\subc{\C}
\def\subm{\M}
\def\subf{\F}
\def\subn{\N}
\def\subh{\H}
\def\subk{\K}
\def\subt{_{\scriptscriptstyle T}}
\def\subr{\R}
\def\subz{\Z}
\def\O{_{\scriptscriptstyle O}}
\def\Q{_{\scriptscriptstyle Q}}
\def\P{_{\scriptscriptstyle P}}
\def\R{_{\scriptscriptstyle R}}
\def\S{_{\scriptscriptstyle S}}
\def\U{_{\scriptscriptstyle U}}
\def\V{_{\scriptscriptstyle V}}
\def\W{_{\scriptscriptstyle W}}
\def\Y{_{\scriptscriptstyle Y}}
\def\Z{_{\scriptscriptstyle Z}}

%%% Superscript stars

\def\superstar{^{\raise 0.5pt\hbox{$\nthinsp *$}}}
\def\SUPERSTAR{^{\raise 0.5pt\hbox{$*$}}}
\def\shrink{\mskip -7mu}  %%% Less space after \xstar if followed by , or .

\def\alphastar{\alpha\superstar}
\def\lamstar  {\lambda\superstar}
\def\lamstarT {\lambda^{\raise 0.5pt\hbox{$\nthinsp *$}T}}
\def\lambdastar{\lamstar}
\def\nustar{\nu\superstar}
\def\mustar{\mu\superstar}
\def\pistar{\pi\superstar}
\def\s{^*}

\def\cstar{c\superstar}
\def\dstar{d\SUPERSTAR}
\def\fstar{f\SUPERSTAR}
\def\gstar{g\superstar}
\def\hstar{h\superstar}
\def\mstar{m\superstar}
\def\pstar{p\superstar}
\def\sstar{s\superstar}
\def\ustar{u\superstar}
\def\vstar{v\superstar}
\def\wstar{w\superstar}
\def\xstar{x\superstar}
\def\ystar{y\superstar}
\def\zstar{z\superstar}

\def\Astar{A\SUPERSTAR}
\def\Bstar{B\SUPERSTAR}
\def\Cstar{C\SUPERSTAR}
\def\Fstar{F\SUPERSTAR}
\def\Gstar{G\SUPERSTAR}
\def\Hstar{H\SUPERSTAR}
\def\Jstar{J\SUPERSTAR}
\def\Kstar{K\SUPERSTAR}
\def\Mstar{M\SUPERSTAR}
\def\Ustar{U\SUPERSTAR}
\def\Wstar{W\SUPERSTAR}
\def\Xstar{X\SUPERSTAR}
\def\Zstar{Z\SUPERSTAR}

%%% bars, hats, tildes  (\skew4 is the default)
\def\del{\partial}
\def\psibar{\bar \psi}
\def\alphabar{\bar \alpha}
\def\alphahat{\widehat \alpha}
\def\alphatilde{\skew3\tilde \alpha}
\def\betabar{\skew{2.8}\bar\beta}
\def\betahat{\skew{2.8}\hat\beta}
\def\betatilde{\skew{2.8}\tilde\beta}
\def\ftilde{\skew{2.8}\tilde f}
\def\delbar{\bar\delta}
\def\deltilde{\skew5\tilde \delta}
\def\deltabar{\delbar}
\def\deltatilde{\deltilde}
\def\lambar{\bar\lambda}
\def\etabar{\bar\eta}
\def\gammabar{\bar\gamma}
\def\lamhat{\skew{2.8}\hat \lambda}
\def\lambdabar{\lambar}
\def\lambdahat{\lamhat}
\def\mubar{\skew3\bar \mu}
\def\muhat{\skew3\hat \mu}
\def\mutilde{\skew3\tilde\mu}
\def\nubar{\skew3\bar\nu}
\def\nuhat{\skew3\hat\nu}
\def\nutilde{\skew3\tilde\nu}
\def\Mtilde{\skew3\tilde M}
\def\omegabar{\bar\omega}
\def\pibar{\bar\pi}
\def\pihat{\skew1\widehat \pi}
\def\sigmabar{\bar\sigma}
\def\sigmahat{\hat\sigma}
\def\rhobar{\bar\rho}
\def\rhohat{\widehat\rho}
\def\taubar{\bar\tau}
\def\tautilde{\tilde\tau}
\def\tauhat{\hat\tau}
\def\thetabar{\bar\theta}
\def\xibar{\skew4\bar\xi}
\def\UBbar{\skew4\overline {UB}}
\def\LBbar{\skew4\overline {LB}}
\def\xdot{\dot{x}}

%%%%%%%%%%%%%%%%%%%%%%%%%%%%%%%%%%%%%
%  Math italic and calligraphic fonts
%%%%%%%%%%%%%%%%%%%%%%%%%%%%%%%%%%%%%
\def\Deltait{{\mit \Delta}}
\def\Deltabar{{\bar \Delta}}
\def\Gammait{{\mit \Gamma}}
\def\Lambdait{{\mit \Lambda}}
\def\Sigmait{{\mit \Sigma}}
\def\Lambarit{\skew5\bar{\mit \Lambda}}
\def\Omegait{{\mit \Omega}}
\def\Omegaitbar{\skew5\bar{\mit \Omega}}
\def\Thetait{{\mit \Theta}}
\def\Piitbar{\skew5\bar{\mit \Pi}}
\def\Piit{{\mit \Pi}}
\def\Phiit{{\mit \Phi}}
\def\Ascr{{\cal A}}
\def\Bscr{{\cal B}}
\def\Fscr{{\cal F}}
\def\Dscr{{\cal D}}
\def\Iscr{{\cal I}}
\def\Jscr{{\cal J}}
\def\Hscr{{\cal H}}
\def\Lscr{{\cal L}}
\def\Mscr{{\cal M}}
\def\Tscr{{\cal T}}
\def\lscr{\ell}
\def\lscrbar{\ellbar}
\def\Oscr{{\cal O}}
\def\Pscr{{\cal P}}
\def\Qscr{{\cal Q}}
\def\Sscr{{\cal S}}
\def\Sscrhat{\widehat{\mathcal{S}}}
\def\Uscr{{\cal U}}
\def\Vscr{{\cal V}}
\def\Wscr{{\cal W}}
\def\Mscr{{\cal M}}
\def\Nscr{{\cal N}}
\def\Rscr{{\cal R}}
\def\Gscr{{\cal G}}
\def\Cscr{{\cal C}}
\def\Zscr{{\cal Z}}
\def\Xscr{{\cal X}}
\def\Yscr{{\cal Y}}
\def\Kscr{{\cal K}}
\def\Mhatscr{\cal{\widehat{M}}}
\def\abar{\skew3\bar a}
\def\ahat{\skew2\widehat a}
\def\atilde{\skew2\widetilde a}
\def\Abar{\skew7\bar A}
\def\Ahat{\widehat A}
\def\Atilde{\widetilde A}
\def\bbar{\skew3\bar b}
\def\bhat{\skew2\widehat b}
\def\btilde{\skew2\widetilde b}
\def\Bbar{\bar B}
\def\Bhat{\widehat B}
\def\cbar{\skew5\bar c}
\def\chat{\skew3\widehat c}
\def\ctilde{\widetilde c}
\def\Cbar{\bar C}
\def\Chat{\widehat C}
\def\Ctilde{\widetilde C}
\def\dbar{\bar d}
\def\dhat{\widehat d}
\def\dtilde{\widetilde d}
\def\Dbar{\bar D}
\def\Dhat{\widehat D}
\def\Dtilde{\widetilde D}
\def\ehat{\skew3\widehat e}
\def\ebar{\bar e}
\def\Ebar{\bar E}
\def\Ehat{\widehat E}
\def\fbar{\bar f}
\def\fhat{\widehat f}
\def\ftilde{\widetilde f}
\def\Fbar{\bar F}
\def\Fhat{\widehat F}
\def\gbar{\skew{4.3}\bar g}
\def\ghat{\skew{4.3}\widehat g}
\def\gtilde{\skew{4.5}\widetilde g}
\def\Gbar{\bar G}
\def\Ghat{\widehat G}
%\def\hbar{\skew{4.2}\bar h}
\def\hhat{\skew2\widehat h}
\def\htilde{\skew3\widetilde h}
\def\Hbar{\skew5\bar H}
\def\Hhat{\widehat H}
\def\Htilde{\widetilde H}
\def\Ibar{\skew5\bar I}
\def\Itilde{\widetilde I}
\def\Jbar{\skew6\bar J}
\def\Jhat{\widehat J}
\def\Jtilde{\widetilde J}
\def\kbar{\skew{4.4}\bar k}
\def\Khat{\widehat K}
\def\Kbar{\skew{4.4}\bar K}
\def\Ktilde{\widetilde K}
\def\ellbar{\bar \ell}
\def\lhat{\skew2\widehat l}
\def\lbar{\skew2\bar l}
\def\Lbar{\skew{4.3}\bar L}
\def\Lhat{\widehat L}
\def\Ltilde{\widetilde L}
\def\mbar{\skew2\bar m}
\def\mhat{\widehat m}
\def\Mbar{\skew{4.4}\bar M}
\def\Mhat{\widehat M}
\def\Mtilde{\widetilde M}
\def\Nbar{\skew{4.4}\bar N}
\def\Ntilde{\widetilde N}
\def\nbar{\skew2\bar n}
\def\pbar{\skew2\bar p}
%\def\phat{\skew2\widehat p}
\def\phat{\skew2\widehat p}
\def\ptilde{\skew2\widetilde p}
\def\Pbar{\skew5\bar P}
\def\phibar{\skew5\bar \phi}
\def\Phat{\widehat P}
\def\Ptilde{\skew5\widetilde P}
\def\qbar{\bar q}
\def\qhat{\skew2\widehat q}
\def\qtilde{\widetilde q}
\def\Qbar{\bar Q}
\def\Qhat{\widehat Q}
\def\Qtilde{\widetilde Q}
\def\rbar{\skew3\bar r}
\def\rhat{\skew3\widehat r}
\def\rtilde{\skew3\widetilde r}
\def\Rbar{\skew5\bar R}
\def\Rhat{\widehat R}
\def\Rtilde{\widetilde R}
\def\sbar{\bar s}
\def\shat{\widehat s}
\def\stilde{\widetilde s}
\def\Shat{\widehat S}
\def\Sbar{\skew2\bar S}
\def\tbar{\bar t}
\def\ttilde{\widetilde t}
\def\that{\widehat t}
\def\th{^{\rm th}}
\def\Tbar{\bar T}
\def\That{\widehat T}
\def\Ttilde{\widetilde T}
\def\ubar{\skew3\bar u}
\def\uhat{\skew3\widehat u}
\def\utilde{\skew3\widetilde u}
\def\Ubar{\skew2\bar U}
\def\Uhat{\widehat U}
\def\Utilde{\widetilde U}
\def\vbar{\skew3\bar v}
\def\vhat{\skew3\widehat v}
\def\vtilde{\skew3\widetilde v}
\def\Vbar{\skew2\bar V}
\def\Vhat{\widehat V}
\def\Vtilde{\widetilde V}
\def\UBtilde{\widetilde {UB}}
\def\Utilde{\widetilde {U}}
\def\LBtilde{\widetilde {LB}}
\def\Ltilde{\widetilde {L}}
\def\wbar{\skew3\bar w}
\def\what{\skew3\widehat w}
\def\wtilde{\skew3\widetilde w}
\def\Wbar{\skew3\bar W}
\def\What{\widehat W}
\def\Wtilde{\widetilde W}
\def\xbar{\skew{2.8}\bar x}
\def\xhat{\skew{2.8}\widehat x}
\def\xtilde{\skew3\widetilde x}
\def\Xhat{\widehat X}
\def\Xbar{\bar X}
\def\ybar{\skew3\bar y}
\def\yhat{\skew3\widehat y}
\def\ytilde{\skew3\widetilde y}
\def\Ybar{\skew2\bar Y}
\def\Yhat{\widehat Y}
\def\zbar{\skew{2.8}\bar z}
\def\zhat{\skew{2.8}\widehat z}
\def\ztilde{\skew{2.8}\widetilde z}
\def\Zbar{\skew5\bar Z}
\def\Zhat{\widehat Z}
\def\shatbar{\skew3 \bar \shat}
\def\Ztilde{\widetilde Z}
\def\MINOS{{\small MINOS}}
\def\NPSOL{{\small NPSOL}}
\def\QPSOL{{\small QPSOL}}
\def\LUSOL{{\small LUSOL}}
\def\LSSOL{{\small LSSOL}}

\def\Mbf{{\bf M}}
\def\Nbf{{\bf N}}
\def\Dbf{{\bf D}}
\def\Nbf{{\bf N}}
\def\Pbf{{\bf P}}
\def\vbf{{\bf v}}
\def\xbf{{\bf x}}
\def\shatbf{{\bf \shat}}
\def\ybf{{\bf y}}
\def\rbf{{\bf r}}
\def\sbf{{\bf s}}
\def \mbf{{\bf m}}
\def \kbf{{\bf k}}
\def\supp{{\rm supp}}
\def\aur{\;\textrm{and}\;}
\def\where{\;\textrm{where}\;}
\def\eef{\;\textrm{if}\;}
\def\ow{\;\textrm{otherwise}\;}
\def\del{\partial}
\def\non{\nonumber}
\def\SOL{{\rm SOL}}
\def\VI{{\rm VI}}
\def\AVI{{\rm AVI}}

\def\VIs{{\rm VIs}}
\def\QVI{{\rm QVI}}
\def\QVIs{{\rm QVIs}}
\def\LCP{{\rm LCP}}
\def\GNE{{\rm GNE}}
\def\VE{{\rm VE}}
\def\GNEs{{\rm GNEs}}
\def\VEs{{\rm VEs}}
\def\SYS{{\rm (SYS)}}
\def\bfI{{\bf I}}
\let\forallnew\forall
\renewcommand{\forall}{\forallnew\ }
\let\forall\forallnew
\newcommand{\mlmfg}{multi-leader multi-follower game\@\xspace}
\newcommand{\mlmfgs}{multi-leader multi-follower games\@\xspace}
\def\ds{\displaystyle}
		\def\bkE{{\rm I\kern-.17em E}}
		\def\bk1{{\rm 1\kern-.17em l}}
		\def\bkD{{\rm I\kern-.17em D}}
		\def\bkR{{\rm I\kern-.17em R}}
		\def\bkP{{\rm I\kern-.17em P}}
		\def\bkY{{\bf \kern-.17em Y}}
		\def\bkZ{{\bf \kern-.17em Z}}
		\def\bkC{{\bf  \kern-.17em C}}
%		\usepackage[ruled,vlined,boxed]{algorithm2e}
		%\setlength{\baselineskip}{24pt}
		%\setlength{\topmargin}     { 0.0in}    % One inch top margins
		%\setlength{\headsep}       {24.0pt}    % with Header
		%\setlength{\textheight}    { 8.5in}
		%\setlength{\textwidth}     { 6.5in}    % One inch side margins
		%\setlength{\oddsidemargin} { 0.0in}
		%\setlength{\evensidemargin}{ 0.0in}
		%\setlength{\marginparsep}  { 0.0pt}    % No Margin notes
		%\setlength{\marginparwidth}{ 0.0pt}
		%\pagestyle{headings}

%%%%%%%%%%%%%%%%%%%%%
%Colours
%%%%%%%%%%%%%%%%%%%%%
%\definecolor{violet}{rgb}{.5,0.0,.80}
%\def\violet{{ \color{violet} }}
%\definecolor{dgreen}{rgb}{0,0.5,0.0}


\def\shiftpar#1{\noindent \hangindent#1}
\newenvironment{indentpar}[1]%
{\begin{list}{}%
         {\setlength{\leftmargin}{#1}}%
         \item[]%
}
{\end{list}}
\def\Psf{\mathsf{P}}
\def\Tsf{\mathsf{T}}
\def\Esf{\mathsf{E}}
\def\Dsf{\mathsf{D}}
\def\Wsf{\mathsf{W}}
\def\Bsf{\mathsf{B}}
\def\Rsf{\mathsf{R}}
		% -- Environments --
		\def\bsp{\begin{split}}
		\def\beq{\begin{eqnarray}}
		\def\bal{\begin{align*}}
		\def\bc{\begin{center}}
		\def\be{\begin{enumerate}}
		\def\bi{\begin{itemize}}
		\def\bs{\begin{small}}
		\def\bS{\begin{slide}}
		\def\ec{\end{center}}
		\def\ee{\end{enumerate}}
		\def\ei{\end{itemize}}
		\def\es{\end{small}}
		\def\eS{\end{slide}}
		\def\eeq{\end{eqnarray}}
		\def\eal{\end{align*}}
		\def\esp{\end{split}}
		\def\qed{ \vrule height7.5pt width7.5pt depth0pt}  %width4.17pt depth0pt}
		\def\problemlarge#1#2#3#4{\fbox
		 {\begin{tabular*}{.88\textwidth}
			{@{}l@{\extracolsep{\fill}}l@{\extracolsep{6pt}}l@{\extracolsep{\fill}}c@{}}
				#1 & $\minimize{#2}$ & $#3$ & $ $ \\[5pt]
					 & $\subject\ $    & $#4$ & $ $
			\end{tabular*}}
			}
		\def\problemsuperlarge#1#2#3#4{\fbox
           {\tiny      {\begin{tabular*}{.95\textwidth}
                        {@{}l@{\extracolsep{\fill}}l@{\extracolsep{6pt}}l@{\extracolsep{\fill}}c@{}}
                              #1 & $\minimize{#2}$ &  $#3$ & $ $ \\[12pt]
                                         & $\subject\ $    & $#4$ & $ $
                        \end{tabular*}} }
                        }
	\def\maxproblemlarge#1#2#3#4{\fbox
		 {\begin{tabular*}{1.05\textwidth}
			{@{}l@{\extracolsep{\fill}}l@{\extracolsep{6pt}}l@{\extracolsep{\fill}}c@{}}
				#1 & $\maximize{#2}$ & $#3$ & $ $ \\[5pt]
					 & $\subject\ $    & $#4$ & $ $
			\end{tabular*}}
			}
	\def\maxproblemsmall#1#2#3#4{\fbox
		 {\begin{tabular*}{0.47\textwidth}
			{@{}l@{\extracolsep{\fill}}l@{\extracolsep{6pt}}l@{\extracolsep{\fill}}c@{}}
				#1 & $\maximize{#2}$ & $#3$ & $ $ \\[5pt]
					 & $\subject\ $    & $#4$ & $ $
			\end{tabular*}}
			}



\newcommand{\boxedeqnnew}[2]{\vspace{12pt}
 \noindent\framebox[\textwidth]{\begin{tabular*}{\textwidth}{p{.8\textwidth}p{.15\textwidth}}
                             \hspace{.05 \textwidth}    #1 & \raggedleft #2
                                \end{tabular*}
}\vspace{12pt}
 }

\newcommand{\boxedeqn}[3]{\vspace{-.2in}
\begin{center}
  \fbox{%
      \begin{minipage}{#2 \textwidth}%
      \begin{equation}#1 #3 \end{equation}%
      \end{minipage}%
    }%
    \end{center}
}
\newcommand{\boxedeqnsmall}[6]{\vspace{12pt}
 \noindent\framebox[#3 \textwidth]{\begin{tabular*}{#3 \textwidth}{p{#5 \textwidth }p{#6 \textwidth }}
                               \hspace{#4 \textwidth}  #1 & \raggedleft #2
                                \end{tabular*}
}\vspace{12pt}
 }


\newcommand{\anybox}[2]
{\fbox{
\begin{minipage}
\begin{tabular*}{#2 \texwidth}{l}
#1
\end{tabular*}
\end{minipage}
}}

		\def\problemsmall#1#2#3#4{\fbox
		 {\begin{tabular*}{0.47\textwidth}
			{@{}l@{\extracolsep{\fill}}l@{\extracolsep{6pt}}l@{\extracolsep{\fill}}c@{}}
				#1 & $\minimize{#2}$ & $#3$ & $ $ \\[5pt]
					  $\sub\ $ &    & $#4$ & $ $
			\end{tabular*}}
			}
\def\problemsmallsingcol#1#2#3#4{\fbox
		 {\begin{tabular*}{0.95\columnwidth}
			{@{}l@{\extracolsep{\fill}}l@{\extracolsep{6pt}}l@{\extracolsep{\fill}}c@{}}
				#1 & $\minimize{#2}$ & $#3$ & $ $ \\[5pt]
					  $\sub\ $ &    & $#4$ & $ $
			\end{tabular*}}
			}

                        \def\problemsmallsingcolwb#1#2#3#4{
		 \begin{tabular*}{0.95\columnwidth}
			{@{}l@{\extracolsep{\fill}}l@{\extracolsep{6pt}}l@{\extracolsep{\fill}}c@{}}
				#1 & $\minimize{#2}$ & $#3$ & $ $ \\[5pt]
					  $\sub\ $ &    & $#4$ & $ $
			\end{tabular*}
			}

		\def\problem#1#2#3#4{\fbox
		 {\begin{tabular*}{0.80\textwidth}
			{@{}l@{\extracolsep{\fill}}l@{\extracolsep{6pt}}l@{\extracolsep{\fill}}c@{}}
				#1 & $\minimize{#2}$ & $#3$ & $ $ \\[5pt]
					 & $\subject\ $    & $#4$ & $ $
			\end{tabular*}}
			}

		\def\minmaxproblem#1#2#3#4#5{\fbox
		 {\begin{tabular*}{0.80\textwidth}
			{@{}l@{\extracolsep{\fill}}l@{\extracolsep{6pt}}l@{\extracolsep{\fill}}c@{}}
				#1 & $\minimize{#2}$ $ \maximize{#3}$& $#4$ & $ $ \\[5pt]
					 & $\subject\ $    & $#5$ & $ $
			\end{tabular*}}
			}
		\def\minmaxproblemsmall#1#2#3#4#5{\fbox
		 {\begin{tabular*}{0.47\textwidth}
			{@{}l@{\extracolsep{\fill}}l@{\extracolsep{6pt}}l@{\extracolsep{\fill}}c@{}}
				#1 & $\minimize{#2}$ $ \maximize{#3}$& $#4$ & $ $ \\[5pt]
					 & $\subject\ $    & $#5$ & $ $
			\end{tabular*}}
			}
		\def\minmaxproblemcompactsmall#1#2#3#4#5{\fbox
		 {\begin{tabular*}{0.47\textwidth}
			{@{}l@{\extracolsep{\fill}}l@{\extracolsep{6pt}}l@{\extracolsep{\fill}}c@{}}
				#1 & $\minimize{#2}$ $ \maximize{#3}$& $#4$ & 			 $\ \subject$     $#5$
			\end{tabular*}}
			}


\def\maxproblem#1#2#3#4{\fbox
		 {\begin{tabular*}{0.80\textwidth}
			{@{}l@{\extracolsep{\fill}}l@{\extracolsep{6pt}}l@{\extracolsep{\fill}}c@{}}
				#1 & $\maximize{#2}$ & $#3$ & $ $ \\[5pt]
					 & $\subject\ $    & $#4$ & $ $
			\end{tabular*}}
			}

\def\maxxproblem#1#2#3#4{\fbox
		 {\begin{tabular*}{1.1\textwidth}
			{@{}l@{\extracolsep{\fill}}l@{\extracolsep{6pt}}l@{\extracolsep{\fill}}c@{}}
				#1 & $\maximize{#2}$ & $#3$ & $ $ \\[5pt]
					 & $\subject\ $    & $#4$ & $ $
			\end{tabular*}}
			}

\def\maxproblemsmallwb#1#2#3#4{
		 \begin{tabular*}{0.47\textwidth}
			{@{}l@{\extracolsep{\fill}}l@{\extracolsep{-4pt}}l@{\extracolsep{\fill}}c@{}}
				#1 &  & $\maximize{#2}$  $#3$ & $ $ \\[4pt]
					  $\sub \ $  &   & $#4$ &  $ $
			\end{tabular*}
			}


\def\maximproblemsmalla#1#2#3#4{\fbox
		 {\begin{tabular*}{0.49\textwidth}
			{@{}l@{\extracolsep{\fill}}l@{\extracolsep{1pt}}l@{\extracolsep{\fill}}c@{}}
				#1 &  & $\maximum{#2}$  $#3$ & $ $ \\[4pt]
					  $\sub \ $  &   & $#4$ &  $ $
			\end{tabular*}}
			}

\def\problemsmallawb#1#2#3#4{
		 \begin{tabular*}{0.47\textwidth}
			{@{}l@{\extracolsep{\fill}}l@{\extracolsep{-4pt}}l@{\extracolsep{\fill}}c@{}}
				#1 &  & $\minimize{#2}$  $#3$ & $ $ \\[4pt]
					  $\sub \ $  &   & $#4$ &  $ $
			\end{tabular*}
                      }
\def\problemsmallabb#1#2#3#4{\fbox{
		 \begin{tabular*}{0.56\textwidth}
			{@{}l@{\extracolsep{\fill}}l@{\extracolsep{-4pt}}l@{\extracolsep{\fill}}c@{}}
				#1 &  & \quad $= $ $\minimize{#2}$  $#3$&   \\[4pt]
                        $\sub \ $  &   & $#4$ &  $ $
			\end{tabular*}}
			}
\def\maxproblemsmallbb#1#2#3#4{\fbox{
		 \begin{tabular*}{0.56\textwidth}
			{@{}l@{\extracolsep{\fill}}l@{\extracolsep{-4pt}}l@{\extracolsep{\fill}}c@{}}
				#1 &  & \quad $ = $ $\maximize{#2}$  $#3$ & $ $ \\[4pt]
					  $\sub \ $  &   & $#4$ &  $ $
			\end{tabular*}
			}}

\def\unconmaxproblemsmall#1#2#3#4{\fbox
		 {\begin{tabular*}{0.47\textwidth}
			{@{}l@{\extracolsep{\fill}}l@{\extracolsep{6pt}}l@{\extracolsep{\fill}}c@{}}
				#1 & $\maximize{#2}$ & $#3$ & $ $ \\[5pt]
			\end{tabular*}}
			}

\def\unconmaxproblem#1#2#3#4{\fbox
		 {\begin{tabular*}{0.95\textwidth}
			{@{}l@{\extracolsep{\fill}}l@{\extracolsep{6pt}}l@{\extracolsep{\fill}}c@{}}
				#1 & $\maximize{#2}$ & $#3$ & $ $ \\[5pt]
			\end{tabular*}}
			}
	\def\problemcompactsmall#1#2#3#4{\fbox
		 {\begin{tabular*}{0.47\textwidth}
			{@{}l@{\extracolsep{\fill}}l@{\extracolsep{6pt}}l@{\extracolsep{\fill}}c@{}}
				#1 &$\minimize{#2}$\ $#3$ &{\rm s.t.} $#4 $
			\end{tabular*}}}
	\def\maxproblemcompactsmall#1#2#3#4{\fbox
		 {\begin{tabular*}{0.47\textwidth}
			{@{}l@{\extracolsep{\fill}}l@{\extracolsep{6pt}}l@{\extracolsep{\fill}}c@{}}
				#1 &$\maximize{#2}$\ $#3$ &{\rm s.t.}\ $#4 $
			\end{tabular*}}}


\def\unconproblem#1#2#3#4{\fbox
		 {\begin{tabular*}{0.95\textwidth}
			{@{}l@{\extracolsep{\fill}}l@{\extracolsep{6pt}}l@{\extracolsep{\fill}}c@{}}
				#1 & $\minimize{#2}$ & $#3$ & $ $ \\[5pt]
			\end{tabular*}}
			}
	\def\cpproblem#1#2#3#4{\fbox
		 {\begin{tabular*}{0.95\textwidth}
			{@{}l@{\extracolsep{\fill}}l@{\extracolsep{6pt}}l@{\extracolsep{\fill}}c@{}}
				#1 & & $#4 $
			\end{tabular*}}}
				\def\viproblem#1#2#3{\small\fbox
					 {\begin{tabular*}{0.47\textwidth}
						{@{}l@{\extracolsep{\fill}}l@{\extracolsep{2pt}}l@{\extracolsep{\fill}}c@{}}
							#1 &#2 & $#3 $
						\end{tabular*}}}

	\def\cp2problem#1#2#3#4{\fbox
		 {\begin{tabular*}{0.9\textwidth}
			{@{}l@{\extracolsep{\fill}}l@{\extracolsep{6pt}}l@{\extracolsep{\fill}}c@{}}
				#1 & & $#4 $
			\end{tabular*}}}
\def\z{\phantom 0}
\newcommand{\pmat}[1]{\begin{pmatrix} #1 \end{pmatrix}}
		\newcommand{\botline}{\vspace*{-1ex}\hrulefill}
		\newcommand{\clin}[2]{c\L(#1,#2)}
		\newcommand{\dlin}[2]{d\L(#1,#2)}
		\newcommand{\Fbox}[1]{\fbox{\quad#1\quad}}
		\newcommand{\mcol}[3]{\multicolumn{#1}{#2}{#3}}
		\newcommand{\assign}{\ensuremath{\leftarrow}}
		\newcommand{\ssub}[1]{\mbox{\scriptsize #1}}
		\def\bkE{{\rm I\kern-.17em E}}
		\def\bk1{{\rm 1\kern-.17em l}}
		\def\bkD{{\rm I\kern-.17em D}}
		\def\bkR{{\rm I\kern-.17em R}}
		\def\bkP{{\rm I\kern-.17em P}}

		\def\bkZ{{\bf{Z}}}
		\def\bkI{{\bf{I}}}

\def\boldp{\boldsymbol{p}}
\def\boldy{\boldsymbol{y}}


\newcommand {\beeq}[1]{\begin{equation}\label{#1}}
\newcommand {\eeeq}{\end{equation}}
\newcommand {\bea}{\begin{eqnarray}}
\newcommand {\eea}{\end{eqnarray}}
\newcommand{\mb}[1]{\mbox{\boldmath $#1$}}
\newcommand{\mbt}[1]{\mbox{\boldmath $\tilde{#1}$}}
\newcommand{\mbs}[1]{{\mbox{\boldmath \scriptsize{$#1$}}}}
\newcommand{\mbst}[1]{{\mbox{\boldmath \scriptsize{$\tilde{#1}$}}}}
\newcommand{\mbss}[1]{{\mbox{\boldmath \tiny{$#1$}}}}
\newcommand{\dpv}{\displaystyle \vspace{3pt}}
\def\texitem#1{\par\smallskip\noindent\hangindent 25pt
               \hbox to 25pt {\hss #1 ~}\ignorespaces}
\def\smskip{\par\vskip 7pt}
%\def\example{{\bf Example :\ }}
\def\st{\mbox{subject to}}

%%% Local Variables:
%%% mode: latex
%%% TeX-master: "report"
%%% End:

\def\bsp{\begin{split}}
		\def\beq{\begin{eqnarray}}
		\def\bal{\begin{align*}}
		\def\bc{\begin{center}}
		\def\be{\begin{enumerate}}
		\def\bi{\begin{itemize}}
		\def\bs{\begin{small}}
		\def\bS{\begin{slide}}
		\def\ec{\end{center}}
		\def\ee{\end{enumerate}}
		\def\ei{\end{itemize}}
		\def\es{\end{small}}
		\def\eS{\end{slide}}
		\def\eeq{\end{eqnarray}}
		\def\eal{\end{align*}}
		\def\esp{\end{split}}
		\def\qed{ \vrule height7.5pt width7.5pt depth0pt}  %width4.17pt depth0pt}
		\def\problemlarge#1#2#3#4{\fbox
		 {\begin{tabular*}{0.47\textwidth}
			{@{}l@{\extracolsep{\fill}}l@{\extracolsep{3pt}}l@{\extracolsep{\fill}}c@{}}
				#1 & $\minimize{#2}$ & $#3$ & $ $ \\[2pt]
					$\subject\ $  &    & $#4$ & $ $
			\end{tabular*}}
			}
		\def\problemsuperlarge#1#2#3#4{\fbox
           {\tiny      {\begin{tabular*}{0.485\textwidth}
                        {@{}l@{\extracolsep{\fill}}l@{\extracolsep{8pt}}l@{\extracolsep{\fill}}c@{}}
                              #1 & $\minimize{#2}$ &  $#3$ & $ $ \\[14pt]
                                         & $\subject\ $    & $#4$ & $ $
                        \end{tabular*}} }
                        }
                        \newenvironment{examplee}{{ \emph{Example:} }}{\hfill $\Box$ \vspace{6pt}}

                        \def\unconproblem#1#2#3{\fbox
		 {\begin{tabular*}{0.47\textwidth}
			{@{}l@{\extracolsep{\fill}}l@{\extracolsep{6pt}}l@{\extracolsep{\fill}}c@{}}
				#1 & $\minimize{#2}$ & $#3$ & $ $ \\[5pt]
			\end{tabular*}}
			}
                        %%%%%%%%%%%%%%%%%%%%%%%%%%%%%%%%%%%%%
%        macros.tex                  %
%%%%%%%%%%%%%%%%%%%%%%%%%%%%%%%%%%%%%%

\usepackage{amsmath, amssymb, xspace}
\usepackage{epsfig}
\usepackage{longtable}
\usepackage{color}
\usepackage{mathrsfs}
\usepackage{subfig}
%\newtheorem{theorem}{Theorem}[section]
%\newtheorem{assumption}{Assumption}[section]
%\newtheorem{lemma}{Lemma}[section]
%\newtheorem{conjecture}{Conjecture}[section]
%\newtheorem{claim}[theorem]{Claim}
%\newtheorem{proposition}[theorem]{Proposition}
%\newtheorem{prop}[theorem]{Proposition}
%\newtheorem{corollary}[theorem]{Corollary}
%\newtheorem{remark}[theorem]{Remark}
%\newtheorem{cor}[theorem]{Corollary}
%\newtheorem{result}[theorem]{Result}
%\newtheorem{definition}{Definition}[section]
%\newenvironment{proof}[1][]{{\noindent \emph {Proof} #1: }}{\hfill \qed \vspace{3pt}\\ }
\def\ext{{\rm ext}}
%\newenvironment{examplee}{{\noindent \bf Example}}{\hfill $\square$\hspace{-4.5pt}\vspace{6pt}}
%\newcommand{\ie}{i.e.\@\xspace}
\def\Cscr{{\cal C}}
\def\Nscr{{\cal N}}
\def\subject{\hbox{\rm subject to}}
\def\sub{\hbox{\rm s.t}}
\def\such{\hbox{\rm such that}}
\def\conv{{\rm conv}}
%\newcommand{\minim}{\mathop{\hbox{\rm min}}}
%\newcommand{\maximize}[1]{\displaystyle\maxim_{#1}}
%\newcommand{\maxim}{\mathop{\hbox{\rm maximize}}}
\newcommand{\maximum}{\mathop{\hbox{\rm max}}}
%\newcommand{\minimize}[1]{\displaystyle\minim_{#1}}
\def\unconproblem#1#2#3{\fbox
		 {\begin{tabular*}{0.47\textwidth}
			{@{}l@{\extracolsep{\fill}}l@{\extracolsep{6pt}}l@{\extracolsep{\fill}}c@{}}
				#1 & $\minimize{#2}$ & $#3$ & $ $ \\[5pt]
			\end{tabular*}}
			}
%			\def\problemsmall#1#2#3#4{\fbox
%		 {\begin{tabular*}{0.8 \textwidth}
%			{@{}l@{\extracolsep{\fill}}l@{\extracolsep{-3pt}}l@{\extracolsep{\fill}}c@{}}
%				#1 & & $\minimize{#2}$ $#3$ & $ $ \\[2pt]
%				$\sub\ $	 &$ $     & $#4$ & $ $
%			\end{tabular*}}
%			}
			\def\problemsmalla#1#2#3#4{\fbox
		 {\begin{tabular*}{0.47\textwidth}
			{@{}l@{\extracolsep{\fill}}l@{\extracolsep{-4pt}}l@{\extracolsep{\fill}}c@{}}
				#1 &  & $\minimize{#2}$  $#3$ & $ $ \\[4pt]
					  $\sub \ $  &   & $#4$ &  $ $
			\end{tabular*}}
			}
			\def\singlecolproblemsmalla#1#2#3#4{\fbox
		 {\begin{tabular*}{0.85\textwidth}
			{@{}l@{\extracolsep{\fill}}l@{\extracolsep{-4pt}}l@{\extracolsep{\fill}}c@{}}
				#1 &  & $\minimize{#2}$  $#3$ & $ $ \\[4pt]
					  $\sub \ $  &   & $#4$ &  $ $
			\end{tabular*}}
			}
			\def\maximproblemsmalla#1#2#3#4{\fbox
		 {\begin{tabular*}{0.49\textwidth}
			{@{}l@{\extracolsep{\fill}}l@{\extracolsep{-4pt}}l@{\extracolsep{\fill}}c@{}}
				#1 &  & $\maximum{#2}$  $#3$ & $ $ \\[4pt]
					  $\sub \ $  &   & $#4$ &  $ $
			\end{tabular*}}
			}
%\def\maximproblemsmalla#1#2#3#4{\fbox
%		 {\begin{tabular*}{0.49\textwidth}
%			{@{}l@{\extracolsep{\fill}}l@{\extracolsep{-4pt}}l@{\extracolsep{\fill}}c@{}}
%				#1 &  & $\maximum{#2}$  $#3$ & $ $ \\[4pt]
%					  $\sub \ $  &   & $#4$ &  $ $
%			\end{tabular*}}
%			}
%			\def\maxproblemlarge#1#2#3#4{\fbox
%		 {\begin{tabular*}{1.0\textwidth}
%			{@{}l@{\extracolsep{\fill}}l@{\extracolsep{-2pt}}l@{\extracolsep{\fill}}c@{}}
%				#1 & $\max{#2}$ & $#3$ & $ $ \\[3pt]
%					 & $\sub $    & $#4$ & $ $
%			\end{tabular*}}
%			}
	\def\maxproblemsmall#1#2#3#4{\fbox
		 {\begin{tabular*}{0.47\textwidth}
			{@{}l@{\extracolsep{\fill}}l@{\extracolsep{-4pt}}l@{\extracolsep{\fill}}c@{}}
				#1 &  & $\maximize {#2}$ $#3$ & $ $ \\[4pt]
					 $\sub \ $  &   & $#4$ & $ $
			\end{tabular*}}
                    }
	\def\maxproblemsmallsingcol#1#2#3#4{\fbox
		 {\begin{tabular*}{0.95\columnwidth}
			{@{}l@{\extracolsep{\fill}}l@{\extracolsep{-4pt}}l@{\extracolsep{\fill}}c@{}}
				#1 &  & $\maximize {#2}$ $#3$ & $ $ \\[4pt]
					 $\sub \ $  &   & $#4$ & $ $
			\end{tabular*}}
                    }
                    	\def\maxproblemsmallsingcolwb#1#2#3#4{
		 \begin{tabular*}{0.95\columnwidth}
			{@{}l@{\extracolsep{\fill}}l@{\extracolsep{-4pt}}l@{\extracolsep{\fill}}c@{}}
				#1 &  & $\maximize {#2}$ $#3$ & $ $ \\[4pt]
					 $\sub \ $  &   & $#4$ & $ $
			\end{tabular*}
			}
\def\singlecolmaxproblemsmall#1#2#3#4{\fbox
		 {\begin{tabular*}{0.85\textwidth}
			{@{}l@{\extracolsep{\fill}}l@{\extracolsep{-4pt}}l@{\extracolsep{\fill}}c@{}}
				#1 &  & $\maximize {#2}$ $#3$ & $ $ \\[4pt]
					 $\sub \ $  &   & $#4$ & $ $
			\end{tabular*}}
			}
%%%%%%%%%%%%%%%%%%%%%%%%%%%%%
% 	bb, bf and bk
%%%%%%%%%%%%%%%%%%%%%%%%%%%%%




\begin{document}
	\maketitle

	%\tableofcontents


	\begin{abstract}
		Strategic classification is the problem of classifying feature vectors that can be manipulated at the testing phase of the classifier. The classifier aims for its decision rule to be robust against strategic manipulation and also be efficiently learnable. In this paper, we present two main ideas for the classifier to achieve this goal. The first idea is to enrich the classifier's decision-making capabilities in two ways: ignoring feature vectors and making randomized decisions. This approach helps the classifier classify more feature vectors correctly by limiting the sender’s options for manipulation. The second idea is to use the strategic structure of the problem during learning. This approach, in certain contexts, can simplify the learning process.  
        Combining these ideas, we propose and compare two learning algorithms, Vanilla ERM and Strategy-aware ERM, and provide a heuristic for selecting the one that yields a classifier both robust to strategic manipulation and efficiently learnable.

        
	\end{abstract}

	\begin{IEEEkeywords}
		Strategic Classification, Penalty Labels, Non-uniform preferences,
		Stackelberg equilibrium, Screening game
	\end{IEEEkeywords}

%	\section*{Information Sheet}
%
%	1. What is the main claim of the paper? Why is this an important contribution to the autonomous agents and multi-agent systems literature?
%
%	This paper addresses the problem of \textit{strategic classification}, where a classifier must correctly classify feature vectors that a sender may manipulate during the test phase. The paper makes two main claims. First, introducing a penalty label improves classification performance against strategic manipulation by the sender. Additionally, incorporating randomization in classification decisions further enhances accuracy. Second, leveraging the problem’s underlying structure during learning leads to more sample-efficient learning, particularly when domain knowledge is available but sample sizes are limited. To the best of our knowledge, the role of a penalty label in improving the accuracy of deterministic classifiers has not been previously explored. These insights can be applied to enhance existing classifiers across various domains. An immediate application of our work is in designing questionnaires for classification-based tasks. Since respondents may provide misleading answers, carefully integrating a penalty label can improve the accuracy of classification outcomes.
%
%	2. What is the evidence you provide to support your claim? Be precise.
%
%	We theoretically demonstrate that introducing a penalty label reduces classification errors compared to a naive classifier. Furthermore, we establish a theoretical result showing that a stochastic classifier can achieve lower classification error than any deterministic classifier, including those using a penalty label. Additionally, incorporating the strategic structure of the problem during learning can sometimes lead to a more sample-efficient learning process. We formalize these results mathematically and support them with examples to confirm their validity.
%
%	3. What papers by other authors make the most closely related contributions, and how is your paper related to them? Provide journal references whenever possible. Don’t rely only on a conference proceedings.
%
%	The paper \cite{hardt2016strategic} models strategic classification as a screening game between a classifier and a sender and provides an algorithm to learn the strategic classifier. However, it assumes uniform sender preferences for labels. Our work extends this framework by incorporating non-uniform sender preferences and introduces an algorithm designed to learn the strategic classifier in this more general setting.  Similarly, the paper \cite{sundaram2021pac} models the problem as a screening game and proposes an algorithm to learn the strategic classifier. While \cite{sundaram2021pac} accounts for non-uniform label preferences, it considers a different hypothesis space for learning the strategic classifier than ours.
%
%	4. Have you published parts of your paper before, for instance in a conference? If so, give details of your previous paper(s) and a precise statement detailing how your paper provides a significant contribution beyond the previous paper(s).
%
%	Some parts  of this work were previously published in two conference papers. A preliminary version of Section \ref{subsec:Discrete feature space}, which considers the discrete feature  space was presented in the  conference paper \cite{singh22strategic} and a preliminary version of Section \ref{subsec: Stochastic strategy}, which considers a stochastic decision rule, was presented in the conference paper \cite{singh2024optimal}.
%	The main contribution of this paper is the analysis of the decision problem for a general feature vector space, which is given in Section \ref{sec:Decision problem}, and the learning problem in this general setting, which is given in Section \ref{sec:Learning problem}. It also refines  assumptions and proofs the above conference papers, to align it with the results of the rest of the paper.
	
	
	
%	 The first conference paper, \cite{singh22strategic}, studied the decision problem in a finite feature space and is covered in Section \ref{subsec:Discrete feature space}. This paper extends the analysis to a general feature vector space and further explores the learning problem in this broader setting.  The second conference paper, \cite{singh2024optimal}, examined the decision problem when the classifier was allowed to make stochastic predictions and is covered in Section \ref{subsec: Stochastic strategy}. The current paper refines \cite{singh2024optimal}, both in its assumptions and proofs, and aligns it with the rest of the paper’s results.



%	\section*{Statements and Declarations}
%
%	The following statements should be included under the heading "Statements and Declarations" for inclusion in the published paper. Please note that submissions that do not include relevant declarations will be returned as incomplete.
%
%	Competing Interests: Authors are required to disclose financial or non-financial interests that are directly or indirectly related to the work submitted for publication. Please refer to “Competing Interests and Funding” below for more information on how to complete this section.
	\section{Introduction}\label{sec:Introduction}

	Consider the following classification situation. A university assigns students to various programs based on their scores in an aptitude test using a decision rule that maps scores to programs. For example, one such rule could be to assign the Physics program to students who have scored well in the physics section of the exam, regardless of their total score. For transparency, the university must declare its decision rule before the students write the exam. Every student taking the exam has a \textit{true aptitude} corresponding it with a \textit{true program}.
However, the students may have preferences for programs that may differ from their true programs. Knowing the decision rule, they can seek external help, such as coaching, to improve their scores, incurring some costs for such assistance. 
Now suppose the university offers two related programs (say theoretical and experimental physics programs), where the students with an aptitude in one program may prefer the other program and, therefore, prepare for the entrance exam accordingly. The university aims to create a question paper for the entrance exam  in such a manner as to assign the maximum number of students to their true programs. The challenge is that the university cannot measure the true aptitude of a student and must rely on the exam scores, i.e., its \textit{demonstrated aptitude}, to assign the program. Moreover, the university does not know the \textit{true function} (which maps true aptitude to a true program) and must learn it from past data.

The university admission problem is an example of a \textit{strategic} classification problem. The strategic classification problem extends the classification problem to the regime where, during the classifier's testing phase, a strategic data sender can manipulate feature vectors to obtain a preferred label while incurring some cost. A non-strategic or \textit{naive classifier} places the decision boundary precisely at the class boundary. Since this classifier's decision rule is known to the strategic sender, it is prone to manipulation. A \textit{strategic classifier} is robust to such manipulation and, in general, differs non-trivially from the naive classifier.

To the best of our knowledge, all previous works on the strategic classification problem assume that the sender has a uniform preference for the labels (e.g. \cite{hardt2016strategic}). The standard example of this setting is the email classification problem, where the strategic sender prefers the non-spam label over the spam label, irrespective of the underlying email content. The strategic classifier for this case is the one that shifts the decision boundary into the \textit{non-spam} region of the feature space, thereby making the criterion for classification as the non-spam label more stringent.

This paper considers a binary classification problem where the strategic sender has a possibly non-uniform preference for labels, as illustrated in the above example for admission to theoretical and experimental physics programs. This assumption poses a new challenge as preferred feature space for one class overlaps with the space not preferred by the other. In particular, the structure of the strategic classifier obtained by \cite{hardt2016strategic} no longer holds without a universally preferred label.
The problem of constructing a strategic classifier with non-uniform preferences results in two subproblems: a \textit{decision problem} and a \textit{learning problem}. In the decision problem, the university knows the \textit{true function} (map between the true aptitude and the true program) and has to calculate the strategic classifier from it. In the learning problem, the university does not know the true function and has to learn the strategic classifier from a finite sample of the true function.

Our main findings are that classification performance improves significantly when the classifier is either allowed to make randomized decisions or allowed to \textit{not} classify certain feature vectors. Moreover, in certain problem instances, particularly those involving complex hypothesis classes, learning can be made easier by incorporating the sender's strategic behaviour in the classifier's learning phase.

\subsection{Main findings}

We adopt a game theoretic viewpoint for our analysis. We model the binary classification problem with non-uniform preferences as a screening game\footnote{In a screening game, the receiver commits to a strategy first; the sender then observes this commitment and chooses a strategy in response, after which the receiver plays the committed strategy.} between a sender and a receiver (classifier), with the receiver as the leader. The sender wants to maximize the difference between its utility and its cost and the receiver wants to maximize the probability of correctly classifying the feature vectors. The strategic classifier corresponds to the receiver's Stackelberg equilibrium strategy. \textcolor{red}{For most of our results, we work with a scalar feature space. We show that under certain assumptions, the problem with the vector feature space with a linear class boundary can be reduced to the problem we consider.} 

We introduce two main ideas. The first is to expand the classifier's capabilities in two ways. One way is by introducing a new penalty label $\Delta$, by which means the receiver ignores a feature vector assigned to this label. From the sender's perspective, this situation is equivalent to receiving a $-\infty$ utility for the label $\Delta$. In the context of university admission problems, the student is the sender, and the university is the receiver. The exam paper comprises of a subset of questions from a question bank, with the requirement that the student has to attempt exactly one question. The university assigns a label to each question in the exam paper with the implication that a student who correctly answers that question gets assigned that label or program. Assignment of label $\Delta$ to a question amounts to \textit{dropping} that question from the exam paper. This approach is in stark contrast to the uniform preference case, where \textit{all} the questions in the question bank are included in the question paper, and only \textit{cut-off marks} for the preferred program are increased.
Ignoring feature vectors can also be considered a generalization of feature selection.

Another way to expand the classifier's action set is by allowing randomized classification. In the university admission problem, this is analogous to saying that solving a particular question is correlated (but not with certainty) with admission to a specific program. These ideas for expanding the classifier's action set are motivated from mechanism design, where penalties influence a strategic sender's behavior, and game theory, where randomization allows for richer outcomes.

The second main idea involves incorporating the underlying game-theoretic structure of the problem by modifying the hypothesis class used for learning. Learning from the original hypothesis class will be inefficient if the university's inductive bias does not consider students' strategic nature. We show that incorporating the game's strategic structure through hypothesis class modification can make an unlearnable hypothesis class learnable.

Using these two ideas, we address the decision and learning problems of strategic classification in the scalar feature space setting. For the decision problem, our first finding is that a deterministic strategic classifier using a penalty label makes fewer classification errors than any classifier that does not use a penalty label, at least strictly fewer errors in some instances of the problem. 
Furthermore, under the adversarial assumption, i.e., when the feature vector always prefers the label opposite to its true label, and assuming a uniform distribution over the feature space with the Euclidean cost function, a stochastic strategic classifier makes fewer errors than any deterministic classifier, including those that use a penalty label, and at least strictly fewer errors than a classifier with a penalty label in some instances of the problem.

However, the stochastic classifier uses only randomization on the original label set; using the penalty label does not increase its accuracy. Nonetheless, the penalty label becomes relevant when it is constrained to use a deterministic strategic classifier, a common situation encountered in real-life applications. Hence, the penalty label can be thought of as partially compensating the classifier for the absence of randomization.



The classifier described above assigns the label $\Delta$ to feature vectors near the class boundary. Indeed, multiple such classifiers exist, all with the same classification error, but they differ in the following three aspects.
The first aspect is classifier complexity, measured by the number of decision boundaries. A decision boundary is a feature vector at which the classifier changes its assigned label. Thus, the number of decision boundaries corresponds to the number of times the classifier switches labels as the feature vector varies. We show that classifiers with fewer decision boundaries are easier to learn due to their lower VC dimension.
The second aspect is classifier fairness, which refers to fair distribution of errors across classes. If in addition to classification performance, fairness is also in consideration, then the receiver can pick among the many classification rules, the one which proportionately distributes the error across classes. The third aspect is the sender's payoff under a strategic classifier. This aspect becomes relevant if a sender can opt-out and obtain a reservation payoff. Thus, a receiver may select a classifier that gives the sender a payoff higher than its reservation payoff. \textcolor{red}{ Although we work with scalar features, our problem suffices to demonstrate the structural
aspects of strategic classification and the importance of a penalty label. Moreover, in
 many university admission problems like those we consider, the final admission decision
 has to be done based on a single scalar score, whereby the scalar problem is practically relevant.}


We introduce two new algorithms called vanilla ERM and strategy-aware ERM to address the learning problem. In vanilla ERM, the classifier first uses ERM to learn the true function from the original hypothesis class. The decision problem is then applied to the learned hypothesis. In the strategy-aware ERM approach, the hypothesis class is modified to incorporate the underlying game-theoretic structure. This modification involves a two-step process. First, each hypothesis in the original class is replaced with its corresponding strategic classifier, which is obtained by applying the decision problem to each hypothesis. Second, the strategic classifier obtained in the previous step is further replaced with a new hypothesis, formed by composing the strategic classifier with the sender’s worst-case best response. The classifier then uses the Empirical Risk Minimization (ERM) algorithm to learn from this modified hypothesis class. What makes strategy-aware ERM novel is its ability to transform an otherwise unlearnable problem into a learnable one. We show that there exist hypothesis classes that are unlearnable under vanilla ERM but become learnable with strategy-aware ERM, illustrating how incorporating game-theoretic structure can fundamentally improve learnability.

While strategy-aware ERM can learn from a broader range of hypothesis classes, its advantage in the ease of learning over vanilla ERM only sometimes holds. This ease of learning is controlled by the VC dimension of the hypothesis class, which in turn determines the sample complexity in PAC learning. The comparison, therefore, becomes context-dependent if a hypothesis class is learnable via vanilla ERM. \textcolor{red}{We provide a heuristic argument that strategy-aware ERM generally makes learning easier if the original hypothesis class is too complex, as it can reduce the effective VC dimension. On the other hand, vanilla ERM proves more efficient if the original hypothesis class is simple. We show this with the help of an example. A formal result, even for scalar features, remains elusive.}


\subsection{Related Works}
The strategic classification problem has been a topic of interest in recent years. Works such as \cite{bruckner2011stackelberg,ghalme2021strategic,bechavod2022information,cohen2024bayesian,jagadeesan2021alternative,dalvi2004adversarial} explicitly formulate the strategic classification problem as a game between a strategic sender and a classifier. In our work, we model the interaction between the strategic sender and the classifier as a sender-receiver screening game. We introduce a penalty label, which expands the receiver's strategy space and gives rise to equilibrium strategies that are robust against the sender's manipulation. To the best of our knowledge, no such penalty label has been considered in the literature so far.

Works such as \cite{hardt2016strategic,cullina2018pac,zhang2021incentive,sundaram2021pac} consider the impact of a strategic sender on the PAC learnability of the hypothesis class. In \cite{hardt2016strategic,cullina2018pac,zhang2021incentive}, the sender is assumed to have a uniform preference for positive labels, whereas in \cite{sundaram2021pac}, the sender is assumed to have a non-uniform preference for labels for each feature vector. In \cite{hardt2016strategic,cullina2018pac,sundaram2021pac}, strategic manipulation incurs a cost, whereas in \cite{zhang2021incentive}, strategic manipulation does not incur a cost. The works \cite{hardt2016strategic,cullina2018pac,zhang2021incentive,sundaram2021pac} do not explicitly utilize the underlying sender-receiver game in the learning process. The learning process in these works can be equivalently described as using the ERM algorithm to optimize over a modified hypothesis class, obtained by replacing each hypothesis with the composition of the hypothesis and the sender's best response to that hypothesis. In contrast, in our work, the sender has a non-uniform preference over labels for each feature vector, and the structure of this preference differs from that in \cite{sundaram2021pac}. Furthermore, in our work, there is a cost incurred for modifying a feature vector. We explicitly make use of the underlying sender-receiver game in the learning process by allowing the ERM algorithm to optimize over a modified hypothesis class, which is obtained by replacing each hypothesis with its corresponding Stackelberg equilibrium composed with the sender's best response.

There are other interesting formulations of the problem that we do not consider in our work. For instance, some research explores scenarios where the sender lacks knowledge of the classifier. Notable works in this area include \cite{ghalme2021strategic}, \cite{bechavod2022information}, and \cite{cohen2024bayesian}. Additionally, other studies introduce more structure to the sender by incorporating user types, as seen in \cite{jagadeesan2021alternative} and \cite{levanon2022generalized}.

Finally, the social welfare implications of strategic classification have been explored in works such as \cite{milli2019social}, which analyze the social costs imposed on the sender. In our work, we briefly address related issues such as fairness and participation constraints in Section~\ref{subsec:Other equilibria}.





% This line of analysis allows us to discuss (Section \ref{subsec:Other equilibria}) the complexity and fairness of the hypothesis class.

% The strategic classification problem, modeled as a game between a strategic sender and a classifier, has been extensively studied in works such as \cite{dalvi2004adversarial}, \cite{bruckner2012static}, \cite{hardt2016strategic}, and \cite{bruckner2011stackelberg}. Among these, \cite{bruckner2012static} treats the interaction as a simultaneous game, while \cite{bruckner2011stackelberg} and \cite{hardt2016strategic} adopt a screening game framework. Our work aligns with the screening game approach.


% In \cite{hardt2016strategic}, the sender is assumed to have uniform preferences over labels. In contrast, we consider a non-uniform preference structure, introducing additional complexity to the classification problem. Key differences include the use of a different cost function, the inclusion of a penalty label, the nature of the true function, the sender's behavior, and other structural aspects.

% Several theoretical studies have extended PAC learning to strategic settings, including \cite{zhang2021incentive}, \cite{cullina2018pac}, and \cite{sundaram2021pac}. Specifically, an independent study, \cite{sundaram2021pac}, also explores non-uniform sender preferences and the strategic learnability of hypothesis classes. Our contribution lies in proposing a novel definition of strategic learnability (see Definition \ref{def:Strategically learnable class}). Unlike prior definitions that focus solely on sender behavior, our approach accounts for the strategic behavior of both the sender and the receiver. As a result, our proposed Algorithm (\ref{algo:Startegically aware ERM}) for learning a strategic classifier differs from the Incentive-Aware algorithm \cite{zhang2021incentive} and the SERM algorithm \cite{sundaram2021pac}.

% There are other interesting formulations of the problem that differ from ours. For instance, some research explores scenarios where the sender lacks knowledge of the classifier. Notable works in this area include \cite{ghalme2021strategic}, \cite{bechavod2022information}, and \cite{cohen2024bayesian}. Additionally, other studies introduce more structure to the sender by incorporating user types, as seen in \cite{jagadeesan2021alternative} and \cite{levanon2022generalized}.

% Finally, the social welfare implications of strategic classification have been explored in works like \cite{milli2019social}, which analyze social costs imposed on the sender. Within our framework, we briefly address related issues such as fairness and participation constraints in Section \ref{subsec:Other equilibria}.












	\section{Problem Formulation}\label{sec:Problem formulation}

	\input{Problem Formulation}

	\section{Decision Problem}\label{sec:Decision problem}

	\input{Decision Problem}

	\subsection{Motivating example: discrete feature space}\label{subsec:Discrete feature space}

	To motivate the solution of the decision problem, we digress and recall a related problem from \cite{singh22strategic}. The problem formulation is almost the same as described in Section \ref{sec:Problem formulation}, and we only highlight the differences.
Let $\cg{X} = \{1,2,\cdots,n\}$ be a finite set, and $h:\cg{X} \rightarrow \cg{Y}$ be the true function.
Let $U: \cg{Y} \times \cg{X} \rightarrow \mathbb{R}$ be the utility function such that $U(h(x),x) = 0$ and $C: \cg{X} \times \cg{X} \rightarrow \mathbb{R}$ be the cost function such that $C(x,x) = 0$. We now characterise the Stackelberg equilibrium strategy in terms of the independence number of a sender graph.  The \textit{sender graph} is defined as follows.
\begin{defn}\label{def:Sender's graph}
	Let $G_h = (V,E)$ be a graph whose node set $V$ is given by $V = \cg{X}^2$. The edge set $E \subset \cg{X}^2 \times \cg{X}^2$ is described as follows: There is an edge between $(\ibb,\kbb)$ and $(\jbb,\lbb)$ if $h(\ibb) \neq h(\jbb)$ and one of the following conditions holds:
	\begin{enumerate}[label=(\alph*)]
		\item $\kbb = \lbb$.
		\item \begin{equation*}
			\begin{aligned}
				&U(h(\jbb),\ibb) - C(\lbb,\ibb) + C(\kbb,\ibb) \geq 0 \text{ or}\\
				&U(h(\ibb),\jbb) - C(\kbb,\jbb) + C(\lbb,\jbb) \geq 0.
			\end{aligned}
		\end{equation*}
	\end{enumerate}
\end{defn}
Recall that a subset of nodes in a graph is called an independent set if no edge connects any two nodes in that subset. The next theorem characterises the Stackelberg equilibrium strategy in terms of the maximum independent set of the sender graph.
\begin{thm}\label{thm:Stackelberg equilibrium for discrete feature space}
	Let $I$ be an independent set of the graph $G_h$ and $r(\cdot;I)$ be a receiver strategy defined as
	\begin{equation*}
		r(\kbb;I) := \begin{cases}
			h(\ibb) & (\ibb,\kbb) \in I\\
			\Delta & (\ibb,\kbb) \not\in I.
		\end{cases}
	\end{equation*}
	The following statements hold.
	\begin{enumerate}[label=(\alph*)]
		\item If $I^{\star}$ is a maximum independent set of $G_h$, then $\rs(\cdot) := r(\cdot,I^{\star})$ is a Stackelberg equilibrium.	
		\item Conversely, if $\rs(\cdot)$ is a Stackelberg equilibrium, then $I^{\star} := \{(x,\srs(x)): \rs(\srs(x)) = h(x)\}$ is a maximum independent set in $G_h$, and furthermore, $\rs(\cdot) = r(\cdot,I^{\star})$. 
	\end{enumerate}
\end{thm}

We require the following lemma to establish the relationship between an independent set of graph $G_h$ and the correct classification of the feature vector.

\begin{lem}\label{lem:1-1 correspondance between independence set and recoverable feature vector}
	Let $I$ be an independent set of the graph $G_h$ and $r(\cdot;I)$ be a receiver strategy defined as
	\begin{equation*}
		r(\kbb;I) := \begin{cases}
			h(\ibb) & (\ibb,\kbb) \in I\\
			\Delta & (\ibb,\kbb) \not\in I.
		\end{cases}
	\end{equation*}
    By the definition of $I$, $r(\cdot,I)$ is well defined. The following statements hold:
	\begin{enumerate}[label=(\alph*)]
		\item If $I$ is an independent set of $G_h$, then for a receiver strategy $r(x) := r(x,I)$, the set $\{(x,s_r(x)): r(s_r(x)) = h(x)\}$ contains $I$.
		\item If $r(x)$ is a receiver strategy, then $I := \{(x,s_r(x)): r(s_r(x)) = h(x)\}$ is an  independent set in $G_h$. 
	\end{enumerate}
\end{lem}

Now, we proceed with our proof.
\begin{proof}
	\begin{enumerate}[label=(\alph*)]
		\item Let $(\ibb,\kbb) \in I$. If $s_r(\ibb) = \lbb \neq \kbb$, then it follows that $r(\lbb) \neq \Delta$. Hence their exists a $\jbb$ such that $(\jbb,\lbb) \in I$. Since $s_r(\ibb) = \lbb$, it implies $U(h(\ibb),\ibb) - C(\kbb,\ibb) \leq U(h(\jbb),\ibb) - C(\lbb,\ibb).$ But from Definition \ref{def:Sender's graph}, this implies that there is an edge between $(\ibb,\kbb)$ and $(\jbb,\lbb)$ which contradicts the independence of the set $I$. Hence $s_r(\ibb) = \kbb$ and therefore $r(s_r(\ibb)) = h(\ibb)$. Since $\ibb$ was an arbitrary element in $I$, therefore $\{(x,s_r(x)): r(s_r(x)) = h(x)\} \supseteq I$.
		
		\item Let $\ibb, \jbb$ be such that $r(s_r(\ibb)) = h(\ibb)$ and $r(s_r(\jbb)) = h(\jbb)$. Let $\kbb := s_r(\ibb)$ and $\lbb := s_r(\jbb)$. We now show there is no edge between $(\ibb,\kbb)$ and $(\jbb,\lbb)$. If $h(\ibb) = h(\jbb)$, then we are trivially done with the proof. Therefore, assume $h(\ibb) \neq h(\jbb)$. Now, we check for the first edge conditions. If $\kbb = \lbb$, then $r(\kbb) = r(\lbb)$ which implies $h(\ibb) = h(\jbb)$ which is a contradiction. Therefore, the first edge condition does not hold. Next, we check for the second edge condition. If $U(h(\jbb),\ibb) - C(\lbb,\ibb) + C(\kbb,\ibb) \geq 0$ (same reasoning holds for the other equations in the second edge condition), then it implies that the sender receives greater payoff when it modifies $s_r(\ibb)$ to $\lbb$ instead of $\kbb$ which is a contradiction since we assumed $s_r(\ibb) = \kbb$. Therefore, the second edge condition does not hold.
		Hence there is no edge between $(\ibb,\kbb) =(\jbb,\lbb)$. Since $(\ibb,\kbb)$ and $(\jbb,\lbb)$ were arbitrary elements of $G_h$, therefore the set $\{(x,s_r(x):r(s_r(x)) = h(x)\}$ is an independent set in $G_h$.
		
	\end{enumerate}
\end{proof}
We now conclude the proof of Theorem \ref{thm:Stackelberg equilibrium for discrete feature space}.
\begin{proof}(Theorem \ref{thm:Stackelberg equilibrium for discrete feature space}).
	\begin{enumerate}[label=(\alph*)]
		\item If $I^{\star}$ is a maximum  independent set, then $r(x; I^{\star})$ is a Stackelberg equilibrium because if not, then there exists another strategy $r^{\star}$ such that $\{(x,\srs(x)):\rs(\srs(x)) = h(x)\}$ contains $\{(x,s_r(x)): r(s_r(x)) = h(x)\}$ which in turn contains $I^\star$. Since $\{(x,\srs(x)):\rs(\srs(x)) = h(x)\}$ and $\{(x,s_r(x)): r(s_r(x)) = h(x)\}$ are both independent set in $G_h$ from Lemma \ref{lem:1-1 correspondance between independence set and recoverable feature vector} b, therefore this contradict the maximality of $I^\star$.
		\item If $r(x)$ is a Stackelberg equilibrium, then $I := \{(x,s_r(x))\}$ is a maximum independent set because if not, then there exists another independent set $J$ such that $|J| < |I|$. This implies the for the receiver strategy $\rt(x;J)$, $\{(x,s_{\rt}(x))\} \supseteq J \supset I$. This equation contradicts the assumption that $r(x;I)$ is a Stackelberg equilibrium. 
	\end{enumerate}
	
\end{proof}

We have shown that for the decision problem for a discrete feature vector set, the Stackelberg equilibrium strategy $r$ is better than the naive classifier $h$. One of the reasons for this is that $r$ can strategically choose to ignore feature vectors that are manipulation-prone. 

	\subsection{Deterministic strategies for the receiver}\label{subsec:Deterministic strategy}

	In this subsection, we address the decision problem formulated in Section \ref{sec:Problem formulation}.
%The best response of the sender, as given in Definition \ref{defn:Best response}, is defined in terms of the expected payoff of the sender. However, in the gameplay, the sender commits to a strategy after observing the realized feature vector $x$, which is sampled according to a probability law $D \in \fk{D}$. Consequently, instead of maximizing the expectation of his payoff, the sender could maximize his payoff function pointwise for each $x \in \cg{X}$. The next lemma states this observation.
%
%\begin{lem} Let $D \in \fk{D}$ be probability law on $\cg{X}$ and let $r$ be a receiver strategy. Then $s \in B(r)$ (best response set) if and only if
%\[s(x) := \arg\max_{s' \in \msf{S}} U(r(s'(x))) - C(s'(x),x)\; \text{almost surely}.\]
%\end{lem}
%
%\begin{proof}
%Let $\dot{\psi}$ be the probability density function of the probability law $D$. Let $s \in B(r)$. Let $Z \subset \cg{X}$ be such that for $x \in Z$, we have $s(x) \neq \arg\max_{s' \in \msf{S}} \left( U(r(s'(x))) - C(s'(x),x) \right)$. Suppose $D(Z) \neq 0$. Now, consider an alternative sender strategy $\bar{s}$ defined as
%\[\bar{s}(x) = \begin{cases}
%	s(x) & \text{for} x \in \cg{X}-Z \\
%	\arg\max_{s' \in \msf{S}} U(r(s'(x))) - C(s'(x),x) & \text{for} x \in Z.\\
%	\end{cases}
%	\]
%It then follows that
%\[
%\int_Z \mathbb{U}(x;r,s) \dot{\psi}(x) \, dx < \int_Z \mathbb{U}(x;r,\bar{s}) \dot{\psi}(x) \, dx
%\]
%and
%\[\int_{\cg{X} \setminus Z} \mathbb{U}(x;r,s) \dot{\psi}(x) \, dx = \int_{\cg{X} \setminus Z} \mathbb{U}(x;r,\bar{s}) \dot{\psi}(x) \, dx.\]
%Hence, we conclude that $\cg{U}(r,s) < \cg{U}(r,\bar{s})$. However, this contradicts the assumption that $s \in B(r)$. Therefore, we must have $D(Z) = 0$.
%Conversely, let $Z \subset \cg{X}$ such that $D(Z) = 0$, and let $\bar{s}$ be such that $\bar{s}(x) = \arg\max_{s' \in \msf{S}} \left( U(r(s'(x))) - C(s'(x),x) \right)$ for $x \in Z$. Let $s \in B(r)$. By the definition of $\bar{s}$, we have:
%\[
%\int_Z \mathbb{U}(x;r,s) \dot{\psi}(x) \, dx \leq \int_Z \mathbb{U}(x;r,\bar{s}) \dot{\psi}(x) \, dx
%\]
%and
%\[\int_{\cg{X} \setminus Z} \mathbb{U}(x;r,s) \dot{\psi}(x) \, dx = \int_{\cg{X} \setminus Z} \mathbb{U}(x;r,\bar{s}) \dot{\psi}(x) \, dx = 0.\]
%Thus, $\cg{U}(r,s) \leq \cg{U}(r,\bar{s})$. Since $s \in B(r)$, the only way to avoid a contradiction is if
%\[
%\int_Z \mathbb{U}(x;r,s) \dot{\psi}(x) \, dx = \int_Z \mathbb{U}(x;r,\bar{s}) \dot{\psi}(x) \, dx.
%\]
%This implies that the sender obtains the same payoff with $\bar{s}$ as with $s$. Therefore, $\bar{s} \in B(r)$.
%\end{proof}
We analyze the problem for various cases of utility functions.\\
%I. \textbf{$\um, \up \leq 0$:} \todo{make subsubsection instead of manual labels}\\
\subsubsection{$\um, \up \leq 0$}
We show that this case is uninteresting as the objectives of the sender and the receiver are aligned as seen in the next proposition.


\begin{thm} Let $\up,\um \leq 0$. Then,
$r(x) \equiv h(x)$ is a Stackelberg equilibrium.
	
%	 \hltodo{$h$ is a Stackelberg equilibrium.}{strategy for receiver? write $r=h$}
\end{thm}
\begin{proof}
	First, we will show that $s_h(x) = x$. Let $x \leq \mathbf{h}$, then $h(x) = \minus$.
    Since $\up \leq 0$, it follows that
\[U(h(x),x) = 0 \geq \max_{x' \in \cg{X}}U(h(x'),x) - |\phi(x')-\phi(x)|.\]
Therefore $s_h(x) = x$. Similar reasoning holds for $x > \mathbf{h}$.
	Now, the set of feature vectors whose label is correctly classified by $h$ is given by the set
	$$E(h,s_h;h) = \{x \in \cg{X}| h(s_h(x)) = h(x)\} = \cg{X}.$$
	Since $h$ correctly classifies all the feature vectors, it is a Stackelberg equilibrium.
\end{proof}


%\noindent II.  \textbf{$\up \geq 1 \geq \um$:}\\
\subsubsection{$\up \geq 1 \geq \um$}

For this case we show that any receiver $r$ whose range contains the $\plus$ label, can at most only correctly classify  feature vectors whose true labels is $\plus$.



\begin{thm}\label{prop:constant is equib}
The receiver strategy $r(c) \equiv c$ for all $x \in \cg{X}$, where $c$ is the majority label is a Stackelberg equilibrium.
	% Let $r$ be the receiver strategy such that their exists $a \in  \cg{X}$ such that
	% $r(a) = \plus$,
	% then $E(r) \supset [0,\textbf{h}]$.
\end{thm}

\begin{proof}
Let $r$ be the receiver Stackelberg equilibrium.
Then, either there does not exist any $x \in \cg{X}$ such that $r(x) = \plus$, or there exists some $t \in \cg{X}$ such that $r(t) = \plus$.
Suppose there does not exist any $x \in \cg{X}$ such that $r(x) = \plus$. Then $r(s_r(x)) \equiv \minus$ for all $x \in \cg{X}$. Thus, $E(r) \subset [\h,1]$. However, the receiver strategy $r'(x) \equiv \minus$ for all $x \in \cg{X}$ yields $E(r') = [\h,1]$. Since $r$ is a Stackelberg equilibrium, it follows that $\E(r) = [\h,1]$.
Now, suppose there exists $t \in \cg{X}$ such that $r(t) = \plus$. Since $\up \geq 1$, it follows that for all $x \in [0,\h)$, $r(s_r(x)) = \plus$. Therefore, $E(r) = [0,\h)$.
Comparing these two cases, it follows that $r(x) \equiv c$ for all $x \in \cg{X}$, where $c$ is the majority label in the Stackelberg equilibrium.
\end{proof}


%\noindent III.  \textbf{$\um,\up \geq 1$:\\}
\subsubsection{$\um,\up \geq 1$}
For this case, we show that for any receiver strategy $r$ whose range contains both $\minus$ and $\plus$ labels does not correctly classify any feature vector.

\begin{thm}
The receiver strategy $r(x) \equiv c$ for all $x \in \cg{X}$, where $c$ is the majority label is a Stackelberg equilibrium.
	
\end{thm}

\begin{proof}
Let $r$ be a receiver Stackelberg equilibrium such that there exist $a, b \in \cg{X}$ with $r(a) = \plus$ and $r(b) = \minus$. Since $\up > 1$, for all $x \in [0, \h)$, the worst-case best response is $s_r(x) = a$, so $r(s_r(x)) = \plus$. Similarly, for $x \in [\mathbf{h}, 1)$, $s_r(x) = b$, so $r(s_r(x)) = \minus$. Therefore, $E(r) = \cg{X}$ and $\cg{E}(r) = 1$.
However, for a receiver strategy $r'(x) \equiv c$ for all $x \in \cg{X}$, where $c$ is the majority label, we have $\cg{E}(r') < 0.5$. This contradicts the optimality of the Stackelberg equilibrium strategy $r$. Therefore, no such $a, b$ exist. This implies that $r$ must be a constant label strategy. Thus, $r(x) \equiv c$ for all $x \in \cg{X}$, where $c$ is the majority label, is a receiver Stackelberg equilibrium strategy.
\end{proof}


\subsubsection{$1 \geq \up \geq 0$ and  \textbf{$|\um| < \up$}}
This case includes scenarios where the sender has a strictly non-uniform preference($\um \geq 0$) as well as scenarios where the sender has a uniform preference ($-\up \leq \um < 0$).
We show that in each of these cases, there exists a receiver Stackelberg equilibrium strategy that belongs to the family of strategies $\sr{R}$ defined as follows.

Let $\ib,\jb$ be such that $0 \leq \ib \leq \h < \jb \leq 1$ and $\phi(\jb) - \phi(\ib) = \up$. Since $\phi$ is a continuously increasing function, such $\ib, \jb$ exist. Let $r(\cdot;\ib,\jb)$ denote a receiver strategy defined as follows
	\begin{equation*}\label{eqn:simple strategies}
		r(x;\ib,\jb) := \begin{cases}
			h(x) & \text{if} x \in [0,\ib] \cup [\jb,1] \\
			\Delta & \text{if} x \in (\ib,\jb)
		\end{cases}.
	\end{equation*}
Let $$\sr{R}_1 := \{r(\cdot;\ib,\jb): 0 \leq \ib \leq \h < \jb \leq 1, \phi(\jb) - \phi(\ib) = \up\}.$$
The strategy $r(\cdot,\ib,\jb) \in \sr{R}_1$ is identical to the true function $h$ except in the interval $(\ib,\jb)$, where it maps the feature vectors to label $\Delta$.
Let$$\sr{R}_2 := \{r(\cdot): r(\cdot) \equiv \minus\} \cup \{r(\cdot): r(\cdot) \equiv \plus\},$$  be the set of constant-label receiver strategies. Let $\sr{R} := \sr{R}_1 \cup \sr{R}_2$.


We define dominance between receiver strategies as follows: for a given probability law $D$ on $\cg{X}$, a strategy $r$ dominates another strategy $\rt$ if $\cg{E}(r) \leq \cg{E}(\rt)$. The next theorem shows that the family of receiver strategies $\sr{R}$ dominates every other receiver strategy. Hence, a receiver can restrict its search for a Stackelberg equilibrium to the set $\sr{R}$ itself.
\begin{thm}\label{thm:Deterministic Stackelberg equilibrium strategy} Let $D \in \fk{D}$ be a law on $\cg{X}$ and let $h(x)$ be a true function. Let $\phi$ be a cost kernel and $\up > 0$. Then their exists a  Stackelberg equilibrium strategy $r(\cdot) \in \sr{R}$.
%	\begin{enumerate}[label=(\alph*)]
%		\item If $\up < \min\{\phi(\h),1-\phi(\h)\}$, then there exist $\ib,\jb \in (0,1)$ such that $r(x;\ib,\jb) \in \sr{R}_1$ is a Stackelberg equilibrium.
%		\item If $\up \geq \min\{\phi(\h),1-\phi(\h)\}$, then $r \in \sr{R}_2$, $r(\cdot) \equiv c$, where $c$ is the majority label, is a Stackelberg equilibrium.
%	\end{enumerate}
\end{thm}



\begin{proof} The proof is given in the Proof \ref{proof:Best response r1} in the Appendix.
\end{proof}



% \begin{proof} The proof is given in the Proof \ref{proof:Deterministic domination of a family of strategies} in the Appendix.
% \end{proof}



%\noindent \textbf{V.} \textbf{$1 \geq \up \geq 0$} and \textbf{$\um \leq -\up$:}\\
\subsubsection{$1 \geq \up \geq 0$} and \textbf{$\um \leq -\up$}

This case corresponds to the sender having an uniform preference. The next theorem specifies a Stackelberg equilibrium.

\begin{thm}\label{thm:u_ < - u_+}
	Let $D \in \fk{D}$ be a probability law on $\cg{X}$ and let $h(x)$ be a true function. Let $\jb \in \cg{X}$ be such that $\phi(\jb) - \phi(\h) = \up$. Then the receiver strategy
	\[ r(x;\h,\jb):=
	\begin{cases}
		\minus & \text{for } x \in [0,\h],\\
		\Delta & \text{for } x \in (\h,\jb),\\
		\plus & \text{for } x \in [\jb,1],
	\end{cases}\]
	is a Stackelberg equilibrium.
\end{thm}

\begin{proof} The proof is given in the Proof \ref{proof:u_ < - u_+} in the Appendix.
	\end{proof}

        \subsection{Multi-dimensional decision problem}

        \input{Multidimension decision problem}
        
	\subsection{Stochastic strategies for the receiver}\label{subsec: Stochastic strategy}

	In this section, we allow the receiver to use a randomized classifier and characterize the Stackelberg equilibrium of the game. This work improves upon our previous study in \cite{singh2024optimal} by refining key assumptions and proofs. Let the cost kernel  $\phi(x)$ be $\phi(x) = x$, and let $D$, the probability law of $X$, be uniform. Let $\cg{P}(\cg{Y} \cup \Delta)$ be the set of probability laws on $\cg{Y} \cup \Delta$. A receiver strategy $r$ is a randomized classifier if it is of the form $r: \cg{X} \to \cg{P}(\cg{Y} \cup \Delta)$. This means that for each $x$, $r(x)$ is a probability law on $\cg{Y} \cup \Delta$. In the strategic classification problem, the receiver observes the values of $\cg{X}$ as communicated by the sender. Let $X'$ denote the random variable representing the receiver's observation of $\cg{X}$ that is communicated by the sender, and let $D'$ denote the corresponding probability law. For a sender strategy $s \in \msf{S}$ , the random variable $X'$ is given by $X' = s(X)$, and the probability law $D'$ is given by
\[
D'_s(X' \in A) = D(X \in \{x : s(x) \in A\}) = \int_{\{x : s(x) \in A\}} dx, \quad \text{for } A \subseteq \cg{X}.
\]
We define a joint random variable $(X',Y)$ over $(\cg{X} \times \cg{Y})$ whose probability law is given by
\[\prob_r(X' \in [x_i,x_j],Y=y) := \int_{x_i}^{x_j}\prob_r(Y=y|X'=x)dD'
\]
where
\begin{equation}\label{eqn:conditional}
  \prob_{r}(Y=y|X'=x) := r(x)(\{y\}), y \in \cg{Y}\cup \Delta.
\end{equation}
% Note that in \eqref{eqn:conditional}, $\prob_r(Y=y|X' = x)$ depends only on $x$, i.e., for two different choices of sender strategies $s_1$ and $s_2$, $\prob_r(Y=y|X'_{s_1} = x) = \prob_r(Y=y|X'_{s_2} = x)$.
We further constrain $r$ such that
\begin{equation}\label{eqn:receiver condition 1}
    \exists x \in \cg{X} \text{ for which } \prob_r(\Delta \mid X' = x) = 0.
\end{equation}
and the sets
\begin{equation}\label{eqn:receiver_condition_2}
\begin{aligned}
    &\{x : \prob_r(- \mid X' = x) + \prob_r(+ \mid X' = x) = 1, x \leq \h\} \text{ and } \\
    &\{x : \prob_r(- \mid X' = x) + \prob_r(+ \mid X' = x) = 1, x > \h\} \text{ are closed.}
\end{aligned}
\end{equation}
The receiver's strategy set is defined as
\begin{equation}\label{defn:stochastic_receiver}
	\msf{R} := \{r \mid r: \cg{X} \to \cg{P}(\cg{Y} \cup \Delta), \; r \text{ satisfies } \eqref{eqn:receiver condition 1} \text{ and } \eqref{eqn:receiver_condition_2}\}.
\end{equation}
For a receiver strategy $r$, let
$$p(x) := \prob_r(Y=\minus|X'=x) \text{ and } q(x) := \prob_r(\{Y=\plus|X'=x\})$$
be the probabilities of assigning labels $-$ and $+$ to the feature vector $x$, respectively. Consequently, the strategy $r(x)$ be represented as the triplet $(p(x), q(x), 1 - p(x) - q(x))$. Let $$\fk{U}(x;r,s) := \sum_{y \in \cg{Y}\cup \Delta}\prob_r(Y=y|X'=s(x))U(y,x),$$ where $U$ is defined in \eqref{eqn:Utility}.
Since $r$ is a randomized classifier, the sender's payoff for the receiver's strategy $r$ at the feature vector $x$ is modified to
\begin{equation*}
	\mathbb{U}(x; r, s) := \fk{U}(x;r,s)-|s(x),x|.
\end{equation*}
We further assume that $0 \leq \up = \um = 2\uu \ll \min\{\h,1-\h\}$. Let the error function $T(x; r, s, h): \cg{X} \to [0,1]$ be defined as:
\begin{equation}\label{eqn:error function}
    T(x; r, s, h) := \prob_{r} \{Y\neq h(x) \mid X' = s(x)\}.
\end{equation}
The error function quantifies the error made at each feature vector $x$ by the sender strategy $s$ and the receiver strategy $r$ under the true function $h$. Since $r$ is a randomized classifier, the receiver's payoff is modified to
\begin{equation}\label{eqn:Receiver_average_error}
	\cg{E}(r, s; h) := \int_\cg{X} T(x; r, s, h) \, dx.
\end{equation}

\begin{notn}
    We sometimes drop some of the arguments from $T(\cdot, \cdot, \cdot, \cdot)$ when it is apparent from the context.
\end{notn}

For a receiver strategy $r$, the sender's worst-case lazy best response is defined as follows.

\begin{defn}[Worst-case Lazy Best Response]\label{def:Lazy_best_response}
Let $r$ be a receiver strategy. Let the set $V(x; r)$ be given as
	\[
	V(x; r) \in \arg\max_{x' \in \cg{X}}\Big(\sum_{y \in \cg{Y} \cup \Delta} \prob_r(Y=y|X'=x')(U(y,x)-|x'-x|)\Big).
	\]
	Let $\bar{B}(r)$ be the set of sender strategies $s$ satisfying
	\begin{equation}
	\begin{aligned}
		s(x) = \arg\min_{x' \in V(x; r)} |x'-x| \quad \text{ for all } x \in \cg{X}.
	\end{aligned}
	\end{equation}
	A sender strategy $s_r$ is called a worst-case lazy best response to a receiver strategy $r$ if
	\begin{equation*}
		s_r := \arg\min_{s' \in \bar{B}(r)} \cg{E}(r, s'; h),
	\end{equation*}
	and there does not exist $s' \in \bar{B}(r)$ such that
	\[
	T(x; r, s') \leq T(x; r, s_r) \quad \forall x \in \cg{X}.
	\]
\end{defn}
Next, we define the Stackelberg equilibrium strategy.
\begin{defn}
	Let $h$ be the true function. A strategy $(r^\star_h,s_{r^\star_h})$ is  a Stackelberg equilibrium if
	\[r^\star_{h} := \arg\max_{r \in \msf{R}}\cg{E}(r, s_r; h)\] and
    $s_{r^\star_h}$ is its worst-case lazy best response.
\end{defn}
\begin{notn}
    Sometimes, we refer $r^\star$ itself as Stackelberg equilibrium if $s_{r^\star}$ and $h$ are apparent from the context.
\end{notn}

There are several reasons for considering the setup where the sender adopts the worst-case lazy best response. \textcolor{red}{First, the existence of a Stackelberg equilibrium under the worst-case best response is hard to guarantee.} This stands in stark contrast to the worst-case lazy best response, where we explicitly construct a Stackelberg equilibrium. Second, the error in the Stackelberg equilibrium under the worst-case lazy best response provides a lower bound on the error of the Stackelberg equilibrium (if it exists) under the worst-case best response. This, combined with other results, allows us to calculate the error of the Stackelberg equilibrium under the worst-case best response without having to explicitly construct the equilibrium itself.
The following proposition and Corollary \ref{cor:Stackelberg worst case best response error} illustrates these points.

\begin{proposition}\label{prop:comparing sender strategy types}
    Let $r \in \msf{R}$ be a receiver strategy such that its worst-case best response and worst-case lazy best response exist. Let $s$ denote the worst-case best response to $r$, and let $s'$ denote the worst-case lazy best response to $r$. Then, the receiver's error satisfies
    \[
    \cg{E}(r, s; h) \geq \cg{E}(r, s'; h).
    \]
\end{proposition}

\begin{proof}
    By definition, both $s$ and $s'$ are best responses to $r$. Suppose
    \[
    \cg{E}(r, s'; h) > \cg{E}(r, s; h).
    \]
    But this contradicts the assumption that $s$ is the worst-case best response, as it implies there exists another strategy $s'$ that results in a higher error for the receiver.
    Therefore, it follows that
    \[
    \cg{E}(r, s; h) \geq \cg{E}(r, s'; h),
    \]
    as claimed.
\end{proof}


\begin{notn}
In this subsection, $s_r$ denotes the worst-case lazy best response, not the worst-case best response unless mentioned otherwise.
\end{notn}
The next theorem gives the structure of a Stackelberg equilibrium.

\begin{thm}\label{thm:Stochastic Stackelberg equilibrium}
Let $h$ be the true function and let $\ib := \h - \uu$ and $\jb := \h + \uu$. Then the receiver strategy
\begin{equation}\label{eqn:stochastic_stackelberg}
		r(x;\ib,\jb) := \begin{cases}
			(1,0,0) & x \in [0, \ib], \\
			(1-\frac{x-\ib}{2\uu},\frac{x-\ib}{2\uu},0) & x \in (\ib,\jb), \\
			(0,1,0) & x \in [\jb,1],
		\end{cases}
	\end{equation}
is a Stackelberg equilibrium and $\cg{E}(r) = \frac{\uu}{2}$.
\end{thm}



% \begin{proof}
%   The proof is given in Proof \ref{proof:Stackelberg worst case best response error} in the Appendix.
% \end{proof}




% \begin{proof} The proof is given in Proof \ref{proof:delta null wlog} in the Appendix.
% \end{proof}



% \begin{proof} The proof is given in Proof \ref{proof:error is half} in the Appendix.

% \end{proof}

% We next show that if a feature vector $x'$ is the worst-case lazy best response of some other feature vector $x$ and both have the same true label, then $j$ is also the worst-case lazy best response of itself.





% Let $r \equiv (p,q,1-p-q) \in \msf{R}$ be the receiver strategy. For feature vectors $x_1$ and $x_2 \in \cg{X}$ such that $x_1 \leq \h < x_2$, we calculate the contribution of error from the intervals $[x_1, x_1 - 2\uu q(x_1)]$ and $[x_2, x_2 + 2\uu q(x_2)]$.




% \begin{proof} The proof is given in Proof \ref{proof:Error more than half} in the Appendix.

% \end{proof}

% We now finally prove Theorem Theorem \ref{thm:Stochastic Stackelberg equilibrium}.

\begin{proof}
    From Lemma \ref{lem:error is half}, we have $\cg{E}(r) = \frac{\uu}{2}$. Furthermore, from Lemma \ref{lem:Error more than half}, we have $\cg{E}(r') \geq \frac{\uu}{2}$ for all $r' \in \msf{R}$. Hence, $r$ is a Stochastic Stackelberg equilibrium.
\end{proof}



	\subsection{Examples and comparison}\label{subsec:Examples and comparision}

	\input{Examples and Comparison}

	\section{Learning the Stackelberg equilibrium}\label{sec:Learning problem}

	In this section, we address the learning problem, where the receiver is unaware of the true function and is limited to employing only deterministic strategies. Let $\cg{F} \subseteq \{f \mid f:\cg{X} \rightarrow \cg{Y}\}$  denote a hypothesis class, and let  $$ l(x,y; f) := \mathbbm{1}_{f(x) \neq y} $$
be the \textit{zero-one loss function}, where $x \in \cg{X}, y \in \cg{Y}$ and $f \in \cg{F}$.
Let $(X,Y)$ be a random variable over $(\cg{X},\cg{Y})$ and $D$ be its corresponding probability law.
Let $L_{D}(f)$ denote the \textit{risk function} corresponding to the zero-one loss function and be defined as 
\begin{equation}\label{eqn:risk function}
    L_D(f) :=
    \E_{D}[l(X,Y;f)].
\end{equation}
Suppose $Y = h(X)$, where $h$ is the true function. Then we modify the definition of risk function given in \eqref{eqn:risk function}.
Let $X$ random variable over $\cg{X}$ and $D$ be it corresponding probability law. Then $L_D(f;h)$ denotes the \textit{$h$-risk function} corresponding to the zero-one loss function and be defined as 
\begin{equation}
L_D(f;h) :=   \E_{D}[l(X,h(X);f)].
\end{equation}
In the absence of strategic interactions between the sender and the receiver, a hypothesis class is PAC (Probably Approximately Correct) learnable if it satisfies the following definition (which is the same as Definition 3.1 in \cite{shalev-shwartz_ben-david_2014}).

\begin{defn}\label{def:PAC learning}
A hypothesis class $\cg{F}$ is PAC learnable if there exists a function $m_{\cg{F}}: (0,1)^2 \to \mathbb{N}$ and a learning algorithm $A: \cup_{n=1}^\infty (\cg{X} \times \cg{Y})^n \to \cg{F}$ such that for every $(\eps, \delta) \in (0,1)^2$, for every true function $h \in \cg{F}$ and every probability law $D$ over $\cg{X}$, the following holds:
Let $S = \{(x_i, h(x_i))\}_{i=1}^m$ is a sample where $x_i$'s are drawn i.i.d. from  $D$ and $m \geq m_{\cg{F}}(\eps,\delta)$. Then
\[
D^{\otimes m}\left( L_{D}(A(S);h)  < \eps \right) > 1 - \delta.
\]
\end{defn}
There exists a more general notion of PAC learning called agnostic PAC learning which is defined in the following definition (which is the same as Definition 3.3 in \cite{shalev-shwartz_ben-david_2014}).
\begin{defn}\label{def:agnostic PAC learning}
A hypothesis class $\cg{F}$ is agnostic PAC learnable if there exists a function $m_{\cg{F}}: (0,1)^2 \to \mathbb{N}$ and a learning algorithm $A: \cup_{n=1}^\infty (\cg{X} \times \cg{Y})^n \to \cg{F}$ such that for every $(\eps, \delta) \in (0,1)^2$ and every probability law $D$ over $\cg{X} \times \cg{Y}$, the following holds: Let $S = \{(x_i, y_i)\}_{i=1}^m$ be a sample where 
$(x_i, y_i)$'s are drawn i.i.d. from  $D$ and 
$m \geq m_{\cg{F}}(\eps, \delta)$. Then
\[
D^{\otimes m}\left( L_{D}(A(S)) - \inf_{f \in \cg{F}} L_{D}(f) < \eps \right) > 1 - \delta.
\]
\end{defn}
An algorithm $A$ is a \textit{successful PAC/agnostic PAC learner} of $\cg{F}$ if $\cg{F}$ is PAC/agnostic PAC learnable and $A$ is the corresponding learning algorithm defined in Definition \ref{def:PAC learning}/\ref{def:agnostic PAC learning}. Consider the learning algorithm Empirical risk minimization(ERM) 
\[\text{ERM}:\cup_{n=1}^\infty(\cg{X}\times \cg{Y})^n \rightarrow \cg{F} \text{defined as} 
\text{ERM}(\{(x_i, y_i)_{i=1}^m\}) := \arg\min_{f \in \cg{F}} \frac{1}{m} \sum_{i=1}^m l(x_i, y_i; f).
\]
The following theorem shows that a hypothesis class is PAC/agnostic PAC learnable if and only if it is learnable by an ERM algorithm.

\begin{thm}[Fundamental theorem of statistical learning]\label{thm:Fundamental theorem of statistical learning}
Let  $\cg{F} \subseteq \{f \mid f:\cg{X} \rightarrow \cg{Y}\}$ be a hypothesis class. Then, the following statements are equivalent:

\begin{enumerate}[label=\alph*)]
    \item $\cg{F}$ is agnostic PAC learnable.
    \item $\cg{F}$ is PAC learnable
    \item Any ERM rule is a successful PAC learner for $\cg{F}$.
    \item $\cg{F}$ has finite VC dimension.
\end{enumerate}
Furthermore, for a given $(\eps,\delta) \in (0,1)^2$ and  VC dim$(\cg{F}) = d$, the sample complexity $m_{\cg{F}}$ is bounded as follows:
\begin{equation}\label{eqn:sample complexity:thm:Fundamental theorem of statistical learning}
C_1\frac{d + \log\left(\frac{1}{\delta}\right)}{\epsilon^2} \leq m_{\cg{F}} \leq C_2\frac{d + \log\left(\frac{1}{\delta}\right)}{\epsilon^2},
\end{equation}
where $C_1, C_2 > 0$ are constants.
\end{thm}

\begin{proof}
Refer to Theorems 6.7 and 6.8 in \cite{shalev-shwartz_ben-david_2014}.
\end{proof}

We now establish a similar learning criterion for a hypothesis class $\cg{F}$ in the context of strategic classification. Recall that for a true function $f \in \cg{F}$ and probability law $D \in \fk{D}$, $r^\star_{f,D}$ denotes the Stackelberg equilibrium as defined in Definition \ref{def:Stackelberg equilibrium}. Define $\widetilde{\cg{F}} := \{r^\star_{f,D} \circ s_{r^\star_{f,D}} : f \in \cg{F}, D \in \fk{D}\}$, referred to as the \emph{Stackelberg hypothesis class}, which represents the space of receiver Stackelberg equilibria derived from $\cg{F}$ and $\fk{D}$. A hypothesis class $\cg{F}$ is strategically learnable if it satisfies the following definition.

\begin{defn}[Strategically learnable class]\label{def:Strategically learnable class}
A hypothesis class $\cg{F}$ is strategically learnable if there exists a function $m_{\cg{F}}: (0,1)^2 \to \mathbb{N}$ and a learning algorithm $A: \cup_{n=1}^\infty (\cg{X} \times \cg{Y})^n \to \widetilde{\cg{F}}$ such that for any true function $h \in \cg{F}$, any probability law $D \in \fk{D}$, and any $(\eps, \delta) \in (0,1)^2$, the following holds: Let $S = \{(x_i, h(x_i))_{i=1}^m\}$ be a sample where $x_i$ are drawn i.i.d. from $D$ and $m \geq m_{\cg{F}}(\eps, \delta),$
then
\begin{equation}\label{def: strategically learnable}
D^{\otimes m}\Big( \cg{E}(A(S); h, D) - \cg{E}(r^\star_{h, D}; h, D) < \eps \Big) > 1 - \delta.
\end{equation}
The function $m_{\cg{F}}(\eps, \delta)$ is called the sample complexity of the algorithm.
\end{defn}


Our definition of strategic learnability is different from \cite{sundaram2021pac} because we let the algorithm search over the Stackelberg equilibrium hypothesis class (i.e., over $r_f, f \in \mathcal{F}$) rather than the original hypothesis class (i.e., over $f \in \mathcal{F}$). 


We propose two algorithms for learning a strategic classifier: \textit{vanilla ERM} and \textit{strategy aware ERM} (defined in Section \ref{subsec:Vanilla ERM} and \ref{subsec:Stragically aware ERM} respectively). We show that a hypothesis class is strategically learnable if and only if it is learnable under the strategy-aware ERM algorithm. 
However, there exist strategically learnable hypothesis classes which are not
learnable under the vanilla ERM algorithm.

Assuming the hypothesis class is learnable by each algorithm, we compare their ease of learning (compared in Section \ref{subsec:comparision}). This is controlled by the VC dimension of the hypothesis class, which determines the sample complexity in PAC learning. We observe that when the hypothesis class is complex (i.e., contains functions with many decision boundaries), strategy-aware ERM makes learning easier, as it can reduce the VC dimension of the effective learning space. Conversely, for simpler hypothesis classes, vanilla ERM is often the the more direct and efficient approach.


\subsection{Some results on hypothesis class with multiple decision boundaries}

A point is called a \textit{decision boundary}, if it assigns different labels to its left and right neighbours. 
The hypothesis class we consider for our learning algorithms comprises of functions with multiple decision boundaries, requiring us to extend the results established in Theorem \ref{thm:Deterministic Stackelberg equilibrium strategy} (which only had one decision boundary). To achieve this, we examine a generalized version of a decision problem where the true function has multiple decision boundaries. Let $\cg{X} = [0,1]$ represent the feature vector space, and let $\cg{Y} = \{-1, +1\}$ be the label space, and $D \in \fk{D}$ a probability law on $\cg{X}$. Decompose $\cg{X}$ as $\cg{X} = O_1 \cup O_2 \cup \cdots \cup O_{d-1}$, where each $O_n$ is a disjoint interval defined as $O_n = [o_{(n)}, o^{(n)})$ for $n < d-1$, $O_{d-1} = [o_{d-1}, 1]$, and $o^{(n)} = o_{(n+1)}$. Define the true function $h(x)$ as follows:
\begin{equation}\label{eqn:multiple boundary true function}
    h(x) = \begin{cases}
        \minus & \text{if } x \in O_n, \; n \text{ is even}, \\
        \plus & \text{if } x \in O_n, \; n \text{ is odd}.
    \end{cases}
\end{equation}
Let the cost kernel be $\phi(x) = x$, and let $\Delta$ denote the penalty label (as considered in Section \ref{sec:Problem formulation}). Define the utility function $U(y,x)$ as:
\begin{equation}\label{eqn:Utility2}
    U(y,x) := \begin{cases}
        \uu & \text{if } x \in O_n, \; y = \plus, \; n \text{ is even}, \\
        \uu & \text{if } x \in O_n, \; y = \minus, \; n \text{ is odd}, \\
        0 & \text{if } y = h(x), \\
        -\infty & \text{if } y = \Delta.
    \end{cases}
\end{equation}
The structure of the Stackelberg equilibrium strategy generalizes Theorem \ref{thm:Deterministic Stackelberg equilibrium strategy} (when $\up = \um$) and is stated in the following proposition.

\begin{prop}\label{prop:Deterministic Stackelberg equilibrium for multiple decision boundaries}
	Let $D$ be a probability law on $\cg{X}$, and let $h(x)$ be the true function. For any $n$, let $\Theta_n$ be an interval such that $\Theta_n \subseteq O_n$. Define the receiver strategy $r$ as follows:
	\begin{equation}\label{eqn:stack multiple}
		r(x;\Theta_1,\cdots,\Theta_{d-1}) = \begin{cases}
			h(x) & x \in \cup_{n=1}^{d-1} \Theta_n, \\
			\Delta & x \not\in \cup_n \Theta_n.
		\end{cases}
	\end{equation}
	Then their exist $\Theta_n$'s, such that $r$ is a Stackelberg equilibrium, and $r \circ s_r$ is of the form:
	\begin{equation*}
		r(s_r(x)) = \begin{cases}
			h(x) & x \in \cup_n \Theta_n, \\
			h(x)^c & x \not\in \cup_n \Theta_n.
		\end{cases}
	\end{equation*}
\end{prop}


\begin{proof} The proof of the Proposition is given in Proof \ref{proof:Deterministic Stackelberg equilibrium for multiple decision boundaries} in the Appendix.


\end{proof}

Our next theorem establishes the VC dimension of a hypothesis class containing hypotheses with multiple decision boundaries.

\begin{prop}\label{prop: VC dimension}
    Let the hypothesis class $\mathcal{F} \subset \{f| f:\cg{X} \rightarrow \cg{Y}\}$ be a collection of all functions with at most $d-1$ decision boundaries. Then VC dim$(\mathcal{F}) = d$.
\end{prop}

\begin{proof}
Suppose VC dim $(\mathcal{F}) < d$. 
This means that there exist $(x_1, y_1), \dots, (x_d, y_d) \in \cg{X} \times \cg{Y}$ such that for each $f \in \mathcal{F}$, there exists an index $i_f$ such that $f(x_{i_f}) \neq y_{i_f}$. Without loss of generality, we assume $x_1 < x_2 <\cdots < x_d$.
Let $J := \{j : y_j \neq y_{j+1}, j < d\}$. Let $j_i$ denote the $i$-th smallest element of $J$.
Now consider the following hypothesis $f$ defined as 
\[
f(x) := 
\begin{cases}
    y_{j_1} &\text{if} x \in [0,\frac{x_{j_1}+x_{j_2}}{2}),\\
    y_{j_{i}} &\text{if} x \in [\frac{x_{j_i} - x_{j_{i-1}}}{2}, \frac{x_{j_{i+1}} - x_{j_i}}{2}),\\    
    y_{j_{|J|}}, & \text{if} x \in [\frac{j_{|J|-1} + j_{|J|}}{2},1].
\end{cases}
\]
The hypothesis $f$ has at most $d-1$ decision boundaries and hence $f \in \mathcal{F}$. Furthermore, $f(x_k) = y_k$ for all $k \in \{1, \dots, d\}$ by construction. But this contradicts the assumption. Hence, VC dim $(\mathcal{F}) \geq d$.
Now, consider $(x_1, y_1), \dots, (x_{d+1}, y_{d+1})$, where  $y_j \neq y_{j+1}$ for all $j \leq d$. This means that any function $f'$ with the property that $f'(x_i) = y_i$ must have at least $d$ decision boundaries. But functions in $\mathcal{F}$ have at most $d-1$ decision boundaries. Therefore, $\mathcal{F}$ cannot shatter $\{x_i\}_{i=1}^{d+1}$. Hence, VC dim$(\mathcal{F}) \not\geq d+1$. Therefore, VC dim $(\mathcal{F}) = d$.

\end{proof}


	\subsection{Vanilla ERM}\label{subsec:Vanilla ERM}
%
	One intuitive algorithm for the receiver to learn the strategic classifier is to estimate the true function from the sample data. The receiver then plays the screening game under this estimated function, treating it as the true function, and outputs the corresponding Stackelberg equilibrium. The vanilla ERM algorithm expresses this idea.
Let $D \in \fk{D}$ be a probability law on $\cg{X}$, $h$ be the true function and let $S = \{(x_i, h(x_i))_{i=1}^m\}$ be a sample where $x_i$ are drawn i.i.d. from $D$. Let $D(S)$ be the empirical law of sample $S$ defined as:
\begin{equation}\label{eqn:empirical measure}
D(S)(B) := \frac{1}{m}\sum_{i=1}^m \mathbbm{1}_{x_i \in B}  
\end{equation}
where $B$ is a Borel measurable set in $\cg{X}$. 

\begin{center}
	\begin{algorithm}[H]\label{algo:Vanilla ERM}
		\SetAlgoLined
		\DontPrintSemicolon
		\KwData{$\text{Hypothesis class }\cg{F}, \text{ Sample } S = \{x_i,h(x_i)\}_{i=1}^m$,\\ $\text{Zero-one loss function } l_h(f,x) := \mathbbm{1}_{\{f(x) = h(x)\}}(x)$}
		
		\KwResult{Hypothesis $\hat{r}_{erm}$}
		\BlankLine
		\tcp{Estimation}
		$A(S) := \arg\min_{f \in \cg{F}} \sum_{i=1}^n \mathbbm{1}_{\{f(x_i) \neq y_i\}}$\\
		$D(S) :=\text{Empirical law }$ \\
		\BlankLine
		\tcp{Decision}
		$\hat{r} := r^\star_{A(S),D(S)}$, i.e., Stackelberg equilibrium under true function $\hat{h}$ and probability law $D(S)$ \;
		\BlankLine
		\caption{Vanilla ERM algorithm}
	\end{algorithm}
\end{center}

The Algorithm \ref{algo:Vanilla ERM} is initialized by providing a hypothesis class $\cg{F}$ and a sample $S$. In the first step, the algorithm calculates $A(S)$, which estimates the true function $h$ from the sample $S$ via the ERM subroutine. In the second step, it calculates the empirical law $D(S)$ of sample $S$, which is an estimate of the input law $D$ on $\cg{X}$ from which the $x_i$'s in $S$ were chosen. In the third step, it calculates the Stackelberg equilibrium $r^\star_{A(S),D(S)}$ (refer Definition \ref{def:Stackelberg equilibrium}) under the (estimated) true function $A(S)$ and the probability law $D(S)$ on $\cg{X}$, as given in Theorem \ref{thm:Deterministic Stackelberg equilibrium strategy}. 

We show that if the hypothesis class $\cg{F}$ has finite VC dimension, then it is strategically learnable under the Algorithm \ref{algo:Vanilla ERM}.

\begin{thm}\label{thm:Startegic learning for vanilla ERM}
Let $\mathcal{F} \subseteq \{ f \mid f: \cg{X} \rightarrow \cg{Y} \}$ be a hypothesis class consisting of all functions with at most $d-1$ decision boundaries, where $d <\infty$. Let $A: \cup_{n=1}^\infty (\cg{X} \times \cg{Y})^n \rightarrow \mathcal{F}$ denote the Algorithm \ref{algo:Vanilla ERM} (vanilla ERM algorithm) and let $B: \cg{F} \rightarrow \widetilde{\cg{F}}$ map a hypothesis to its corresponding Stackelberg equilibrium. Then $\cg{F}$ is strategically learnable (cf. Definition \ref{def:Strategically learnable class}) under $B \circ A$.
\end{thm}

We prove the theorem with the help of the following two lemmas. The first lemma states that for a fixed true function $f \in \cg{F}$, the classification errors of the Stackelberg equilibrium strategies $r^\star_{f,D}$ and $r^\star_{f,D(S)}$ are similar for all $D \in \fk{D}$.

\begin{lem}\label{lem:same true function}
Let $(\epsilon, \delta) \in (0,1)^2$. Let $f \in \cg{F}$ be a true function with upto $d-1$ decision boundaries, and let $D \in \fk{D}$ be a probability law on $\cg{X}$. Let $S = \{(x_i, h(x_i))\}_{i=1}^m$ be a sample of size $m$, where the $x_i$'s are drawn i.i.d. according to a probability law $D$, and let $D(S)$ be the empirical law of the sample $S$. Let $r^{\star}_{D(S)}$ and $r^{\star}_{D}$ be the Stackelberg equilibrium under the true function $f$ for probability law $D(S)$ and $D$ respectively. Then, there exists an $\bar{m}(\eps,\delta)$ such that for all $m > \bar{m}(\eps,\delta)$ the following holds:
    \begin{equation*}
        D^{\otimes m}\big(\cg{E}(r^\star_{f,D(S)};f, D) - \cg{E}(r^\star_{f,D};f, D) < 2(d-1)\epsilon\big) > 1 - 2e^{-2m\epsilon^2}.
    \end{equation*}
\end{lem}

\begin{proof} The proof is given in Proof \ref{proof:same true function} in the Appendix.
\end{proof}

The next lemma states that for a fixed probability law $D \in \fk{D}$, the classification errors of the Stackelberg equilibrium strategies $r^\star_{A(S),D}$ and $r^\star_{h,D}$ are similar. 
\begin{rem}
Notice that $r_{A(S),D}$ is the Stackelberg equilibrium when the estimated true function $A(S)$ is considered the true function. To avoid confusion with the actual true function $h$, which generated the samples, we will refer to $h$ as the \textit{ground true function} and $A(S)$ as the \textit{estimated true function} whenever there is potential for ambiguity.
\end{rem}

\begin{lem}\label{lem:same distribution}
Let $(\epsilon, \delta) \in (0,1)^2$.
Let $\mathcal{F} \subseteq \{ f \mid f: \cg{X} \rightarrow \cg{Y} \}$ be a hypothesis class consisting of all functions with at most $d-1$ decision boundaries. Let $A: \cup_{n=1}^\infty (\cg{X} \times \cg{Y})^n \rightarrow \mathcal{F}$ denote the algorithm described in Algorithm \ref{algo:Vanilla ERM} (vanilla ERM algorithm).
For every ground true function $h \in \mathcal{F}$ and every probability law $D \in \fk{D}$ on $\cg{X}$, the following holds. Let $S = \{(x_i, h(x_i))\}_{i=1}^m$ be a sample of size $m$, where $x_i$ are drawn i.i.d. according to $D$. Let $r^\star_{A(S),D}$ and $r^\star_{h,D}$ be the Stackelberg equilibria under $D$ and the estimated true function $A(S)$ and the ground true function $h(x)$, respectively. Then there exists $\bar{m}(\eps,\delta)$ such that for all $m \geq \bar{m}(\eps,\delta)$, the following holds.

\begin{enumerate}[label=(\alph*)]
    \item \label{lem:same distribution (a)} For any receiver strategy $r \in \msf{R}$ and (arbitrary) true function $f \in \cg{F}$ ($s_r$ depends on $f$), we have \begin{equation*}
        D^{\otimes m}\big(|\mathcal{E}(r; h,D) - \mathcal{E}(r; A(S), D)| < \epsilon\big) > 1 - \delta,
    \end{equation*}

    \item \label{lem:same distribution (b)} For receiver Stackelberg equilibria $r^\star_{A(S),D}$ and $r^\star_{h,D}$, we have \begin{equation}\label{eqn:lem:vanilla ERM PAC bound}
    D^{\otimes m}\big(\mathcal{E}(r^\star_{A(S),D}; h,D) - \mathcal{E}(r^\star_{h,D}; h, D) < 2\epsilon\big) > 1 - \delta.
\end{equation}

Furthermore, the sample complexity $\underbar{m}(\eps,\delta)$ satisfies
\begin{equation}\label{eqn:sample complexity:lem:same distribution}
    \begin{aligned}
        C_1 \frac{d + \log\left(\frac{1}{\delta}\right)}{\epsilon} &\leq \underbar{m}(\eps,\delta) \leq C_2 \frac{d \log\left(\frac{1}{\epsilon}\right) + \log\left(\frac{1}{\delta}\right)}{\epsilon}
    \end{aligned}.
\end{equation}
Here $C_1, C_2 > 0$ are constants.
\end{enumerate}

\end{lem}


\begin{proof} The proof is given in Proof \ref{proof:same distribution} in the Appendix.
\end{proof}




Combining Lemmas \ref{lem:same distribution} and \ref{lem:same true function}, we now prove our theorem
\begin{proof}[Proof of Theorem \ref{thm:Startegic learning for vanilla ERM}]
Let $D \in \fk{D}$ be the probability law on $\cg{X}$ and let $(\eps,\delta) \in (0,1)^2$. Let $h \in \cg{F}$ be the ground true function and let $S = \{(x_i,h(x_i))\}_{i=1}^m$ be a $m$-sized sample where $x_i$'s are drawn i.i.d. from $D$. Let $m \geq m_\cg{F},\bar{m}(\eps,\delta),\underbar{m}(\eps,\delta)$ where $m_\cg{F}, \bar{m}(\eps,\delta)$ and $\underbar{m}(\eps,\delta)$ are defined in Theorem \ref{thm:Fundamental theorem of statistical learning}, and Lemma \ref{lem:same true function} and \ref{lem:same distribution} respectively.
Let $D(S)$ be the corresponding empirical distribution of $S$.
Let $r_1 := r^\star_{h,D}, 
r_2 := r^\star_{A(S),D},
r_3 := r^\star_{A(S),D(S)} = B(A(S)).$

Since VC dim$(\cg{F}) < \infty$,  $\cg{F}$ is PAC learnable from Theorem \ref{thm:Fundamental theorem of statistical learning}. Furthermore, from Theorem  \ref{thm:Fundamental theorem of statistical learning} we have $A(S)$ as the estimated true function such that 
\begin{equation}\label{eq:concentration bound3}
    D^{\otimes m}(L_{D}(A(S);h) < \epsilon) > 1 - \delta.
\end{equation}
Let $$\cg{S}_1 := \{S: L_D(A(S);h) < \eps\}$$ 
and 
$$\cg{S}_2 := \{S: \sup_{x \in [0,1]}|\psi(x) - \psi_m(x)| < \eps \},$$
where $\psi(x)$ and $\psi_m(x)$ is the distribution function of $D$ and $D(S)$ respectively.
Then from \eqref{eq:concentration bound3}, we get $D^{\otimes m}(\cg{S}_1) > 1-\delta$ and from \eqref{eqn:DKW inequality:lem: same true function}, we get $D^{\otimes m}(\cg{S}_2) > 1-2e^{-2m\eps^2}$. Let 
\[\cg{S} := \cg{S}_1 \cap \cg{S}_2.\]
Then, from union bound, we have $D^{\otimes m}(\cg{S}) >  1-2e^{-2m\eps^2} - \delta$.
Fix $S \in \cg{S}$.
Using Lemma  \ref{lem:same distribution} (a) by identifying $r = r_3$ and $f = A(S)$ as the (estimated) true function, we get 
\begin{equation}\label{eqn:r3 h to A(S)}
    \mathcal{E}(r_3; h, D) - \mathcal{E}(r_3; A(S), D) < \epsilon.
\end{equation}
Using Lemma \ref{lem:same true function} by identifying $f = A(S)$ as the (estimated) true function and $D$ as the probability law, we obtain
\begin{equation}\label{eq:thm:1st triangle}
    \mathcal{E}(r_3; A(S), D) - \mathcal{E}(r_2; A(S), D) < 2(d-1)\epsilon.
\end{equation}
Using Lemma  \ref{lem:same distribution} (a) by identifying $r = r_2$ and $f = A(S)$ as the (estimated) true function, we get
\begin{equation}\label{eqn:r2 A(S) to h}
    \mathcal{E}(r_2; A(S), D) - \mathcal{E}(r_2; h, D) < \epsilon.
\end{equation}
Finally, using Lemma \ref{lem:same distribution} $(b)$ by identifying $D$ as the probability law, we have
\begin{equation}\label{eq:thm:3rd triangle}
    \mathcal{E}(r_2; h, D) - \mathcal{E}(r_1; h, D) < \epsilon.
\end{equation}
Thus combining \eqref{eqn:r3 h to A(S)}, \eqref{eqn:r2 A(S) to h}, \eqref{eq:thm:1st triangle}, and \eqref{eq:thm:3rd triangle}, we get
\[
\mathcal{E}(r_3; h, D) - \mathcal{E}(r_1; h, D) < (2d + 1)\epsilon.
\]
This equation holds for all $S \in \cg{S}$.   But, we have $D^{\otimes m}(\cg{S}) > 1-\delta - 2e^{-2m\eps^2}$. Let $\delta' = \delta + 2e^{-2m\eps^2}$. Thus we have 
\[
D^{\otimes m}(\mathcal{E}(r_3; h, D) - \mathcal{E}(r_1; h, D) < (2d + 1)\epsilon) >  1- \delta'.
\]
The choice of $D$, $(\eps,\delta)$ and $h$ were arbitrary, therefore this shows that $\cg{F}$ is strategically learnable under $B \circ A$.
\end{proof}

The previous theorem only provides sufficient conditions for strategic learnability. We will now present an algorithm that is both a necessary and sufficient condition for a hypothesis class to be strategically learnable.










	\subsection{Strategy aware ERM}\label{subsec:Stragically aware ERM}
%
	The learning algorithm suggested in Algorithm \ref{algo:Vanilla ERM} does not account for strategic behaviour during the testing phase. This is a potential source of inefficiency in determining the strategic learnability of a hypothesis class. An alternate algorithm which takes into account the strategic behaviour and which we show is both a necessary and sufficient condition for strategic learnability is described as follows.
Let $\cg{F} \subseteq \{f \mid f:\cg{X} \rightarrow \cg{Y}\}$ be the hypothesis class consisting of all functions with at most $d-1$ decision boundaries. Corresponding to this hypothesis class, consider the hypothesis class $\widetilde{\cg{F}}$ defined as
\[
\widetilde{\cg{F}} := \{r \circ s_r \mid r := r^\star_{f,D}, D \in \fk{D}, f \in \cg{F}\}.
\]
We learn the Stackelberg equilibrium directly from $\widetilde{\cg{F}}$. The strategy-aware ERM algorithm expresses this idea.
\begin{center}
	\begin{algorithm}[H]\label{algo:Startegically aware ERM}
		\SetAlgoLined
		\DontPrintSemicolon
		\KwData{Hypothesis class $\cg{F}$, Sample $\{(x_i,y_i)\}_{i=1}^m$\\
			\text{Zero - one loss function} $l_h(r(s_r),x) := \mathbbm{1}_{r(s_r(x)) = h(x)}$}
		\KwResult{Hypothesis $\hat{r}$}
		\BlankLine
		\tcp{Decision}
		$\widetilde{\cg{F}} := \{r \circ s_r \mid r := r^\star_{f,D}, D \in \fk{D}, f \in \cg{F}\}$  \;
		\BlankLine
		\tcp{Estimation}
		$\hat{r}_{\textrm{serm}}\circ s_{\hat{r}_{\textrm{serm}}}:= \arg\min_{r \circ s_r \in \widetilde{\cg{F}}} \sum_{i=1}^m \mathbbm{1}_{\{r(s_r(x_i)) \neq y_i\}}$ \\
		
		\BlankLine
		
		\caption{Strategy aware ERM algorithm}
	\end{algorithm}
\end{center}
The algorithm \ref{algo:Startegically aware ERM} is initialised with a hypothesis class $\cg{F}$ and a sample $S$. In the first step, we modify the hypothesis class $\cg{F}$ to $\widetilde{\cg{F}}$. In the second step, we run the ERM subroutine on the modified hypothesis class $\widetilde{\cg{F}}$ and sample $S$ to get a hypothesis $\hat{r}_{\textrm{serm}}\circ s_{\hat{r}_{\textrm{serm}}}$. We now show that VC dim$(\widetilde{\cg{F}}) < \infty$ if and only if $\cg{F}$ is strategically learnable under Algorithm \ref{algo:Startegically aware ERM}. 

\begin{thm}[The fundamental theorem of strategic learning ]\label{thm:Fundamental theorem of strategic learning}
Let $\mathcal{F} \subseteq \{ f \mid f: \cg{X} \rightarrow \cg{Y} \}$ be a hypothesis class consisting of all functions with at most $d-1$ decision boundaries. Let $A: (\cg{X} \times \cg{Y})^n \rightarrow \widetilde{\cg{F}}$ denote the algorithm described in Algorithm \ref{algo:Startegically aware ERM}.  The $\cg{F}$ is strategically learnable under $A$ (refer Definition \ref{def:Strategically learnable class}) iff VC dim$(\widetilde{\cg{F}}) < \infty$.
\end{thm}

\begin{proof}
Let  $D \in \fk{D}$ be the probability law on $\cg{X}$ and let $(\eps,\delta) \in (0,1)^2$. Let $h \in \cg{F}$ be the true function and let $S = \{(x_i,y_i)\}_{i=1}^m$ be the $m$-sized sample where $x_i$'s are drawn i.i.d. from $D$. Let $\cg{E}(r,f) := \cg{E}(r,s_r,f,D)$. The Strategy aware ERM algorithm is similar to the vanilla ERM algorithm but applied to a different hypothesis class $\widetilde{\cg{F}}$ under the zero-one loss function $
		l_h(r \circ s_r, x) = \mathbbm{1}_{\{r(s_r(x)) = h(x)\}}$.
	We first prove the if part.
    Let $A(S)$ be the hypothesis obtained from the strategy aware ERM algorithm. Since VC dim$(\widetilde{\cg{F}}) < \infty$, we know that it is agnostic PAC learnable from Theorem \ref{thm:Fundamental theorem of statistical learning}. Furthermore, from Theorem \ref{thm:Fundamental theorem of statistical learning} we have for $m \geq \bar{m}(\eps,\delta)$,
	\begin{equation}\label{eqn:lem:startegic learnable}
	   \begin{aligned}
			D^{\otimes m}(\cg{E}(A(S), h) - \min_{r \circ s_r \in \widetilde{\cg{F}}} \cg{E}(r, h) < \epsilon) > 1 - \delta.
		\end{aligned}
	\end{equation}
	Note that by construction of $\widetilde{\cg{F}}$, we have
	\[
		\min_{f \in \cg{F}} \cg{E}(r_f, s_{r_f}, h, D) = \min_{r \circ s_r \in \widetilde{\cg{F}}} \cg{E}(r, s_r, h, D).
	\]
	Substituting this in \eqref{eqn:lem:startegic learnable}, we get with probability $1 - \delta$,
    \begin{equation}\label{eqn:thm:fundamental theorem of strategic learning}
        \begin{aligned}
            D^{\otimes m}(\cg{E}(A(S), h) - \min_{f \in \widetilde{\cg{F}}} \cg{E}(r, h) < \epsilon) > 1 - \delta.
        \end{aligned}
    \end{equation}
The choice of $D$, $(\eps,\delta)$ and $h$ were arbitrary. Hence \eqref{eqn:thm:fundamental theorem of strategic learning}
shows that $\cg{F}$ is strategically learnable.

\noindent Next, we prove the only if part. If $\cg{F}$ is strategically learnable, there exists an algorithm $A': (\cg{X} \times \cg{Y}) \rightarrow \widetilde{\cg{F}}$ such that 

\begin{equation}\label{eqn:thm:fundamental strategic learning bound}
D^{\otimes m}\Big( \cg{E}(A'(S); h, D) - \cg{E}(r^\star_{h, D}; h, D) < \eps \Big) > 1 - \delta
\end{equation}
holds.
Note that from the definition of $r^\star_{h,D}$, $\cg{E}(r^\star_{h,D};h,D) \leq \cg{E}(r;h,D)$ for all $r \in \widetilde{\cg{F}}$. Furthermore, from the definition of $\cg{E}(\cdot;h)$ and $L(\cdot;h)$, we have $\cg{E}(A(S),h) = L(A(S);h)$.
Therefore, \eqref{eqn:thm:fundamental strategic learning bound} reduces to 
\begin{equation}\label{eqn:them:fundamanetal strategic learning bound}
D^{\otimes m}\Big( L(A'(S);h) - \arg\min_{r \in \widetilde{\cg{F}}}L(r;h) < \eps \Big) > 1 - \delta.
\end{equation}
This shows that $\widetilde{\cg{F}}$ is agnostic PAC learnable under $A'(S)$. Therefores from the Theorem \ref{thm:Fundamental theorem of statistical learning} we have VC dim$(\widetilde{\cg{F}}) < \infty$.
\end{proof}

\begin{cor}\label{cor:SERM is more widely applicable than vanilla ERM}
	Let $\cg{F}$ be a hypothesis class used in Algorithm \ref{algo:Vanilla ERM} and $\widetilde{\cg{F}}$ be the hypothesis in Algorithm \ref{algo:Startegically aware ERM} obtained from $\cg{F}$. Then
	\[
		\text{VC dim}(\cg{F}) < \infty \implies \text{VC dim}(\widetilde{\cg{F}}) < \infty.
	\]
\end{cor}



We now demonstrate that certain hypothesis classes are not strategically learnable via Algorithm \ref{algo:Vanilla ERM} but become strategically learnable via Algorithm 
\ref{algo:Startegically aware ERM}.  Specifically, we show the existence of a hypothesis class $\mathcal{F}$ such that VC dim$(\mathcal{F}) = \infty$ but VC dim$(\tilde{\mathcal{F}}) < \infty$. This establishes that the converse of Corollary \ref{cor:SERM is more widely applicable than vanilla ERM} does not hold.

\begin{exg}\label{exg:counter example}
Consider the deterministic decision problem described in Section \ref{sec:Problem formulation}. Let $h$ be the true function, $D$ be the uniform probability law on $\cg{X}$ and $\um = \up = 0.1$. Let $\mathbf{f} \in [0,1)$, and let $F \subseteq ([\mathbf{f},1] \cap \mathbb{Q})$. Define the hypothesis $f_F: \cg{X} \to \cg{Y}$ as follows:
\[
f_F(x) :=
\begin{cases}
    \minus & x \in [0, \mathbf{f}), \\
    \plus & x \in F \setminus [0, \mathbf{f}), \\
    \minus & x \in F^c \setminus [0, \mathbf{f}) \text{ where} F^c = X \backslash F.
\end{cases}
\]
Now, consider the following hypothesis class:
\[
\begin{aligned}
    \mathcal{F} &:= \{f_F : \mathbf{f} \in (0,1), \, F \subseteq ([\mathbf{f}, 1] \cap \mathbb{Q})\}
\end{aligned}
\]
Then VC dim$(\mathcal{F}) = \infty$ because $\mathcal{F}$ can shatter every countable subset of $[\mathbf{f}, 1] \cap \mathbb{Q}$. We now claim that VC dim$(\widetilde{\cg{F}}) = 0$. Let $f_F \in \mathcal{F}$ be the true function. The receiver strategy $r(x) \equiv \minus$ is a Stackelberg equilibrium. This is because $\cg{E}(r) = 0$.
Thus, we have \( \widetilde{\cg{F}} = \{r(\cdot) \equiv \minus\} \), implying that the VC dimension of \( \widetilde{\cg{F}} \) is $0$. In conclusion, while \( \mathcal{F} \) is not strategically learnable by Algorithm \ref{algo:Vanilla ERM} due to its infinite VC dimension, it becomes strategically learnable by Algorithm \ref{algo:Startegically aware ERM}, as the hypothesis space is effectively reduced to one with finite VC dimension.
\end{exg}








	\subsection{Comparison between algorithms}\label{subsec:comparision}

	\input{Comparision between algorithms}

	\section{Conclusions and Limitations}

	\input{Conclusion}

	\section{Appendix: Omitted Proofs}

        \subsection{Omitted Proofs of Section III.B}

        We analyze the worst-case best response and error of a strategy $r \in \sr{R}_2$.

\begin{lem}\label{lem:constant strategy}
	Let $c \in \cg{Y}$ be a label and $r \in \sr{R}_2$ be a constant-label  receiver strategy $r(\cdot) \equiv c$. Then $s_r(x) = x$ and $E(r)^c = \{x: h(x) =c\}$ (here $^c$ denotes the set complement).
\end{lem}

\begin{proof}
	Since the receiver uses a constant label strategy, for a given $x \in \cg{X}$, the sender's payoff is $U(r(s_r(x)),x) - C(s_r(x),x) = U(c,x) - C(s_r(x),x)$.
	Since, the utility term remains fixed irrespective of the sender strategy, the sender will focus on minimizing the cost term. Consequently, $s_r(x) = x$ incurs the least cost. The receiver strategy $r$ correctly classify feature vectors whose true label is $c$, hence $E(r)^c = \{x:h(x = c)\}$.
\end{proof}


We begin by analyzing the worst-case best response and error of a strategy $r$ in $\sr{R}_1$.
\begin{lem}\label{lem:Best response r1}
	Let $r(\cdot;\ib,\jb) \in \sr{R}_1$.
	\begin{enumerate}[label=(\alph*)]
		\item If \[\phi(\h) - \phi(\jb) \leq \um + \phi(\ib) - \phi(\h),\]
			then $s_r(\cdot)$ is given by
		\begin{equation*}
			s_r(x) := \begin{cases}
				x & \text{if} x \in [0,\ib) \cup (\jb,1], \\
				\ib & \text{if} x \in (\h,\lb), \\
				\jb & \text{if}x \in [\ib,\h] \cup [\lb,\jb],
			\end{cases}
		\end{equation*}
		where $\lb := \phi^{-1}\Big(\frac{\phi(\ib) + \phi(\jb) + \um}{2}\Big)$. Furthermore, $E(r) = [\ib,\lb]$.
		\item If \[\phi(\h) - \phi(\jb) > \um + \phi(\ib) - \phi(\h),\]
		then $s_r(\cdot)$ is given by
		\begin{equation*}
			s_r(x) := \begin{cases}
				x & \text{if} x \in [0,\ib) \cup (\jb,1], \\
				\jb & \text{if}x \in [\ib,\jb],
			\end{cases}
		\end{equation*}
\noindent Furthermore $E(r) = [\ib,\h]$ and $\up \geq \phi(\h) - \phi(\ib) \geq \frac{\up + \um}{2}$.
	\end{enumerate}
\end{lem}


\begin{proof}\label{proof:Best response r1}
    We divide the task of finding $s_r$ over $\cg{X}$ into various subintervals of $\cg{X}$.\\
\textbf{Case 1: }{$x \in [0, \ib)$\\} The sender's payoff is given by $U(r(s_r(x)),x) - |\phi(s_r(x))-\phi(x)|$.
	In the utility term, $r(s_r(x))$ can result in three labels: $\minus, \plus$ and $\Delta$. The $\Delta$ label is unpreferable because it gives a $-\infty$ utility.
	Hence the sender will choose $s_r(x)$ such that $r(s_r(x)) = \minus$ or $ \plus$.
	For $r(s_r(x)) = \minus$, from the definition of $r$, $s_r(x) \in [0,\ib]$. Among these options, the feature vector which incurs the least cost to the sender is $x$.
	For $r(s_r(x)) = \plus$, from the definition of $r$, $s_r(x) \in [\jb,1]$. Among these options, the feature vector which incurs the least cost to the sender is $\jb$.
	Accordingly, the sender's payoff at $x$ is
		\[\senpayoff(x;r,s_r) := \begin{cases}
			\up + \phi(x) - \phi(s_r(x)) & \text{ if } s_r(x)= \jb \\
			0& \text{ if }s_r(x) = x.\\
		\end{cases}\]
	 Since $\phi$ is an increasing function and $r \in \sr{R}_1$ satisfies $\phi(\jb) - \phi(\ib) = \up > 0$, therefore for $x \in [0, \ib)$, we have $$\up + \phi(x) - \phi(s_r(x))  < \up + \phi(\ib) - \phi(\jb) = \up -\up  = 0.$$ Thus the payoff for modifying $x$ to $\jb$ is negative and hence $s_r(x) = x$.\\

\textbf{Case 2:}
		 $x \in (\h,\jb]$.\\ Arguing as in the Case $1$, the $s_r(x)$ can be equal to $\ib$ or $\jb$. Hence the sender's payoff at $x$ is
		\[\senpayoff(x;r,s_r) := \begin{cases}
			\phi(x) - \phi(\jb)  & \text{ if } s_r(x) = \jb \\
		\um + \phi(\ib) - \phi(x)& \text{ if }s_r(x)=\ib.\\
		\end{cases}\]
	If $\phi(\h) - \phi(\jb) > \um + \phi(\ib) - \phi(x)$, then it follows that $\phi(x) - \phi(\jb) > \um + \phi(\ib) - \phi(x)$. Hence $s_r(x) = \jb$.
	If $\phi(\h) - \phi(\jb) \leq \um + \phi(\ib) - \phi(\h)$, then their exists $\lb \in (\h,\jb]$ (same $\lb$ that is defined in the statement of this lemma) such that
	\[\phi(\lb) - \phi(\jb) = \um + \phi(\ib) - \phi(x).\] Hence, if $\phi(x) \geq \phi(\lb),$ then $s_r(x) = \jb$, and if $\phi(x) < \phi(\lb),$ then $s_r(x) = \ib$.
	Since $\phi$ is an increasing function, this implies that
	 \[s_r(x) = \begin{cases}
	 	\ib, \text{if} x \in (\h,\lb)\\
	 	\jb, \text{if} x \in [\lb,\jb].
	 \end{cases}\]\\


\textbf{Case 3: }
		$x \in [\ib,\h]$.\\ Arguing as in the Case $1$, the $s_r(x)$ can be equal to $\ib$ or $\jb$.
		Accordingly, the sender's payoff at $x$ is
		\[\senpayoff(x;r,s_r) := \begin{cases}
			\up + \phi(x) - \phi(\jb)  & \text{ if } s_r(x) = \jb \\
			\phi(\ib) - \phi(x)& \text{ if }s_r(x) = \jb.\\
		\end{cases}\]
 Since $$\up + \phi(x) - \phi(\jb)  \geq 0 \text{ and } \phi(\ib) - \phi(x) <0,$$ therefore $s_r(x) = \jb$. \\


\textbf{Case 4:}
		Let $x \in (\jb,1]$. Arguing as in the Case $1$, the $s_r(x)$ can be equal to $\ib$ or $x$. Accordingly, the sender's payoff at $x$ is
			\[\senpayoff(x;r,s_r) := \begin{cases}
			\um + \phi(\ib) - \phi(x) & \text{ if } s_r(x) = \ib \\
			0& \text{ if }s_r(x) = x.\\
		\end{cases}\]
		Since $\phi$ is an increasing function and $r \in \sr{R}_1$ satisfies $\phi(\jb) - \phi(\ib) = \up > 0$, therefore for $x \in [\jb,1]$, we have
		\[\um + \phi(\ib) - \phi(x) < \um + \phi(\ib) - \phi(\jb) < \um -\up \leq 0.\]
 		Therefore $s_r(x) = x$.

\noindent Hence combining all the cases, we get that if $\phi(\h) - \phi(\jb) \leq  \um + \phi(\ib) - \phi(x)$, then
\begin{equation*}
   s_r(x) := \begin{cases}
       x &\text{ for } x \in [0,\ib) \cup (\jb,1]\\
       \ib &\text{ for } x \in [\ib,\lb)\\
       \jb &\text{ for } x \in [\ib,\h] \cup [\lb,\jb],
   \end{cases}
\end{equation*}
and $E(r) = [\ib,\lb]$, where $g := \phi^{-1}(\frac{\phi(\ib) + \phi(\jb) + \um}{2})$ which proves part (a).\\
\noindent Similarly combining all the cases, we get if $\phi(\h) - \phi(\jb) > \um + \phi(\ib) - \phi(x)$, then
\begin{equation*}
   s_r(x) := \begin{cases}
       x &\text{ for } x \in [0,\ib) \cup (\jb,1]\\
       \jb &\text{ for } x \in [\ib,\jb],
   \end{cases}
\end{equation*}
and $E(r) = [\ib,\jb]$ which proves part (b).
\end{proof}

\begin{lem}\label{lem:helping lemma}
	Let $r$ be a receiver strategy. Let $h(x)$ be a true function, and let $x_1, x_2 \in E(r)^c$ ($^c$ denote the set complement) be feature vectors such that $x_1 \leq \h \leq x_2$. Let $s_r$ be the worst-case best response of $r$. Let $z_1,z_2 \in \cg{X}$ be such that $z_1 = s_r(x_1)$, $z_2 = s_r(x_2)$ and, $r(z_1) \neq r(z_2)$ and $r(z_1),r(z_2) \neq \Delta$. Now consider the following two sender strategies \begin{equation}\label{eqn:s hat and s bar}
		\hat{s}(t) :=
		\begin{cases}
			z_2 & \text{for } t = x_1  \\
			s_{r}(t) & \text{otherwise}
		\end{cases},\quad
		\bar{s}(t) :=
		\begin{cases}
			z_1 & \text{for } t = x_2 \\
			s_{r}(t) & \text{otherwise}.
		\end{cases}
	\end{equation}
Then
\[\mathbb{U}(x_1;r,s_r) > \mathbb{U}(x_1;r,\hat{s}) \text{and} \mathbb{U}(x_2;r,s_r) > \mathbb{U}(x_2;r,\bar{s}).\]
\end{lem}

\begin{proof}\label{proof:helping lemma}
    Since $s_{r}$ is a worst-case best response to $r$, we have $$\mathbb{U}(x; r, s_r) \geq \mathbb{U}(x; r, \hat{s}) \text{ and } \mathbb{U}(x; r, s_r) \geq \mathbb{U}(x; r, \bar{s}) \text{ for all } x \in \cg{X}.$$ Since $\hat{s}(t) = s_{r}(t)$ for $t \in \cg{X} \backslash  \{x_1\}$ and $\bar{s}(t) = s_{r}(t)$ for $t \in \cg{X} \backslash \{x_2\}$, we have
	\[\mathbb{U}(t; r, s_{r}) = \mathbb{U}(t; r, \hat{s}) \text{for} t \in \cg{X} \backslash \{x_1\} \text{ and } \mathbb{U}(t; r, s_{r}) = \mathbb{U}(t; r, \bar{s}) \text{for} t \in \cg{X} \backslash \{x_2\}.\]
	If $\mathbb{U}(x_1; r, s_r) = \mathbb{U}(x_1; r, \hat{s})$, then $\hat{s} \in B(r)$ (the best response set of $r$) and $E(r, \hat{s}) = E(r, s_{r}) \cup \{x_1\}$. This contradicts the definition of $s_{r}$. Hence, $\mathbb{U}(x_1; r, s_{r}) \neq \mathbb{U}(x_1; r, \hat{s})$. Similarly, if $\mathbb{U}(x_2; r, s_{r}) = \mathbb{U}(x_2; r, \bar{s})$, then $\bar{s} \in B(r)$ and $E(r, \hat{s}) = E(r, s_{r}) \cup \{x_2\}$. This contradicts the definition of $s_{r}$. Hence, $\mathbb{U}(x_2; r, s_{r}) \neq \mathbb{U}(x_2; r, \bar{s})$.
	Therefore, it follows that
	\begin{equation*}
		\mathbb{U}(x_1; r, s_{r}) > \mathbb{U}(x_1; r, \hat{s}) \quad \text{and} \quad \mathbb{U}(x_2; \rt, s_{r}) > \mathbb{U}(x_2; r, \bar{s}).
	\end{equation*}
	which implies
	\begin{equation*}
		-|\phi(z_1) - \phi(x_1)| > \up -|\phi(z_2) - \phi(x_1)| \text{and} -|\phi(z_2) - \phi(x_2)| > \um - |\phi(z_1) - \phi(x_2)|.
	\end{equation*}
\end{proof}

We say a strategy $r$ dominates another receiver strategy $\rt$ if $E(r,s_r,h,D) \subseteq E(\rt,s_{\rt},h,D)$. We now show that the family of receiver strategies $\sr{R}$ dominates every other receiver strategy.

\begin{lem}\label{lem:Deterministic domination of a family of strategies}
	Let $D \in \fk{D}$ be a probability law on $\cg{X}$ and let $h(x)$ be a true function. Let $\rt$ be any receiver strategy. Then there exists a strategy $r \in \sr{R}$ such that $\cg{E}(\rt) \geq \cg{E}(r)$.
\end{lem}

\begin{proof}(Proof of Lemma \ref{lem:Deterministic domination of a family of strategies})\label{proof:Deterministic domination of a family of strategies}\\
    Let $\ib := \sup (E(\rt)^c \cap [0,\h])$ and $\tilde{\lb} := \inf(E(\rt)^c \cap (\h,1])$ (here $^c$ represents the set complement). It follows that $\ib$ represents the supremum of all the feature vectors with true label $\minus$ that get correctly classified by $\rt$. Similarly, $\tilde{\lb}$ represents the infimum of all the feature vectors with true label $\plus$ that get correctly classified by $\rt$. Hence $E(\rt) \supset (\ib,\tilde{\lb}).$ Let $\jb \in \cg{X}$ be such that $\h < \jb \leq 1$ and $\phi(\jb) - \phi(\ib) = \up$. Such a feature vector $\jb$ exists due to the continuity of $\phi$.\\
	\textbf{Case 1: }\label{lem:Deterministic domination of a family of strategies:a}{$\ib = 0$:}\\
		In this case, $E(\rt) \supset (0,\h)$, i.e., it misclassifies all the feature vectors with the true label $\minus$. Hence, for a receiver strategy $r \in \sr{R}_2$ which assigns the majority label, $\cg{E}(r) \leq \cg{E}(\rt)$.


\textbf{Case 2: }\label{lem:Deterministic domination of a family of strategies:b}{$\tilde{\lb} = 1$:}\\
		In this case, $E(\rt) \supset (\h,1)$ i.e., it misclassifies all the feature vectors with the true label $\plus$. Hence, for a receiver strategy $r \in \sr{R}_2$ which assigns the majority label, $\cg{E}(r) \leq \cg{E}(\rt)$.



\textbf{Case 3: }{$\ib > 0, \tilde{\lb} < 1$ and $\phi(\tilde{\lb}) - \phi(\ib) \geq \frac{\up + \um}{2}$:}\\
From the assumption on $\tilde{\lb}$, it follows that $\h \leq \tilde{\lb} \leq \jb$. Now consider a receiver strategy $r(x;\ib,\jb) \in \sr{R}_1$. If $\phi(\h) - \phi(\jb) \leq \um + \phi(\ib) - \phi(\h)$, then from Lemma \ref{lem:Best response r1} (a), $E(r) = [\ib,\lb]$, where
$\lb := \phi^{-1}\Big(\frac{\phi(\ib) + \phi(\jb) + \um}{2}\Big)$.
Since $\phi(\lb) - \phi(\ib) =\frac{\up+\um}{2}$, therefore $\lb \leq \tilde{\lb}$.
Hence
$$E(r) = [\ib,\lb] \text{ and } E(\rt) \supseteq (\ib,\tlb)$$ which implies $\cg{E}(r) \leq \cg{E}(\rt)$.
If $\phi(\h) - \phi(\jb) > \um + \phi(\ib) - \phi(\h)$, then from Lemma \ref{lem:Best response r1} (b), $E(r) = [\ib,\h]$. Since $\tilde{\lb} \geq \h$, therefore we have
\[E(r) = [\ib,\h] \text{ and } E(\rt) \supset (\ib,\tlb).\]
Hence $\cg{E}(r) \leq \cg{E}(\rt)$.


\textbf{Case 4: }{$\ib > 0, \tilde{\lb} < 1$ and $\phi(\tilde{\lb}) - \phi(\ib) < \frac{\up + \um}{2}$:}\\
	We will show that this case does not occur. Let $\eps$ be such that
	$0 < 2\eps < \frac{\up + \um}{2} -  \phi(\tlb) + \phi(\ib)$. Let
	$\ib', \tlb' \in E(\rt)^c$ be such that $\ib' \leq \ib \leq \h < \tlb \leq \tlb'$ and
	\begin{equation}\label{eqn:hdash}
		\phi(\ib) - \phi(\ib') < \eps \quad \text{and} \quad \phi(\tlb') - \phi(\tlb) < \eps.
	\end{equation}
	Such feature vectors $\ib', \tlb'$ exist by the definition of $\ib$ and $\tlb$ and the continuity of $\phi$. Adding the inequalities in \eqref{eqn:hdash}, and using the definition of $\eps$, we get
	\begin{equation}\label{eqn:gap between g and h}
		\phi(\tlb') - \phi(\ib') < \frac{\up + \um }{2}.
	\end{equation}
	Let $y_1 := s_{\rt}(\ib')$ and $y_2 := s_{\rt}(\tlb')$. It follows that
	\begin{equation}\label{eqn:rt at y1 y2}
		\rt(y_1) = h(\ib') = \minus \quad \text{and} \quad \rt(y_2) = h(\tlb') = \plus.
	\end{equation}
	Furthermore, the sender's payoff at $\ib'$ and $\tlb'$ under the worst-case best response $s_{\rt}$ is
	\begin{equation}\label{eqn:y_1 best response payoff}
		\mathbb{U}(\ib'; \rt, s_{\rt}) = U(\rt(y_1), \ib') - |\phi(y_1) - \phi(\ib')| = - |\phi(y_1) - \phi(\ib')| \leq 0
	\end{equation}
	and
	\begin{equation}\label{eqn:y_2 best response payoff}
		\mathbb{U}(\tlb'; \rt, s_{\rt}) = U(\rt(y_2), \tlb') - |\phi(y_2) - \phi(\tlb')| = - |\phi(y_2) - \phi(\tlb')| \leq 0
	\end{equation}
	respectively.\\
Additionally applying Lemma \ref{lem:helping lemma} by taking $r = \rt, x_1 = \ib'$, $x_2 = \tlb'$, we get $z_1 = y_1, z_2 = y_2$ and,
	\begin{equation}\label{eqn:s hat bad}
	-|\phi(y_1) - \phi(\ib')| > \up -|\phi(y_2) - \phi(\ib')| \text{and} -|\phi(y_2) - \phi(\tlb')| > \um - |\phi(y_1) - \phi(\tlb')|.
\end{equation}



\begin{claims}{$\rt(x) = \Delta$ for $x \in [\ib', \tlb']$.}\\
	Let $\kb^+, \kb^- \in \cg{X}$ be such that $\kb^- < \ib' < \kb^+$, $\phi(\kb^+) - \phi(\ib') = \up$ and $\phi(\ib') - \phi(\kb^-) = \up$. Such feature vectors $\kb^-, \kb^+$ exist due to the continuity of $\phi$. We now show that
	\begin{equation}\label{eqn:ordering g'}
		\kb^- < \ib' < \tlb' < \kb^+.
	\end{equation}
	To show this, we only need to prove that $\tlb' < \kb^+$. Recall from the assumption of \textbf{IV} that $\up \geq |\um|$. From this, we have
	\begin{equation}\label{eqn:gap}
		\phi(\kb^+) - \phi(\tlb') = \phi(\kb^+) - \phi(\ib') + \phi(\ib') - \phi(\tlb') > \up - \frac{\up + \um}{2} = \frac{\up - \um}{2} \geq 0,
	\end{equation}
	where we have used the definition of $\kb^+$ and \eqref{eqn:gap between g and h} to lower bound the middle term. Hence, $\phi(\kb^+) > \phi(\tlb')$. Since $\phi$ is strictly increasing, we have $\kb^+ > \tlb'$ as required. \\
	Now to prove the claim, we first show that $\rt(x) \neq \plus$ for $x \in [\kb^-, \kb^+]$. Suppose there exists $\ab \in [\kb^-, \kb^+]$ such that $\rt(\ab) = \plus$.
	Consider a sender strategy $s$ defined as
	\[
	s(t) = \begin{cases}
		\ab & \text{for } t = \ib' \\
		s_r(t) & \text{otherwise.}
	\end{cases}
	\]
	Then, at the feature vector $\ib'$, the sender can map $\ib'$ to $s(\ib') = \ab$ and obtain a payoff
	\[
	\mathbb{U}(\ib'; \rt, s) = U(\rt(s(\ib')), \ib') - |\phi(s(\ib')) - \phi(\ib')| =
	U(\plus, \ib') - |\phi(\ab) - \phi(\ib')| \geq \up - \up = 0.
	\]
	Using \eqref{eqn:y_1 best response payoff}, it follows that the sender obtains a better payoff at $\ib'$ by playing $s$ instead of $s_{\rt}$. But this contradicts the fact that $s_{\rt}$ is the worst-case best response to $\rt$. Hence, such a sender strategy $s$ cannot exist, which implies that
	\begin{equation}\label{eqn:rtx neq plus}
		\rt(x) \neq \plus \quad \text{for all } x \in [\kb^-, \kb^+].
	\end{equation}
	We next show that $\rt(x) \neq \minus$ for $x \in [\ib', \tlb']$. Suppose there exists $\bbb \in [\ib', \tlb']$ such that $\rt(\bbb) = \minus$.
	Using \eqref{eqn:rtx neq plus} and $\rt(y_2) = \plus$, we have $y_2 \notin [\kb^-, \kb^+]$. Since $\tlb' \in [\kb^-, \kb^+]$ from \eqref{eqn:ordering g'}, we have
	\[
	|\phi(y_2) - \phi(\tlb')| \geq \min\{\phi(\tlb') - \phi(\kb^-), \phi(\kb^+) - \phi(\tlb')\}.
	\]
	We now lower bound the right-hand side. From \eqref{eqn:gap}, we have $\phi(\kb^+) - \phi(\tlb') > \frac{\up - \um}{2}$. Furthermore, from \eqref{eqn:ordering g'}
	and the definition of $\kb^-$, we have  $\phi(\tlb') - \phi(\kb^-) > \phi(\ib') - \phi(\kb^-) = \up$. Therefore,
	\begin{equation}\label{eqn:g' y2 gap}
		|\phi(y_2) - \phi(\tlb')| > \frac{\up - \um}{2}.
	\end{equation}
	Consider a sender strategy $s'$ defined as
	\[
	s'(t) = \begin{cases}
		\bbb & \text{for } t = \tlb' \\
		s_r(t) & \text{otherwise.}
	\end{cases}
	\]
	Then
	\[
	\mathbb{U}(\tlb'; \rt, s') = U(\rt(s'(\tlb')), \tlb') - |\phi(s'(\tlb')) - \phi(\tlb')| = \um - \phi(\tlb') + \phi(\bbb).
	\]
	Since $s_{\rt}$ is the worst-case best response, at the feature vector $\tlb'$, we have
	\[
	\mathbb{U}(\tlb'; \rt, s') \leq \mathbb{U}(\tlb'; \rt, s_{\rt}),
	\]
	which, using \eqref{eqn:g' y2 gap}, gives
	\[
	\um - \phi(\tlb') + \phi(\bbb) \leq  - |\phi(y_2) - \phi(\tlb')| < \frac{\um - \up}{2}.
	\]
	This implies
	\[
	\phi(\tlb') - \phi(\bbb) > \frac{\up + \um}{2}.
	\]
	Since $\bbb \in [\ib', \tlb']$, and $\phi$ is increasing, we have $\phi(\tlb') - \phi(\bbb) \leq \phi(\tlb') - \phi(\ib')$. Hence
	\[
	\frac{\up + \um}{2} <  \phi(\tlb') - \phi(\ib').
	\]
	But this contradicts \eqref{eqn:gap between g and h}. Hence,
	\begin{equation}\label{eqn:rt neq minus}
		\rt(x) \neq \minus \quad \text{for all } x \in [\ib', \tlb'].
	\end{equation}
	From \eqref{eqn:ordering g'}, we have $[\ib', \tlb'] \subset [\kb^+, \kb^-]$. Hence, combining \eqref{eqn:rtx neq plus} with \eqref{eqn:rt neq minus}, it follows that $\rt(x) \neq \plus, \minus$ for $x \in [\ib', \tlb']$. Therefore, $\rt(x) = \Delta$ for $x \in [\ib', \tlb']$. As a consequence, we have that
	\begin{equation}\label{eqn:y1,y2 not delta}
		y_1, y_2 \notin [\ib', \tlb']
	\end{equation}
	because it would result in a $-\infty$ payoff for the sender.
\end{claims}


\begin{claims}{$y_1 < \ib'$.}\\
	We will prove this claim by showing that $y_1 \not> \tlb'$. Combining this with \eqref{eqn:y1,y2 not delta}, we will get $y_1 < \ib'$.\\
	Assume that $y_1 > \tlb'$. If $y_2 > \tlb'$, then from \eqref{eqn:s hat bad}
we get
	\[ \quad \phi(y_1) < \phi(y_2) - \up	\text{and}
	 \phi(y_2) < \phi(y_1) - \um
	\]
	respectively.
	Combining these inequalities, we get
	\[
	\phi(y_2) < \phi(y_2) - \up - \um.
	\]
	Since $\up \geq |\um|$, this results in a contradiction. Therefore, either $y_1 \not> \tlb'$ or $y_2 \not> \tlb'$. If $y_1 \not> \tlb'$, then we have proved our claim. Hence, suppose $y_2 \not> \tlb'$ and $y_1 > \tlb'$. Since $y_2 \not\in [\ib', \tlb']$ from \eqref{eqn:y1,y2 not delta}, therefore suppose $y_2 < \ib'$.
	From \eqref{eqn:s hat bad}, we get
	\[
	\phi(\ib') > \frac{\phi(y_1) + \phi(y_2) + \up}{2}
	\text{and}
 \phi(\tlb') < \frac{\phi(y_1) + \phi(y_2) - \um}{2}.
	\]
	respectively.
	Since $\phi(\tlb') > \phi(\ib')$, combining these inequalities yields
	\[
	\frac{\phi(y_1) + \phi(y_2) - \um}{2} > \frac{\phi(y_1) + \phi(y_2) + \up}{2},
	\]
	which implies
	\[
	-\um > \up.
	\]
	But this contradicts  $\up \geq |\um|$.
	Therefore, either $y_1 \not> \tlb'$ or $y_2 \not< \ib'$. But $y_2 \in \cg{X}$, and previously we have supposed $y_2 \not> \tlb'$ and $y_2 \in [\ib',\tlb']$, therefore $y_1 \not> \tlb'$.


\end{claims}

\begin{claims}{$y_2 > \tlb'$.}\\
	We will prove this claim by showing that $y_1 \not< \ib'$. Combining this with \eqref{eqn:y1,y2 not delta}, we will get $y_1 > \tlb'$.\\
	Assume that $y_2 < \ib'$.
	If $y_1 > \tlb'$, then from \eqref{eqn:s hat bad}, we get
	\[\phi(y_2) > \um + \phi(y_1) \text{and}
 \phi(y_1) > \up + \phi(y_2),
	\]
	respectively.
	Combining these inequalities, we get
	\[
	\phi(y_1) > \phi(y_1) + \up + \um.
	\]
	Since $\up \geq |\um|$, this results in a contradiction.
	Therefore, either $y_1 \not> \tlb'$ or $y_2 \not< \ib'$. If $y_2 \not< \ib'$, then we have proved our claim. Hence, suppose $y_1 \not> \tlb'$ and $y_2 < \ib'$. Since $y_1 \not\in [\ib',\tlb']$ from \eqref{eqn:y1,y2 not delta}, therefore suppose $y_2 < \ib'$. Then \eqref{eqn:s hat bad} reduces to
	\[ \phi(\tlb') < \frac{\phi(y_1) + \phi(y_2) - \um}{2} \text{and}
 \phi(\ib') > \frac{\phi(y_1) + \phi(y_2) + \up}{2},
	\]
	respectively.
	Since $\phi(\tlb') > \phi(\ib')$, combining these inequalities gives
	\[
	\frac{\phi(y_1) + \phi(y_2) - \um}{2} > \frac{\phi(y_1) + \phi(y_2) + \up}{2},
	\]
	which implies
	\[
	-\um > \up.
	\]
	But this contradicts $\up \geq |\um|$. Therefore, either $y_1 \not< \ib'$ or $y_2 \not< \ib'$. But $y_1 \in \cg{X}$, and previously we have supposed $y_1 \not> \tlb'$ and $y_1 \not\in [\ib',\tlb']$, therefore $y_2 \not< \ib'$.

\end{claims}


\begin{claims}{$\phi(\tlb') - \phi(\ib') > \frac{\up + \um}{2}$.}\\
	From Claims 1-3, we have
	\begin{equation}\label{eqn:ordering y1 y2}
		y_1 < \ib' < \h < \tlb' < y_2.
	\end{equation}
	Using \eqref{eqn:ordering y1 y2} in \eqref{eqn:s hat bad}, we get
	\[\phi(\ib') < \frac{\phi(y_1) + \phi(y_2) - \up}{2} \text{and}
 \phi(\tlb') > \frac{\phi(y_1) + \phi(y_2) + \um}{2}.
	\]
	Thus,
	\[
	\phi(\tlb') - \phi(\ib') > \frac{\up + \um}{2}
	\] respectively.
	This result contradicts the assumption made in Case 4. Hence, Case 4 cannot occur.
\end{claims}

We now complete the proof of Theorem \ref{thm:Deterministic Stackelberg equilibrium strategy}.
\begin{apdxproof}(Proof of Theorem \ref{thm:Deterministic Stackelberg equilibrium strategy}).\\
	From Lemma \ref{lem:Deterministic domination of a family of strategies}, every receiver strategy $\rt$ is dominated by a receiver strategy $r(\cdot) \in \sr{R}$. Hence, the Stackelberg equilibrium is a strategy in $\sr{R}$ with the least error. Recall that $\sr{R} = \sr{R}_1 \cup \sr{R}_2$. Let's further partition the subset $\sr{R}_1$ into two parts, $\sr{R}_{10}$ and $\sr{R}_{11}$, i.e., $\sr{R}_1 := \sr{R}_{10} \cup \sr{R}_{11}$, where
	\[
	\begin{aligned}
		\sr{R}_{10} &:= \{r(\cdot;\ib,\jb) \ |\ r(\cdot;\ib,\jb) \in \sr{R}_1, \phi(\h) - \phi(\jb) \leq \um + \phi(\ib) - \phi(\h)\}, \\
		\sr{R}_{11} &:= \{r(\cdot;\ib,\jb) \ |\ r(\cdot;\ib,\jb) \in \sr{R}_1, \phi(\h) - \phi(\jb) > \um + \phi(\ib) - \phi(\h)\}.
	\end{aligned}
	\]
Thus, $\sr{R} = \sr{R}_{10} \cup \sr{R}_{11} \cup \sr{R}_2$. We show that there exists a Stackelberg equilibrium in $\sr{R}$ by showing that each subset, $\sr{R}_{10}$, $\sr{R}_{11}$, and $\sr{R}_2$, contains a strategy that yields the least error within its respective subset. Thus, the Stackelberg equilibrium will be a strategy with the least error among the three strategies that have the least error in their subsets.\\
	For $r(\cdot;\ib,\jb) \in \sr{R}_{10}$, using Lemma \ref{lem:Best response r1}(a), the error set is $E(r(\cdot;\ib,\h)) = [\ib,\h]$. Here, $\up \geq \phi(\h) - \phi(\ib) \geq \frac{\up + \um}{2}$. Now consider the strategy $r(\cdot;\ib^*,\h)$, where $\ib^*$ is such that $\phi(\h) - \phi(\ib^*) = \frac{\um + \up}{2}$. Such a feature vector $\ib^*$ exists due to the continuity of $\phi$. The error set of $r(\cdot;\ib^*,\h)$ is $E(r(\cdot;\ib^*,\h)) = [\ib^*,\h]$. Since $\phi(\h) - \phi(\ib) \geq \phi(\h) - \phi(\ib^*)$, and $\phi$ is increasing, we have $\ib \leq \ib^*$, and hence, $[\ib^*,\h] \subset [\ib,\h]$. Since $r(\cdot;\ib,\h)$ was an arbitrary strategy, $r(\cdot;\ib^*,\h)$ is the strategy with the least error in $\sr{R}_{10}$.\\
	For $r(\cdot;\ib,\jb) \in \sr{R}_{11}$, using Lemma \ref{lem:Best response r1}(b), it follows that the error set of the strategy is $E(r(\cdot;\ib,\jb)) = [\ib,\tlb]$. Since we seek a strategy with the least error, we choose $\ib^*$ such that the following holds:
	\[
	\ib^* := \arg\inf_{\ib \in [\phi^{-1}(\phi(\h) - \up), \h]} D(\{x : x \in [\ib, \lb]\}).
	\]
	We now show that the infimum of the above optimization problem does indeed exist. Recall that $\lb$ is completely determined by $\ib$ and $\up$. Since the probability law $D$ has a continuous  (in fact differentiable) density function $\dot{\psi}$ except at finitely many points, the map
	\[
	\ib \mapsto D(\{x : x \in [\ib, \tlb]\}) = \int_{\ib}^{\tlb} \frac{d\psi(x)}{dx}\, dx
	\]
	is (finite) piecewise continuous for $\ib \in [\phi^{-1}(\phi(\h) - \up), \h]$. Because we are optimizing a (finite) piecewise continuous function over a compact set, the infimum is indeed attained.
	For $\sr{R}_2$, which contains only two elements, the majority class (w.r.t. $D$) strategy is the one with the least error.
\end{apdxproof}


\noindent Based on all the cases, we can see that there always exists a receiver strategy $r \in \sr{R}$ which dominates $\rt$. Therefore, the family of receiver strategies $\sr{R}$ dominates all other receiver strategies.
\end{proof}

\begin{apdxproof}(Proof of Theorem \ref{thm:u_ < - u_+})\label{proof:u_ < - u_+}\\

We will show that $E(r) = \emptyset$, which will prove that $r(\cdot;\h,\jb)$ is a Stackelberg equilibrium. Note that such a $\jb$ exists due to the continuity of $\phi$. We now calculate $s_r$. The calculation is similar to Lemma \ref{lem:Best response r1}.

	\begin{case}{$x \in [0, \h)$.}\\
		The sender's payoff is given by $U(r(s_r(x)),x) - |\phi(s_r(x)) - \phi(x)|$. Based on the structure of $r$, $s_r(x)$ can either be $x$ or $\jb$. Therefore, the sender's payoff is
		\[
		\mathbb{U}(x; r, s_r) = \begin{cases}
			0 & \text{if } s_r(x) = x,\\
		\up - \phi(\jb) + \phi(x) & \text{if } s_r(x) = \jb.
		\end{cases}
		\]
	We have $- \up = \phi(\h) - \phi(\jb) < \phi(x) - \phi(\jb)$. Thus $\up - \phi(x) - \phi(\jb) <0$. Hence, $s_r(x) = x$.
	\end{case}

	\begin{case}{$x \in (\jb, 1]$.}\\
		The sender's payoff is given by $U(r(s_r(x)),x) - |\phi(s_r(x)) - \phi(x)|$. Based on the structure of $r$, $s_r(x)$ can either be $x$ or $\h$. The sender's payoff is
		\[
		\mathbb{U}(x; r, s_r) = \begin{cases}
			0 & \text{if } s_r(x) = x,\\
		- \phi(\h) + \phi(x) & \text{if } s_r(x) = \h.
		\end{cases}
		\]
		Since $x \geq \h$, we have $\phi(\h) - \phi(x) < 0$. Hence, $s_r(x) = x$.
	\end{case}

	\begin{case}{$x \in [\h, \jb]$.}\\
		The sender's payoff is given by $U(r(s_r(x)),x) - |\phi(s_r(x)) - \phi(x)|$. Based on the structure of $r$, $s_r(x)$ can either be $\h$ or $\jb$. Therefore, the sender's payoff is
		\[
		\mathbb{U}(x; r, s_r) = \begin{cases}
			\um - \phi(x) + \phi(\h) & \text{if } s_r(x) = \h,\\
		 - \phi(\jb) + \phi(x) & \text{if } s_r(x) = \jb.
		\end{cases}
		\]
		Since $\jb \geq x \geq \h$, we have $\phi(x) - \phi(\jb) \geq \phi(\h) - \phi(\jb) = -\up$, $\um \leq -\up \leq 0$, and $\phi(\h) - \phi(x) \leq 0$. Hence, $\um - \phi(x) + \phi(\h) \leq 0 - \phi(\jb) + \phi(x)$ for all $x \in [\h, \jb]$, and therefore we have $s_r(x) = \jb$.
	\end{case}

	\noindent	Combining all the cases, we have
	\[
	s_r(x) := \begin{cases}
		x & \text{for } x \in [0, \h) \cup (\jb, 1],\\
		\jb & \text{for } x \in [\h, \jb].
	\end{cases}
	\]
	\noindent Thus, $E(r) = \{x \mid r(s_r(x)) = h(x)\} = \emptyset$.
\end{apdxproof}




        \subsection{Omitted Proofs of Section III.C}

        \input{Proofs Section 3C}

        \subsection{Omitted Proofs of Section III.D}

        The next Corollary show that if a Stackelberg equilibrium exist under worst-case best response setup, then its error will be $\frac{\uu}{2}$.
\begin{cor}\label{cor:Stackelberg worst case best response error}
    Suppose there exists a Stackelberg equilibrium $(r^\star, s_{r^\star})$ under the setup where the sender plays the worst-case best response. Then, $\cg{E}(r^\star,s_{r^\star}) = \frac{\uu}{2}$.
\end{cor}


\begin{proof}\label{proof:Stackelberg worst case best response error}
        Let $s'$ be the worst-case lazy best response to $r^\star$. Let $(\bar{r}^\star, s_{\bar{r}^\star})$ be the Stackelberg equilibrium under the worst-case lazy best response.
    From Proposition \ref{prop:comparing sender strategy types}, we know that
    \[
    \cg{E}(r^\star, s_{r^\star}) \geq \cg{E}(r^\star, s').
    \]
    Additionally, by the definition of $\bar{r}^\star$, we have
    \[
    \cg{E}(r^\star, s') \geq \cg{E}(\bar{r}^\star, s_{\bar{r}^\star}).
    \]
    From Theorem \ref{thm:Stochastic Stackelberg equilibrium}, we know that
    \[
    \cg{E}(\bar{r}^\star, s_{\bar{r}^\star}) = \frac{\uu}{2}.
    \]
    Combining these results, we conclude that
    \[
    \cg{E}(r^\star, s_{r^\star}) \geq \frac{\uu}{2}.
    \]
    Furthermore, from Example \ref{ex:eps stochastic strategy} (which appears in Section \ref{subsec:Examples and comparision}), we know that for every $\eps > 0$,
    \[
    \cg{E}(r_\eps,s_{r_\eps}) = \frac{\uu + \eps}{2}.
    \]
    Thus, if a Stackelberg equilibrium $r^\star$ exists, then it follows that
    \[
    \cg{E}(r^\star,s_{r^\star}) = \frac{\uu}{2}.
    \]
\end{proof}

From the definition of $\msf{R}$ in \eqref{defn:stochastic_receiver}, every receiver strategy $r \in \msf{R}$ has the structure $r(x) \equiv (p(x),q(x),1-p(x)-q(x))$. The first lemma shows that without loss of generality, we can assume the structure of $r$ at each $x$ to be either $r(x) \equiv (1-q(x),q(x),0)$ or $r(x) \equiv (0,0,1)$.


\begin{lem}\label{lem:delta null wlog}
    Let $r \equiv (p,q,1-p-q) \in \msf{R}$ be a receiver strategy. Let $\Theta := \{x': x' \in \cg{X}, p(x') + q(x') <1\}$. Now consider another receiver strategy $\rt \equiv (\pt,\qt,1-\pt - \qt)$ defined as
    \[
    \rt(x) = \begin{cases}
        r(x) &\text{for } x \in \cg{X}\setminus \Theta,\\
        (0,0,1) &\text{for } x \in \Theta.
    \end{cases}
    \]
    Then $\cg{E}(r) = \cg{E}(\rt)$.
\end{lem}

\begin{proof}\label{proof:delta null wlog}
Let $x \in \cg{X}$. We first show that $s_r(x),s_{\rt}(x) \not\in \Theta$. From the definition of $p,q$ and $\Theta$, we have
    \[1-p(x')-q(x') =\prob_r(\Delta|X' = x')  > 0, \text{for all} x' \in \Theta.\]
    \noindent If $s_r(x) \in \Theta$, then $$\prob_r(y = \Delta|X'_{s_r}=s_r(x)) > 0$$ and the sender's payoff is $\mathbb{U}(x;r,s_r) = -\infty$.
    Similarly, if $s_{\rt}(x) \in \Theta$, then \[\prob_{\rt}(y = \Delta|X'_{s_{\rt}} = s_{\rt}(x)) = 1 \] and the sender's payoff is $\mathbb{U}(x;\rt,s_{\rt}) = -\infty$.
    But from \eqref{eqn:receiver condition 1}, the feature vector set $\cg{X}\setminus \Theta \neq \emptyset$. Therefore, there exists $t \in \cg{X}\setminus \Theta$ such that
        $$\prob_r(Y = \Delta|X'=t) = \prob_{\rt}(Y = \Delta|X'=t) = 0.$$
Hence, for the sender strategy $s(x) = t$ for all $x \in \cg{X}$, we have
    $$\prob_r(Y=\Delta|X'=t) = \prob_{\rt}(Y=\Delta|X'=t)) = 0$$ and
    the sender's payoff at $x$ for receiver strategies $r$ and $\rt$ will be $\mathbb{U}(x,r,s) < \infty\text{and} \mathbb{U}(x,\rt,s) < -\infty$ respectively.
    But this contradicts that $s_r$ and $s_{\rt}$ are the worst-case lazy best responses, as $\mathbb{U}(x,r,s) < \mathbb{U}(x,r,s_r)$ and $\mathbb{U}(x,\rt,s) < \mathbb{U}(x,\rt,s_{\rt})$. Hence,
    \begin{equation}\label{lem:eqn:non-Delta best response}
        s_r(x), s_{\rt}(x) \in \cg{X} \setminus \Theta \text{ for all } x \in \cg{X}.
    \end{equation}
    From the construction of $\rt$, we know that $r(x) = \rt(x)$ for $x \in \cg{X} \setminus \Theta$. Therefore, combining it with \eqref{lem:eqn:non-Delta best response}, we get
    \begin{equation}\label{eqn:lem:non-Delta best response 2}
        r(s_r(x)) = \rt(s_r(x)) \text{ and } r(s_{\rt}(x)) = \rt(s_{\rt}(x)) \text{ for all } x \in \cg{X}.
    \end{equation}
    We next show that $s_r$ and $s_{\rt}$ are best responses to both $r$ and $\rt$. Since $s_r$ and $s_{\rt}$ are worst-case lazy best responses of $r$ and $\rt$, respectively, we have
    \begin{equation}\label{eqn:lem:non-Delta best response 3}
        \mathbb{U}(x;r,s_r) \geq \mathbb{U}(x;r,s_{\rt}) \text{ and } \mathbb{U}(x;\rt,s_{\rt}) \geq \mathbb{U}(x;\rt,s_r) \text{ for all } x \in \cg{X}.
    \end{equation}
    Substituting \eqref{eqn:lem:non-Delta best response 2} in \eqref{eqn:lem:non-Delta best response 3} and expanding it, we get



    \begin{equation*}
        \begin{aligned}
            \fk{U}(x;r,s_r) - |s_r(x)-x| &\geq \fk{U}(x;r,s_{\rt})   - |s_{\rt}(x) -x| = \fk{U}(x;\rt,s_{\rt})   - |s_{\rt}(x) -x| &\text{ and }\\
            \fk{U}(x;\rt,s_{\rt}) - |s_{\rt}(x)-x| &\geq \fk{U}(x;\rt,s_r)   - |s_r(x)-x| = \fk{U}(x;r,s_r)   - |s_r(x)-x|.
        \end{aligned}
    \end{equation*}
    which implies
    \[
    |s_{\rt}(x) - x| - |s_r(x) -x| \geq \fk{U}(\rt(s_{\rt}(x)),x) - \fk{U}(r(s_r(x)),x) \geq |s_{\rt}(x) - x| - |s_r(x) -x|
    \]
    or
    \begin{equation}\label{eqn:lem:non-Delta best response 4}
        \fk{U}(\rt(s_{\rt}(x)),x) - \fk{U}(r(s_r(x)),x) = |s_{\rt}(x) - x| - |s_r(x) -x|.
    \end{equation}
    Thus, substituting \eqref{eqn:lem:non-Delta best response 4} back in \eqref{eqn:lem:non-Delta best response 3}, we get
    \begin{equation}\label{eqn:lem:non-Delta best response 5.1}
        \mathbb{U}(x;r,s_r) = \mathbb{U}(x;r,s_{\rt}) \text{ and } \mathbb{U}(x;\rt,s_{\rt}) = \mathbb{U}(x;\rt,s_r).
    \end{equation}
    This shows that $s_{\rt}$ and $s_r$ are best responses to both $r$ and $\rt$.
    Furthermore, since $s_r$ and $s_{\rt}$ are both worst-case lazy best responses to $r$ and $\rt$, respectively, using the laziness constraint on \eqref{eqn:lem:non-Delta best response 5.1}, we have
    \begin{equation}\label{eqn:lem:non-delta best response 6}
        |s_r(x)-x| \leq |s_{\rt}(x) - x| \text{ and } |s_{\rt}(x) - x| \leq |s_r(x) - x| \text{ which implies } |s_r(x)-x|= |s_{\rt}(x) -x|.
    \end{equation}
    As an additional consequence, \eqref{eqn:lem:non-Delta best response 4} reduces to
    \begin{equation}\label{eqn:lem:non-delta best response 7}
        \fk{U}(r(s_r(x)),x) = \fk{U}(\rt(s_{\rt}(x)),x).
    \end{equation}
    Next, we show that $s_r$ and $s_{\rt}$ are worst-case lazy best responses to both $r$ and $\rt$. Since $s_r$ and $s_{\rt}$ are worst-case lazy best responses of $r$ and $\rt$, respectively, we have (using the worst-case constraint)
    \begin{equation}\label{eqn:lem:non-Delta best response 5}
        \begin{aligned}
            &\cg{E}(r,s_r) = \int_\cg{X} T(x;r,s_r) \, dx \geq \int_\cg{X} T(x;r,s_{\rt}) \, dx = \cg{E}(r,s_{\rt})\quad \text{and} \\
            &\cg{E}(\rt,s_{\rt}) = \int_\cg{X} T(x;\rt,s_{\rt}) \, dx \geq \int_\cg{X} T(x;\rt,s_r) \, dx = \cg{E}(\rt,s_r).
        \end{aligned}
    \end{equation}
    Recall from \eqref{eqn:error function} that $T$ is defined as $T(x;r,s) = \prob_r(Y \neq h(x) \mid X'=s(x))$. Furthermore, if $s(x) \in \cg{X} \setminus \Theta$, then $$\fk{U}(x;r,s) = 2\uu\prob_r(Y \neq h(x) \mid X'=s(x))  = 2\uu T(x;r,s).$$ Hence, using this in \eqref{eqn:lem:non-Delta best response 5}, we get
    \[
    \int_\cg{X} \fk{U}(r(s_r(x)),x) \, dx \geq \int_\cg{X} \fk{U}(r(s_{\rt}(x)),x) \, dx \quad \text{and} \quad \int_\cg{X} \fk{U}(\rt(s_{\rt}(x)),x) \, dx \geq \int_\cg{X} \fk{U}(\rt(s_r(x)),x) \, dx.
    \]
    But combining \eqref{eqn:lem:non-Delta best response 2} with \eqref{eqn:lem:non-delta best response 7}, we get
    \[
    \fk{U}(\rt(s_r(x)),x) = \fk{U}(r(s_r(x)),x) = \fk{U}(\rt(s_{\rt}(x)),x) = \fk{U}(r(s_{\rt}(x)),x).
    \]
    Thus,
    \begin{equation*}\label{eqn:lem:non-Delta best response 8}
        \int_\cg{X} T(x;r,s_r) \, dx = \int_\cg{X} T(x;r,s_{\rt}) \, dx \quad \text{and} \quad \int_\cg{X} T(x;\rt,s_{\rt}) \, dx = \int_\cg{X} T(x;\rt,s_r) \, dx.
    \end{equation*}
    Hence, $s_{\rt}$ and $s_r$ are also worst-case lazy best responses of $r$ and $\rt$, respectively. Therefore, $\cg{E}(r) = \cg{E}(\rt)$ as claimed.
\end{proof}

The next lemma calculates the worst-case lazy best response and classification error of a specific receiver strategy $ r $, which will later be shown to be a Stackelberg equilibrium.

\begin{lem}\label{lem:error is half}
    Let $ r \equiv (1-q,q,0) $, where
    \begin{equation}\label{eqn:q defn}
        q(x) = \begin{cases}
            0 & \text{for } x \in [0,\ib], \\
            \frac{x-\ib}{2\uu} & \text{for } x \in (\ib,\jb), \\
            1 & \text{for } x \in [\jb,1]
        \end{cases}
    \end{equation}
    and, $\ib = \h - \uu$ and $\jb = \h + \uu$.
    Then the worst-case lazy best response of $ r $ is $ s_r(x) = x $, and the error function $ T(x;r,s_r) $ is
    \[
    T(x;r,s_r) = \begin{cases}
        0 & \text{for } x \in [0,\ib] \cup [\jb,1], \\
        q(x) & \text{for } x \in (\ib,\h], \\
        1-q(x) & \text{for } x \in (\h,\jb).
    \end{cases}
    \]
    Furthermore, $ \cg{E}(r) = \frac{\uu}{2} $.
\end{lem}

\begin{proof}\label{proof:error is half}
    We divide the task of finding $ s_r $ over $ \cg{X} $ into various subintervals of $ \cg{X} $ and show that $s_r(x) = x$ in all of them.  Let $s_r(x) = x'$.

    \begin{case} $ x \in [0,\ib] $.\\
 We now calculate $\mathbb{U}(x;r,s_r)$ and show that in the current case, $x=x'$.

\begin{equation*}
    \mathbb{U}(x;r,s_r) = \begin{cases}
    -|x-x'| = 0 &\text{if } x'= x,\\
    -|x-x'| < 0 &\text{if } x'\neq x \text{and} x' \in [0,\ib],\\
    2\uu q(x') - x' + x \leq 0 &\text{if } x' \in (\ib,\jb),\\
    2\uu - x' + x \leq 0&\text{if } x' \in [\jb,1].
    \end{cases}
\end{equation*}
Here we have used \eqref{eqn:q defn} to expand $q(x')$ and to show that  $2\uu q(x') - x' + x \leq 0$, and have used $\jb - \ib = 2\uu$ to show that $2\uu -x' + x \leq 0$.
Thus for $x' \neq x$, we have $\mathbb{U}(x;r,s_r) \leq 0$ whereas for $x' = x$, $\mathbb{U}(x;r,s_r) = 0$. Furthermore, $x'=x$ is also the laziest option for $s_r$. Hence $s_r(x) = x$.
\end{case}

\begin{case}$ x \in (\ib,\h] $.\\
 We now calculate $\mathbb{U}(x;r,s_r)$ and show that in the current case, $x=x'$.
 \begin{equation*}
    \mathbb{U}(x;r,s_r) = \begin{cases}
    2\uu q(x') - x + x' = x - \ib &\text{if } x'= x,\\
    -x + x' < 0. &\text{if } x' \in [0,\ib],\\
    2\uu q(x') - x + x' \leq x- \ib &\text{if } x' \in (\ib,x] \text{and} x' \neq x,\\
    2\uu q(x') - x' + x < x-\ib    &\text{if} x' \in (x,\jb),\\
    2\uu - x' + x \leq x - \ib &\text{if } x' \in [\jb,1].
    \end{cases}
\end{equation*}
Here we have used \eqref{eqn:q defn} to expand $q(x')$ wherever it appears and have used $2\uu = \jb - \ib$ to show that $2\uu -x' + x \leq x - \ib$. Note that $x - \ib \geq 0$.
Thus for $x' \neq x$, we have $\mathbb{U}(x;r,s_r) \leq x-\ib$ whereas for $x' = x$, $\mathbb{U}(x;r,s_r) \leq x - \ib$. Furthermore, $x'=x$ is also the laziest option for $s_r$. Hence $s_r(x) = x$.
\end{case}

   \begin{case} $ x \in [\h, \jb) $.\\
 We now calculate $\mathbb{U}(x;r,s_r)$ and show that in the current case, $x=x'$.
 \begin{equation*}
    \mathbb{U}(x;r,s_r) = \begin{cases}
    2\uu(1-q(x')) - x' + x = \jb -x &\text{if } x'= x,\\
    2\uu - x + x' \leq \jb - x &\text{if } x' \in [0,\ib],\\
    2\uu(1-q(x')) - x + x' < \jb- x &\text{if } x' \in (\ib,x],\\
    2\uu(1- q(x')) - x' + x \leq \jb - x    &\text{if} x' \in (x,\jb) \text{and} x' \neq x,\\
    -x' + x \leq  -\jb + x &\text{if } x' \in [\jb,1].        \end{cases}
\end{equation*}
Here we have used \eqref{eqn:q defn} to expand $q(x')$ wherever it appears and have used $2\uu = \jb - \ib$ to show that $2\uu - x + x' \leq \jb -x$. Note that $\jb - x \geq 0$. Thus for $x' \neq x$, we have $\mathbb{U}(x;r,s_r) \leq \jb-x$ whereas for $x' = x$, $\mathbb{U}(x;r,s_r) = \jb-x$. Furthermore, $x'=x$ is also the laziest option for $s_r$. Hence $s_r(x) = x$.
   \end{case}

\begin{case} $ x \in [\jb, 1] $.\\
We now calculate $\mathbb{U}(x;r,s_r)$ and show that in the current case, $x=x'$.

 \begin{equation*}
    \mathbb{U}(x;r,s_r) = \begin{cases}
    -|x'-x| = 0 &\text{if } x'= x,\\
    2\uu - x + x' \leq \jb - x \leq 0 &\text{if } x' \in [0,\ib],\\
    2\uu(1-q(x')) - x + x' < \jb- x < 0 &\text{if } x' \in (\ib,\jb),\\
    -|x'-x| <  0 &\text{if } x' \in [\jb,1] \text{and} x' \neq x.       \end{cases}
\end{equation*}
Here we have used \eqref{eqn:q defn} to expand $q(x')$ to show that $2\uu(1-q(x')) - x + x' < \jb- x < 0$. Thus for $x' \neq x$, we have $\mathbb{U}(x;r,s_r) \leq 0$ whereas for $x' = x$, $\mathbb{U}(x;r,s_r) = 0$. Furthermore, $x'=x$ is also the laziest option for $s_r$. Hence $s_r(x) = x$.
\end{case}
\noindent Hence, by combining all the cases, we conclude that $ s_r(x) = x $. Therefore,
    \[
    T(x; r, s_r) = \prob_{r}\{Y \neq h(x) \mid X' = s_r(x)\} = \begin{cases}
        0 & \text{for } x \in [0, \ib] \cup [\jb, 1], \\
        q(x) & \text{for } x \in (\ib, \h], \\
        1 - q(x) & \text{for } x \in (\h, \jb).
    \end{cases}
    \]
    Since $ \cg{E}(r) = \int_\cg{X} T(x) \, dx $, it follows that $ \cg{E}(r) = \frac{\uu}{2} $.
\end{proof}

We next show that if a feature vector $x'$ is the worst-case lazy best response of some other feature vector $x$ and both have the same true label, then $j$ is also the worst-case lazy best response of itself.

\begin{lem}\label{lem:self bestresponse}
    Let $r \in \msf{R}$ be a receiver strategy, and let $x,x' \in \cg{X}$ be such that $s_r(x) = x'$. Then, for all $\bar{x} \in [\min(x,x'),\max(x,x')] \setminus \{x\}$ such that $h(\bar{x}) = h(x)$, we have $s_r(\bar{x}) = x'$.
\end{lem}

\begin{proof}\label{proof:self bestresponse}
     Let $\bar{x} \in [\min(x,x'),\max(x,x')] \setminus \{x\}$ such that $h(\bar{x}) = h(x)$. Let $s_r(\bar{x}) = z$. We will assume $x' < x$. A similar argument will hold if $x' \geq x$. Consider the sender strategies $s(t) \equiv x'$ and $\hat{s}(t) \equiv z$ for all $t \in \cg{X}$. Since $s_r$ is the worst-case lazy best response of $r$, we have
   $$\mathbb{U}(t,r;s_r) \geq \mathbb{U}(t,r;s) \text{ and } \mathbb{U}(r;s_r) \geq \mathbb{U}(r;\hat{s}) \text{for all} t \in \cg{X}.$$
    This implies for $t= \bar{x}$ and $t=x$, we have
    $$
    \mathbb{U}(\bar{x};r,s_r) \geq \mathbb{U}(\bar{x};r,s) \quad \text{and} \quad \mathbb{U}(x;r,s_r) \geq \mathbb{U}(x,r,\hat{s}) \text{ respectively}.
    $$
    Recall that from the definition of $U$ given in \eqref{eqn:Utility}, $U(k,\bar{x}) = U(k,x)$ for all $k \in \cg{Y} \cup \Delta$ because $h(\bar{x}) = h(x)$. As a consequence, for a fixed receiver strategy $\fk{r}$ and sender strategy $\fk{s}$, we have $\fk{U}(x;\fk{r},\fk{s}) = \fk{U}(\bar{x};\fk{r},\fk{s})$. Hence, this implies
    \begin{equation}\label{lem:eqn:w in y x}
    \begin{aligned}
        \fk{U}(x;r,s_r) - |\bar{x}-z| = \fk{U}(\bar{x};r,s_r) - |\bar{x}-z| &\geq \fk{U}(\bar{x};r,s) - \bar{x} + x' = \fk{U}(x;r,s) - \bar{x} + x' \quad \text{and} \\
        \fk{U}(x;r(x'),s) - x + x' &\geq \fk{U}(x;r,s_r) - |x-z|.
    \end{aligned}
    \end{equation}
    We show that $z = x'$ by the method of contradiction. Suppose $z \neq x'$. Then, consider the following cases.
\textbf{Case 1. }{$z \geq x$.} \\
        We show that this case does not happen. From the case assumption, \eqref{lem:eqn:w in y x} reduces to
        \begin{equation}
            \fk{U}(x;r,s_r) \geq \fk{U}(x;r,s) - 2\bar{x} + x' + z \quad \text{and} \quad \fk{U}(x;r,s) \geq \fk{U}(x;r,s_r) + 2x - x' - z.
        \end{equation}
        Combining them gives
        $
        \fk{U}(x;r,s_r) \geq \fk{U}(x;r,s_r) + 2(x-\bar{x}),
        $
        which gives rise to a contradiction because $x - w > 0$. Hence, this case does not happen as claimed.\\
\textbf{Case 2. }{$x \geq z > \bar{x} \geq x'$.} \\
        We show that this case does not happen. From the case assumption, \eqref{lem:eqn:w in y x} reduces to
        \begin{equation}
            \fk{U}(x,r,s_r) \geq \fk{U}(x;r,s) - 2\bar{x} + z + x' \quad \text{and} \quad \fk{U}(x;r,s) > \fk{U}(x;r,s_r) + z - x'.
        \end{equation}
        The second part has a strict inequality because $|x-z| < |x-x'|$. The reason is that if there were equality, then $z$ would have been a low-cost (lazy) option for the sender at $x$ as compared to $x'$, which would contradict that $x'$ is the worst-case lazy best response of $r$ at $x$.
        Combining them gives
        $
        \fk{U}(r(z),x) > \fk{U}(r(z),x) + 2(z-\bar{x}),
        $
        which is a contradiction because $z > \bar{x}$. Hence, this case does not happen as claimed.\\
\textbf{Case 3. }{$x > \bar{x} \geq z > x'$.} \\
        We show that this case does not happen. From the case assumption, \eqref{lem:eqn:w in y x} reduces to
        \begin{equation}
            \fk{U}(x;r,s_r) \geq \fk{U}(x;r,s) + x' - z \quad \text{and} \quad \fk{U}(x;r,s) > \fk{U}(x;r,s_r) + z - x'.
        \end{equation}
        The second part has a strict inequality because $|x-z| < |x-x'|$. The same reasons as mentioned in Case 2. Combining them gives
        $
        \fk{U}(x;r,s_r) > \fk{U}(x;r,s_r),
        $
        which is a contradiction. Hence, this case does not happen as claimed.\\
\textbf{Case 4. }{$z < x'$.} \\
        We show that this case does not happen. From the case assumption, \eqref{lem:eqn:w in y x} reduces to
        \begin{equation}
            \fk{U}(x;r,s_r) > \fk{U}(x;r,s) + x' - z \quad \text{and} \quad \fk{U}(x;r,s) \geq \fk{U}(x;r,s_r) + z - x'.
        \end{equation}
        The first part has a strict inequality because $|\bar{x}-x'| < |\bar{x}-z|$. The reason is the same as mentioned in Case 2. This further reduces to
        $
        \fk{U}(x;r,s_r) > \fk{U}(x;r,s_r),
        $
        which is a contradiction.
Hence, combining all the cases, we get a contradiction. Therefore, $z = x'$.
\end{proof}

Let $r \equiv (p,q,1-p-q) \in \msf{R}$ be the receiver strategy. For feature vectors $x_1$ and $x_2 \in \cg{X}$ such that $x_1 \leq \h < x_2$, we calculate the contribution of error from the intervals $[x_1, x_1 - 2\uu q(x_1)]$ and $[x_2, x_2 + 2\uu q(x_2)]$.

\begin{lem}\label{lem:conservative estimate of the error outside the delta interval}
    Let $r \equiv (p, q,1-p-q) \in \msf{R}$ be a receiver strategy. Let $x_1, x_2 \in \cg{X}$ be such that $x_1 \leq \h < x_2$, and $p(x_i) + q(x_i) = 1$ for $i = 1, 2$. Then
    \begin{equation*}
    \begin{aligned}
        \int_{x_1 - 2\uu q(x_1)}^{x_1} T(x;r,s_r) \, dx &\geq \uu q(x_1)^2, \quad \text{and}\;
        \int_{x_2}^{x_2 + 2\uu(1-q(x_2))} T(x;r,s_r) \, dx &\geq \uu (1-q(x_2))^2.
    \end{aligned}
    \end{equation*}
\end{lem}

\begin{proof}  Let $q_1 := q(x_1)$ and $q_2 := q(x_2)$. Recall that from \eqref{eqn:error function}, $T$ is defined as $T(x;r,s_r) = \prob(r(s_r(x)) \neq h(x) |X'=s_r(x)  = q(s_r(x))$. We upper bound $q(s_r(x))$ (or $1-q(x)$ depending upon x) to get a upper bound on $T(x;r,s_r)$ and eventually over $\int_{x_1 - 2\uu q(x_1)}^{x_1} T(x;r,s_r)dx$ and $\int_{x_2}^{x_2 + 2\uu(1-q(x_2))} T(x;r,s_r)dx$.
     Let $x \in [x_1 - 2\uu q_1, x_1]$. Consider a sender strategy $s(x) = x_1$. Then, from the definition of the worst-case lazy best response, we have
    \[
    \mathbb{U}(x;r,s_r) \geq \mathbb{U}(x;r,s)
    \]
    which implies
    \begin{equation}\label{eqn:expand T}
        2\uu q(s_r(x)) - |s_r(x) - x| \geq 2\uu q_1 - x_1 + x \quad \text{or} \quad 2\uu q(s_r(x)) \geq 2\uu q_1 - x_1 + x.
    \end{equation}
     Hence, from \eqref{eqn:expand T}, we get $T(x;r,s_r) \geq q_1 - \frac{x_1 - x}{2\uu}$. Therefore,
    \begin{equation}\label{lem:eqn:q_1>q_2 left side}
        \int_{x_1 - 2\uu q_1}^{x_1} T(x;r,s_r) \, dx \geq \int_{x_1 - 2\uu q_1}^{x_1} \left( q_1 - \frac{x_1 - x}{2\uu} \right) \, dx = \uu q_1^2.
    \end{equation}
    Similarly, let $x \in [x_2, x_2 + 2\uu(1-q_2)]$. Consider a sender strategy $\bar{s}(x) = x_2$. Then, from the definition of the worst-case lazy best response, we have
    \[
    \mathbb{U}(x;r,s_r) \geq \mathbb{U}(x;r,\bar{s})
    \]
    which implies
    \[
    2\uu (1-q(s_r(x))) - |s_r(x) - x| \geq 2\uu (1-q_2) - x + x_2 \quad \text{or} \quad 2\uu (1-q(s_r(x))) \geq 2\uu (1-q_2) - x + x_2.
    \]
    Hence, $T(x;r,s_r) = 1 - q(s_r(x)) \geq (1 - q_2) - \frac{x - x_2}{2\uu}$. Therefore,
    \begin{equation}\label{lem:eqn:q_1>q_2 right side}
        \int_{x_2}^{x_2 + 2\uu (1-q_2)} T(x;r,s_r) \, dx \geq \int_{x_2}^{x_2 + 2\uu (1-q_2)} \left( (1 - q_2) - \frac{x - x_2}{2\uu} \right) \, dx = \uu (1-q_2)^2.
    \end{equation}
Hence, combining \eqref{lem:eqn:q_1>q_2 left side} and \eqref{lem:eqn:q_1>q_2 right side}, we get
    \begin{equation*}
    \begin{aligned}
        \int_{x_1 - 2\uu q(x_1)}^{x_1} T(x;r,s_r) \, dx \geq \uu q(x_1)^2, \quad \text{and} \\
        \int_{x_2}^{x_2 + 2\uu(1-q(x_2))} T(x;r,s_r) \, dx \geq \uu (1-q(x_2))^2.
    \end{aligned}
    \end{equation*}
\end{proof}


We can improve upon the result of Lemma \ref{lem:conservative estimate of the error outside the delta interval} if we have an additional constraint that $x_2 - x_1 \leq 2\uu(q(x_2) - q(x_1))$.
\begin{lem}\label{lem:more sharp estimate of the error outside the delta interval}
    Let $r \equiv (p, q, 1 - p - q) \in \msf{R}$ be a receiver strategy. Let $x_1, x_2 \in \cg{X}$ be such that $x_1 \leq \h < x_2$, $p(x_i) + q(x_i) = 1$ for $i = 1, 2$, and $x_2 - x_1 \leq 2\uu(q(x_2) - q(x_1))$. Then
    \begin{equation*}
    \begin{aligned}
        &\int_{x_2 - 2\uu q(x_2)}^{x_1} T(x; r, s_r) \, dx &\geq& \uu \Big( q(x_2) - \frac{x_2 - x_1}{2\uu} \Big)^2 \quad \text{and} \\
        &\int_{x_2}^{x_1 + 2\uu(1 - q(x_1))} T(x; r, s_r) \, dx &\geq& \uu \Big( 1 - q(x_1) - \frac{x_2 - x_1}{2\uu} \Big)^2.
    \end{aligned}
    \end{equation*}
\end{lem}

\begin{proof}\label{proof:more sharp estimate of the error outside the delta}
     Let $q_1 := q(x_1)$ and $q_2 := q(x_2)$. We upper bound $q(s_r(x))$ (or $1-q(x)$ depending upon x) to get a upper bound on $T(x;r,s_r)$ and eventually over $\int_{x_2 - 2\uu q(x_2)}^{x_1} T(x;r,s_r)dx$ and $\int_{x_2}^{x_1 + 2\uu(1-q(x_1))}T(x;r,s_r)dx$.
    Let $x \in [x_2 - 2\uu q_2, x_1]$. Consider the sender strategy $\tilde{s}(t) = x_2$. From the definition of the worst-case lazy best response, we have $\mathbb{U}(x; r, s_r) \geq \mathbb{U}(x; r, \tilde{s})$, which implies
    \[
    2\uu q(s_r(x)) - |s_r(x) - x| \geq 2\uu q_2 - x_2 + x \quad \text{or} \quad 2\uu q(s_r(x)) \geq 2\uu q_2 - x_2 + x.
    \] Therefore,
    \begin{equation}\label{lem:eqn:lhs outside error}
        \int_{x_2 - 2\uu q_2}^{x_1} T(x; r, s_r) \, dx \geq \int_{x_2 - 2\uu q_2}^{x_1} \left( q_2 - \frac{x_2 - x}{2\uu} \right) \, dx = \uu \Big( q_2 - \frac{x_2 - x_1}{2\uu} \Big)^2.
    \end{equation}
    Similarly, let $x \in [x_2, x_1 + 2\uu(1 - q_1)]$ and consider the sender strategy $\hat{s}(t) = x_1$. Then, from the definition of the worst-case lazy best response, we have $\mathbb{U}(x; r, s_r) \geq \mathbb{U}(x; r, \hat{s})$, which implies
    \[
    2\uu(1 - q(s_r(x))) - |s_r(x) - x| \geq 2\uu(1 - q_1) - x + x_1 \quad \text{or} \quad 2\uu(1 - q(s_r(x))) \geq 2\uu(1 - q_1) - x + x_1.
    \]
    Hence, $T(x; r, s_r) = 1 - q(s_r(x)) \geq (1 - q_1) - \frac{x - x_1}{2\uu}$. Therefore,
    \begin{equation}\label{lem:eqn:rhs outside error}
        \int_{x_2}^{x_1 + 2\uu(1 - q_1)} T(x; r, s_r) \, dx \geq \int_{x_2}^{x_1 + 2\uu(1 - q_1)} \left( (1 - q_1) - \frac{x - x_1}{2\uu} \right) \, dx = \uu \Big( (1 - q_1) - \frac{x_2 - x_1}{2\uu} \Big)^2.
    \end{equation}
    Hence, combining \eqref{lem:eqn:lhs outside error} and \eqref{lem:eqn:rhs outside error}, we get
    \begin{equation*}
    \begin{aligned}
        &\int_{x_2 - 2\uu q(x_2)}^{x_1} T(x; r, s_r) \, dx \geq \uu \Big( q(x_2) - \frac{x_2 - x_1}{2\uu} \Big)^2 \quad \text{and} \\
        &\int_{x_2}^{x_1 + 2\uu(1 - q(x_1))} T(x; r, s_r) \, dx \geq \uu \Big( 1 - q(x_1) - \frac{x_2 - x_1}{2\uu} \Big)^2.
    \end{aligned}
    \end{equation*}
\end{proof}


In the next lemma, we show that any arbitrary receiver strategy $r \in \msf{R}$ will make at least $\frac{\uu}{2}$ errors.

\begin{lem}\label{lem:Error more than half}
    Let $r \in \msf{R}$ be a receiver strategy. Then $\cg{E}(r) \geq \frac{\uu}{2}$.
\end{lem}
\begin{proof}\label{proof:Error more than half}
    Let $r \in \msf{R}$ be a receiver strategy. Then, from Lemma \ref{lem:delta null wlog}, $r$ is of the form $r(x) \equiv (p(x), q(x), 1 - p(x) - q(x))$, where either $p(x) + q(x) = 1$ or $p(x) + q(x) = 0$. Let $k_1 := \sup\{x : p(x) + q(x) = 1, x \leq \h\}$ and $k_2 := \inf\{x : p(x) + q(x) = 1, x > \h\}$. From \eqref{eqn:receiver_condition_2}, both the infimum and supremum are attained in the set. Let $s_r(k_1) := l_1$ and $s_r(k_2) := l_2$. Let $q(k_1) := q_1$, $q(k_2) := q_2$, $q(l_1) := \bqo$, and $q(l_2) := \bqt$. We now calculate the receiver errors for various cases.\\
    \textbf{Case 1. }{$q_1 > q_2$.} \\
        From the definition of $\cg{E}(r)$, the following inequality holds:
        \[
        \cg{E}(r) = \int_\cg{X} T(x; r, s_r) \, dx \geq \int_{k_1 - 2\uu q_1}^{k_1} T(x; r, s_r) \, dx + \int_{k_2}^{k_2 + 2\uu(1 - q_2)} T(x; r, s_r) \, dx.
        \]
        Using Lemma \ref{lem:conservative estimate of the error outside the delta interval} by taking $x_1 = k_1$ and $x_2 = k_2$, we get
        \[
        \cg{E}(r) \geq \uu(q_1^2 + (1 - q_2)^2).
        \]
        Since from the case assumption, $q_1 > q_2$, therefore, we have
        \[
        \cg{E}(r) \geq \uu(q_1^2 + (1 - q_1)^2).
        \]
        The minimum value of $q_1^2 + (1 - q_1)^2$ is $\frac{1}{2}$ at $q_1 = \frac{1}{2}$. Hence, we have $\cg{E}(r) \geq \frac{\uu}{2}$.\\
\textbf{Case 2. }{$q_2 > q_1$ and $k_2 - k_1 < \uu(q_2 - q_1)$.} \\
        From the definition of $\cg{E}(r)$, the following inequality holds:
        \begin{equation}\label{lem:eqn:lower bound on error from outside}
            \cg{E}(r) = \int_\cg{X} T(x; r, s_r) \, dx \geq \int_{k_2 - 2\uu q_2}^{k_1} T(x; r, s_r) \, dx + \int_{k_2}^{k_1 + 2\uu (1 - q_1)} T(x; r, s_r) \, dx.
        \end{equation}
        Using Lemma \ref{lem:more sharp estimate of the error outside the delta interval} with $x_1 = k_1$ and $x_2 = k_2$, and using $k_2 - k_1 < \uu(q_2 - q_1)$, we get
        \[
        \cg{E}(r) \geq \uu((1 - q_1)^2 + q_2^2).
        \]
        Since $q_2 > q_1$, therefore,
        \[
        \cg{E}(r) \geq \uu((1 - q_2)^2 + q_2^2).
        \]
        The minimum value of $(1 - q_2)^2 + q_2^2$ is $\frac{1}{2}$ at $q_2 = \frac{1}{2}$. Hence, $\cg{E}(r) \geq \frac{\uu}{2}$.\\
   \textbf{Case 3. }{$q_2 \geq q_1$, $2\uu (q_2 - q_1) \geq k_2 - k_1 \geq \uu(q_2 - q_1)$, and $l_2 \leq k_1 < k_2 \leq l_1$.} \\
    We now estimate $T(x;r,s_r)$ for $x \in [k_1,\h]$. Let $x \in [k_1, \h]$. It follows that $h(x) = h(k_1) = \minus$. Using Lemma \ref{lem:self bestresponse} with $i = k_1$, $j = l_1$, and $w = x$, we have $s_r(x) = l_1$.
        Consider a sender strategy $s(x) := k_2$. Since $s_r$ is the worst-case lazy best response, we have
        \[
        2\uu q(s_r(x)) - l_1 + k_1 \geq 2\uu q_2 - k_2 + k_1.
        \]
        Since $l_1 \geq k_2$, therefore we have $q(s_r(x)) \geq q_2$. This implies
        \begin{equation}\label{lem:eqn:delta interval error lhs}
            T(x; r, s_r) \geq q_2 \quad \text{for } x \in [k_1, \h].
        \end{equation}
        Similarly, we estimate $T(x;r,s_r)$ for $x \in (\h,k_2]$. Let $x \in (\h, k_2]$. It follows that $h(x) = h(k_2) = \plus$. Using Lemma \ref{lem:self bestresponse} with $i = k_2$, $j = l_2$, and $w = x$, we have $s_r(x) = l_2$.
        Consider a sender strategy $\bar{s}(i) := k_2$. Since $s_r$ is the worst-case lazy best response, we have
        \[
        2\uu (1 - q(s_r(x))) - k_2 - l_2 \geq 2\uu (1 - q_1) - k_2 + k_1.
        \]
        Since $l_2 \geq k_1$, therefore we have $1 - q(s_r(x)) \geq 1 - q_1$. This implies
        \begin{equation}\label{lem:eqn:delta interval error rhs}
            T(x; r, s_r) \geq 1 - q_1 \quad \text{for } x \in (\h, k_2).
        \end{equation}
        From the definition of $\cg{E}(r)$, we have
        \begin{equation}\label{lem:eqn:error in delta interval2}
            \cg{E}(r) := \int_\cg{X} T(x; r, s_r) \, dx \geq \int_{k_2 - 2\uu q_2}^{k_1} T(x; r, s_r) \, dx + \int_{k_2}^{k_1 + 2\uu (1 - q_1)} T(x; r, s_r) \, dx + \int_{k_1}^{k_2} T(x; r, s_r) \, dx.
        \end{equation}
        Using Lemma \ref{lem:more sharp estimate of the error outside the delta interval} with $x_1 = k_1$, $x_2 = k_2$, and substituting \eqref{lem:eqn:delta interval error lhs} and \eqref{lem:eqn:delta interval error rhs} into \eqref{lem:eqn:error in delta interval2}, we get
        \[
        \cg{E}(r) \geq \uu \left(q_2 - \frac{k_2 - k_1}{2\uu}\right)^2 + \uu \left((1 - q_1) - \frac{k_2 - k_1}{2\uu}\right)^2 + (1 - q_1)(\h - k_1) + q_2(k_2 - \h).
        \]
        Let $k_2 - k_1 := w\uu (q_2 - q_1)$, where $1 \leq w \leq 2$.
        Suppose $1 - q_1 > q_2$ (a similar argument will hold when $1 - q_1 \leq q_2$). Then
        \[
        (1 - q_1)(\h - k_1) + q_2(k_2 - \h) \geq q_2(k_2 - k_1) = \uu w (q_2 - q_1) q_2,
        \]
        which implies
        \[
        \cg{E}(r) \geq \uu \left(q_2 - w \frac{q_2 - q_1}{2}\right)^2 + \uu \left((1 - q_1) - w \frac{q_2 - q_1}{2}\right)^2 + \uu w (q_2 - q_1) q_2.
        \]
        The minimum value of
        \[
        \uu \left(q_2 - w \frac{q_2 - q_1}{2}\right)^2 + \uu \left((1 - q_1) - w \frac{q_2 - q_1}{2}\right)^2 + \uu w (q_2 - q_1) q_2
        \]
        is $\frac{1}{2}$ at $q_1 = \frac{1}{2}, q_2 = \frac{1}{2}$, and $w = 2$. Hence, $\cg{E}(r) = \frac{\uu}{2}$.\\
\textbf{Case 4. }{$q_2 \geq q_1$, $k_2 - k_1 > 2\uu(q_2 - q_1)$, and $l_2 \leq k_1 < k_2 \leq l_1$.} \\
        Using the same argument used to derive \eqref{lem:eqn:delta interval error lhs} and \eqref{lem:eqn:delta interval error rhs} in Case 3, we get
        \begin{equation}\label{lem:eqn:delta interval error lhs2}
            T(x; r, s_r) \geq q_2 \quad \text{for } x \in [k_1, \h]
        \end{equation}
        and
        \begin{equation}\label{lem:eqn:delta interval error rhs2}
            T(x; r, s_r) \geq (1 - q_1) \quad \text{for } x \in (\h, k_2].
        \end{equation}
        From the definition of $\cg{E}(r)$, we have
        \begin{equation}\label{lem:eqn:error in delta interval}
            \cg{E}(r) := \int_\cg{X} T(x; r, s_r) \, dx \geq \int_{k_1 - 2\uu q_1}^{k_1} T(x; r, s_r) \, dx + \int_{k_2}^{k_2 + 2\uu (1 - q_2)} T(x; r, s_r) \, dx + \int_{k_1}^{k_2} T(x; r, s_r) \, dx.
        \end{equation}
        Using Lemma \ref{lem:conservative estimate of the error outside the delta interval} with $x_1 = k_1$, $x_2 = k_2$, and substituting \eqref{lem:eqn:delta interval error lhs2} and \eqref{lem:eqn:delta interval error rhs2} into \eqref{lem:eqn:error in delta interval}, we get
        \[
        \cg{E}(r) \geq \uu q_1^2 + \uu (1 - q_2)^2 + (1 - q_1)(\h - k_1) + q_2(k_2 - \h).
        \]
        Suppose $q_2 > 1 - q_1$ (a similar argument will hold when $1 - q_1 \leq q_2$). Then
        \[
        (1 - q_1)(\h - k_1) + q_2(k_2 - \h) \geq (1 - q_1) |k_2 - k_1| > 2\uu(q_2 - q_1)(1 - q_1),
        \]
        which implies
        \[
        \cg{E}(r) > \uu \left(q_1^2 + (1 - q_2)^2 + 2(q_2 - q_1)(1 - q_1)\right).
        \]
        The minimum value of $q_1^2 + (1 - q_2)^2 + 2(q_2 - q_1)(1 - q_1)$ is $\frac{1}{2}$ at $q_1 = q_2 = \frac{1}{2}$. Hence, $\cg{E}(r) > \frac{\uu}{2}$.\\
\textbf{Case 5. }{$q_2 \geq q_1$, $k_2 - k_1 > \uu(q_2 - q_1)$, and $l_1 \leq k_1 < k_2 \leq l_2$.} \\
    Here we will bound $T(x;r,s_r)$, To do so, we first show that $\bqo > q_1$ and $\bqt < q_2$.
        From Lemma \ref{lem:self bestresponse} with $i = k_1$, $j = l_1$, it follows that $s_r(l_1) = l_1$ and $s_r(l_2) = l_2$. Recall from the notation mentioned at the start of the proof that $q(l_1) := \bqo$ and $q(l_2) := \bqt$. We now show that $l_2 - l_1 \geq 2\uu (\bqt - \bqo)$. Let $s(i) := l_2$ be a sender strategy. Since $s_r$ is the worst-case lazy best response, we have $\mathbb{U}(l_1, r, s_r) \geq \mathbb{U}(l_1, r, s)$. This implies
        \[
        2\uu \bqo \geq 2\uu \bqt - l_2 + l_1 \quad \text{or} \quad 2\uu (\bqt - \bqo) \leq l_2 - l_1
        \]
        as claimed.
        Now let $\bar{s}(i) := k_1$ be another sender strategy. Then it follows that $\mathbb{U}(k_1, r, s_r) \geq \mathbb{U}(k_1, r, \bar{s})$, which implies
        \begin{equation}\label{eqn:lem:bar greater}
            2\uu \bqo - k_1 + l_1 > 2\uu q_1 \quad \text{or} \quad \bqo > q_1.
        \end{equation}
        This is a strict inequality since $k_1$ is the cheaper option for $k_1$ compared to $l_1$. A similar reasoning shows that
        \begin{equation}\label{eqn:lem:bar lesser}
            \bqt < q_2.
        \end{equation}
        Next, we bound $s_r(x)$ for $x \in [m_1,m_2]$ where $m_1$ and $m_2$ are defined as:
        \[
        m := \frac{k_1 + k_2}{2}, \quad m_1 := m - \uu(q_2 - \bqo), \quad m_2 := m + \uu(\bqt - q_1).
        \]\\
        \textbf{Claim. } $m_1 \geq k_1$ and $m_2 < k_2$. \\
        Let $m \leq \h$ (a similar argument holds when $m > \h$). We now show that $m_1 \geq k_1$. Let $x \in (m_1, \h]$. Suppose $s_r(x) \leq x$, then $s_r(x) \leq k_1$. This is because from the definition of $k_1$, we have $r(x) = \Delta$ for $x \in [k_1, \h]$. But $h(x) = h(s_r(x)) = \minus$, therefore $s_r(s_r(x)) = s_r(x)$ from Lemma \ref{lem:self bestresponse}. Furthermore, we have
        $h(l_1) = h(k_1) = h(s_r(x)) = \minus$. Therefore, from Lemma \ref{lem:self bestresponse} with $i = k_1$, $j = l_1$, and $w = s_r(x)$, we have $s_r(s_r(x)) = s_r(x) = l_1$. Now consider the sender strategy $s(x) := k_2$. From the definition of the worst-case lazy best response, we have $\mathbb{U}(x, r, s_r) \geq \mathbb{U}(x, r, s)$, which implies
        \[
        2\uu \bqo - x + l_1 \geq 2\uu q_2 - k_2 + x.
        \]
        Since $x > m_1$, we have
        \[
        2\uu \bqo - x + l_1 < 2\uu \bqo - m_1 + l_1 \quad \text{and} \quad 2\uu q_2 - k_2 + x > 2\uu q_2 - k_2 + m_1.
        \]
        Hence, we have
        \[
        2\uu \bqo - m_1 + l_1 > 2\uu q_2 -         k_2 + m_1,
        \]
        which implies
        \[
        2\uu (\bqo - q_2) > 2m_1 - l_1 - k_2.
        \]
        Substituting for $m_1$ and simplifying, we get
        \[
        2\uu (\bqo - q_2) > k_1 + k_2 - 2\uu (q_2 - \bqo) - l_1 - k_2,
        \]
        which implies
        \[
        0 > k_1 - l_1.
        \]
        This is a contradiction since $k_1 \geq l_1$. Therefore, our assumption that $s_r(x) \leq x$ is false. Hence,
        \begin{equation}\label{eqn:lem:not crossing boundary 1}
            s_r(x) > x \quad \text{for} \quad x \in (m_1, \h].
        \end{equation}
        If $m_1 < k_1$, then by taking $x = k_1$, we get $s_r(x) = l_1 > k_1$, which is false as $l_1 \leq k_1$ by case assumption. Hence, $m_1 \geq k_1$. A similar reasoning shows
        \begin{equation}\label{eqn:lem:not crossing boundary 2}
            s_r(x) < x \quad \text{for} \quad x \in (\h, m_2)
        \end{equation}
        and that $m_2 \leq k_2$.

        Next, we calculate the error in the interval $(m_1, m_2)$.
        Let $x \in (m_1, \h]$. Consider the sender strategy $s(x) := k_2$. Since $s_r$ is the worst-case lazy best response, we have $\mathbb{U}(x, r, s_r) \geq \mathbb{U}(x, r, s)$, which implies
        \[
        2\uu q(s_r(x)) - s_r(x) + x \geq 2\uu q_2 - k_2 + x \quad \text{or} \quad q(s_r(x)) \geq q_2,
        \]
        where we have used \eqref{eqn:lem:not crossing boundary 1} and the fact that $r(x) = \Delta$ for $x \in (\h, \jb)$ (which implies $s_r(x) \geq k_2$). Thus,
        \begin{equation}\label{eqn:lem:error function case 5 1}
            T(x; r, s_r) \geq q_2 \quad \text{for} \quad x \in (m_1, \h].
        \end{equation}
        Using similar reasoning for $x \in (\h, m_2)$, we have $q(s_r(x)) \leq q_1$, and hence
        \begin{equation}\label{eqn:lem:error function case 5 2}
            T(x; r, s_r) \geq 1 - q_1 \quad \text{for} \quad x \in (\h, m_2].
        \end{equation}
        From the definition of $\cg{E}(r)$, we have
        \begin{equation}\label{eqn:lem:total error case 5}
            \cg{E}(r) := \int_\cg{X} T(x; r, s_r) \, dx \geq \int_{l_1 - 2\uu q_1}^{l_1} T(x; r, s_r) \, dx + \int_{l_2}^{l_2 + 2\uu (1 - q_2)} T(x; r, s_r) \, dx + \int_{m_1}^{m_2} T(x; r, s_r) \, dx.
        \end{equation}
        Using Lemma \ref{lem:conservative estimate of the error outside the delta interval} with $x_1 = l_1$, $x_2 = l_2$, and substituting \eqref{eqn:lem:error function case 5 1} and \eqref{eqn:lem:error function case 5 2} in \eqref{eqn:lem:total error case 5}, we get
        \begin{equation}\label{eqn:lem:total error case 5 1}
            \cg{E}(r) \geq \uu \bqo^2 + \uu(1 - \bqt)^2 + (\h - m_1)q_2 + (m_2 - \h)(1 - q_1).
        \end{equation}
        Suppose $q_2 < 1 - q_1$ (a similar argument will hold when $1 - q_1 < q_2$).
        Then, we have
        \begin{equation}\label{eqn:lem:bar dominates}
            (\h - m_1)q_2 + (m_2 - \h)(1 - q_1) \geq (m_2 - m_1)q_2 = \uu(\bqt - \bqo + q_2 - q_1)q_2.
        \end{equation}
        Using \eqref{eqn:lem:bar greater}, \eqref{eqn:lem:bar lesser}, and \eqref{eqn:lem:bar dominates} in \eqref{eqn:lem:total error case 5 1}, and using $q_2 - q_1 \geq 0$, we obtain
        \[
        \cg{E}(r) \geq \uu \bqo^2 + \uu(1 - \bqt)^2 + 2\uu(\bqt - \bqo)\bqt.
        \]
        The minimum of the expression $\uu \bqo^2 + \uu(1 - \bqt)^2 + 2\uu(\bqt - \bqo)\bqt$ is $\frac{1}{2}$, and it occurs when $\bqo = \bqt = \frac{1}{2}$. Hence, we conclude that
        \[
        \cg{E}(r) = \frac{\uu}{2}.
        \]\\
\textbf{Case 6. }{$q_2 \geq q_1$, $k_2 - k_1 > \uu(q_2 - q_1)$, and either $k_2 \leq \min\{l_1, l_2\}$ or $k_1 \geq \max\{l_1, l_2\}$.} \\
    We will show that this case does not arise. We will prove this for $k_1 \geq \max\{l_1, l_2\}$. The proof for $k_2 \leq \min\{l_1, l_2\}$ is similar.
Consider the following sender strategies:
    \[
    \bar{s}(i) := l_1, \quad s'(i) := l_2.
    \]
    Since $s_r$ is the worst-case lazy best response, we have
    \[
    \mathbb{U}(k_1; r, s_r) \geq \mathbb{U}(k_1; r, s') \quad \text{and} \quad \mathbb{U}(k_2; r, s_r) \geq \mathbb{U}(k_2; r, \bar{s}),
    \]
    which implies
    \begin{equation}\label{lem:eqn:contradiction 1}
        2\uu \bqo - k_1 + l_1 \geq 2\uu \bqt - k_1 + l_2 \quad \text{or} \quad 2\uu (\bqo - \bqt) \geq l_2 - l_1
    \end{equation}
    and
    \begin{equation}\label{lem:eqn:contradiction 2}
        2\uu (1 - \bqt) - k_2 + l_2 \geq 2\uu (1 - \bqo) - k_2 + l_1 \quad \text{or} \quad 2\uu (\bqo - \bqt) \geq l_1 - l_2,
    \end{equation}
    respectively.
Combining \eqref{lem:eqn:contradiction 1} and \eqref{lem:eqn:contradiction 2}, we get $l_1 = l_2$.
Now, consider the strategy $s(i) = k_1$. Since $s_r$ is the worst-case lazy best response, we have
    \[
    \mathbb{U}(k_1; r, s_r) > \mathbb{U}(k_1; r, s) \quad \text{and} \quad \mathbb{U}(k_2; r, s_r) > \mathbb{U}(k_2; r, s).
    \]
    The inequality is strict for the sender, $k_1$ is the less costly (lazy) option compared to $l_1$ and $l_2$ at $k_1$ and $k_2$, respectively. This implies
    \begin{equation}\label{eqn:lem:barq1 >= q1}
        2\uu \bqo - k_1 + l_1 > 2\uu q_1 \quad \text{or} \quad \bqo > q_1
    \end{equation}
    and
    \begin{equation}\label{eqn:lem:barq1 <= q1}
        2\uu (1 - \bqo) - k_2 + l_2 > 2\uu (1 - q_1) - k_2 + k_1 \quad \text{or} \quad 2\uu (q_1 - \bqo) > k_1 - l_2 = k_1 - l_1 \quad \text{or} \quad q_1 > \bqo.
    \end{equation}

    Combining \eqref{eqn:lem:barq1 <= q1} and \eqref{eqn:lem:barq1 >= q1}, we get $q_1 > q_1$, which is a contradiction. Hence, this case does not arise.

\end{proof}

        \subsection{Omitted Calculations in Section III.E}

        \input{Calculation Section 3E}

        \subsection{Omitted Proofs of Section IV.A}

        \begin{apdxproof}(Proof of Proposition \ref{prop:Deterministic Stackelberg equilibrium for multiple decision boundaries})\label{proof:Deterministic Stackelberg equilibrium for multiple decision boundaries}
    \setcounter{claims}{0}
	Let $r^\star$ be a Stackelberg equilibrium, and define:
	\begin{equation*}
		\begin{aligned}
			\theta^{(n)} &:= \sup (E(r^\star)^c \cap O_n), \\
			\theta_{(n)} &:= \inf (E(r^\star)^c \cap O_n), \\
			\Theta_n &:= [\theta_{(n)}, \theta^{(n)}].
		\end{aligned}
	\end{equation*}
	This implies $E(r^\star)^c \subseteq \cup_{n=1}^{d-1} \Theta_n$. We demonstrate that with this choice of $\Theta_n$, the receiver strategy $r(x;\Theta_1,\cdots, \Theta_{d-1})$ given in \eqref{eqn:stack multiple} is a Stackelberg equilibrium. Consider the following cases for $\Theta_n$'s. \\

	\begin{case}{All $\Theta_n$'s are empty.}\\
    This case will not happen because if all $\Theta_n$'s are empty, then $r^\star$ misclassifies all feature vectors. However, a majority-label (described in Section \ref{sec:Problem formulation}) receiver strategy will correctly classify at least half of them, contradicting that $r^\star$ is a Stackelberg equilibrium. Thus, at least one $\Theta_n$ is non-empty.
	\end{case}

\begin{case}{For all odd (or even) $n$, $\Theta_n$ is empty.}\\
Suppose $\Theta_n = \emptyset$ for all odd $n$. For even $n$, denote the corresponding $\Theta_n$ as $\Theta_{2k}$, where $k \in \left\{1, \dots, \left\lfloor \frac{d-1}{2} \right\rfloor \right\}$.
Since $r(x) = \Delta$ for all $x \in \cg{X} \setminus \bigcup_k \Theta_{2k}$, it follows that for $x \in 
 \cg{X}$, $s_r(x) \in \bigcup_k \Theta_{2k}$. Thus, $r(s_r(x)) \equiv h(s_r(x))$ from \eqref{eqn:stack multiple}.
From \eqref{eqn:multiple boundary true function}, we have $ h(t) = \minus$ for $t \in \cup_kO_{2k}$. Therefore, $r(s_r(x)) \equiv \minus$ and
$r(s_r(x)) = h(x)$ for $x \in \cup_kO_{2k}.$ Hence 
\[
E(r)^c \supseteq \bigcup_{2k} O_k \supseteq \bigcup_{2k} \Theta_k.
\]
Since $r^\star$ is a Stackelberg equilibrium, it follows that $E(r^\star) \subseteq \bigcup_k \Theta_{2k}$ by construction. Consequently, $E(r)^c = E(r^\star) = \bigcup_k O_{2k}$, and $r$ is a Stackelberg equilibrium.
A similar argument applies to show that $r$ is a Stackelberg equilibrium when $n$ for which $\Theta_n = \emptyset$ are all even.
\end{case}

\begin{case}{Their exist an odd $k$ and even $l$ such that $\Theta_k, \Theta_l \neq \emptyset$.}\\
We will first show that $\theta_{(\max\{l,k\})} - \theta^{(\min\{l,k\})} \geq \uu$. Suppose $k < l$ (similar reasoning holds when $k>l$).
Suppose $\theta_{(l)} - \theta^{(k)} < \uu$. Let $\eps >0$ be such that it satisfies $0 < 2\eps < \uu - (\theta_{(l)} - \theta^{(k)})$. Let $\theta^{(k')}, \theta_{(l')} \in \cg{X}$ be feature vectors such that $\theta^{(k')}, \theta_{(l')} \in E(r^\star)^c$,  $0 \leq \theta_{(l')} - \theta_{(l)} < \eps$, and $0 \leq \theta^{(k)} - \theta^{(k')} < \eps$.  Such feature vectors exist by the definition of $\theta^{(k)}$ and $\theta_{(l)}$. It follows that 
$$\theta_{(l')} - \theta^{(k')} < 2\eps + \theta_{(l)} - \theta^{(k)} < \uu.$$ Note that $h(\theta_{(l')}) = \minus$ and $h(\theta^{(k')}) = \plus$. Suppose $\mathbb{U}(\theta^{(k')};r^\star,s_{r^\star}) \leq \mathbb{U}(\theta_{(l')};r^\star,s_{r^\star})$ (similar arguments applies when $\mathbb{U}(\theta^{(k')};r^\star,s_{r^\star}) > \mathbb{U}(\theta_{(l')};r^\star,s_{r^\star})$).
Let $s$ be a sender strategy defined as $s(x) := s_{r^\star}(\theta_{(l')})$ for all $x \in \cg{X}$. Then $$\mathbb{U}(\theta^{(k')}; r^\star, s) = \uu - |\theta^{(k')} -s_{r^\star}(\theta_{(l')})| \geq   \uu - |\theta_{(l')} - \theta^{(k')}|- |\theta_{(l')} - s_{r^\star}(\theta_{(l')})| > -|\theta_{(l')} - s_{r^\star}(\theta_{(l')})| = \mathbb{U}(\theta_{(l')};r^\star,s_{r^\star}).$$ This implies 
$\mathbb{U}(\theta^{(k')};r^\star,s) > \mathbb{U}(\theta^{(k')};r^\star,s_{r^\star})$, which contradicts the fact that $s_{r^\star}$ is the worst-case best response. Therefore, $\theta_{(l)} - \theta^{(k)} \geq \uu$, proving the claim.

Now consider the receiver strategy $r(x; \{\Theta_n\})$. We calculate its worst-case best response $s_r$. Let $t \in \Theta_i$. Since $r(x) = \Delta$ for $x \notin \cup_n \Theta_n$, we have $s_r(x) \in \cup_n \Theta_n$. Suppose $i$ is even (similar arguments hold for odd $i$) and $s_r(t) \geq t$ (a similar arguments hold for $s_r(t) < t$). The sender's payoff is:
\[
\mathbb{U}(t; r, s_r) = 
\begin{cases}
-s_r(t) + t & \text{if } s_r(t) \in \Theta_j, \, j \text{ even}, \\
\uu - s_r(t) + t & \text{if } s_r(t) \in \Theta_j, \, j \text{ odd}.
\end{cases}
\]
If $j$ is even, $\mathbb{U}(t; r, s_r) = -s_r(t) + t \leq 0$, with the maximum value $0$ at $s_r(t) = t$. If $j$ is odd, then using $\theta_{(l)} - \theta^{(k)} \geq \uu$, we get $s_r(t) - t \geq \uu$, and hence $\mathbb{U}(t; r, s_r) \leq 0$. Thus $s_r(t) = t$ gives the maximum payoff to the sender. Therefore, $r(s_r(t)) = h(t)$ for all $t \in \cup_n \Theta_n$. This gives $E(r)^c \supseteq \cup_n \Theta_n \supseteq E(r^\star)$, proving that $r$ is a Stackelberg equilibrium.
\end{case}
\end{apdxproof}

        \subsection{Omitted Proofs of Section IV.B}

        \begin{apdxproof}(Proof of Lemma \ref{lem:same true function})\label{proof:same true function}
Let $r_1 := r^\star_{f,D},\; r_2 := r^\star_{f,D(S)}$. Let $\psi(x)$ be the distribution function of the probability law $D$ and $\psi_m(x)$ be the distribution function of the empirical probability law $D(S)$ for a sample $S$, i.e.,
    \begin{equation*}
        \psi_m(x) := \frac{1}{m} \sum_{x_i \in S} \mathbbm{1}_{[0,x]}.
    \end{equation*}
According to the Dvoretzky–Kiefer–Wolfowitz inequality (\cite{Dvoretzky-Kiefer-Wolfowitz-1956}), for any $(\epsilon, \delta) \in (0,1)^2$, there exists an $\bar{m}(\eps,\delta)$ such that for all $m > \bar{m}(\eps,\delta)$, we have:
    \begin{equation}\label{eqn:DKW inequality:lem: same true function}
        D^{\otimes m}\big(\sup_{x \in \cg{X}} |\psi(x) - \psi_m(x)| < \epsilon\big) > 1 - 2e^{-2m\epsilon^2}.
    \end{equation}
Let 
\begin{equation}\label{eqn:DKW inequality:lem: same true function good set}
    \cg{S} := \{S: \sup_{x \in \cg{X}}|\psi(x) - \psi_m(x)| < \eps\}.
\end{equation}
Then from \eqref{eqn:DKW inequality:lem: same true function}, we know that $D^{\otimes m}(\cg{S}) > 1-2e^{-2em\eps^2}$. Fix a $S \in \cg{S}$ and let $\psi_m$ be the corresponding empirical distribution function. 
Consider the receiver strategies $r_1$ and $r_2$. From Proposition \ref{prop:Deterministic Stackelberg equilibrium for multiple decision boundaries}, we have $r_1(x) \equiv r_1(x; \Theta^{(1)}_1,\cdots,\Theta^{(1)}_{d-1})$ and $r_2(x) \equiv r_2(x;\Theta^{(2)}_1,\cdots,\Theta^{(2)}_{d-1})$ for some $\Theta^{(i)}_n$, where $i = 1, 2$ and $n \in \{0,1,\cdots,d-1\}$, and where $\Theta^{(i)}_n := [\theta_{(i,n)}, \theta^{(i,n)}]$ is a closed interval. Let $E_n^{(i)} := (\theta^{(i,n)}, \theta_{(i,n+1)})$. From Lemma \ref{prop:Deterministic Stackelberg equilibrium for multiple decision boundaries}, $E_n^{(i)}$ is the set of feature vectors misclassified by $r_i$ near the decision boundary $n$. From the definition of $\psi_m$ and $\psi$, we have
    \begin{equation*}
        D(S)(E_n^{(i)}) = \psi_m(\theta^{(i,n+1)}) - \psi_m(\theta_{(i,n)}),\quad D(E_n^{(i)}) = \psi(\theta^{(i,n+1)}) - \psi(\theta_{(i,n)}),\quad i \in \{1, 2\}, n \in \{1,\cdots,d-1\}.
    \end{equation*}
Since $S \in \cg{S}$, using \eqref{eqn:DKW inequality:lem: same true function good set}, we get 
    \begin{equation}\label{lem:eqn:actual dominate emp}
        \begin{aligned}
            |D(S)(E_n^{(i)}) - D(E_n^{(i)})| &\leq |\psi_m(\theta^{(i,t+1)}) - \psi(\theta^{(i,n+1)})| + |\psi(\theta_{(i,n)}) - \psi_m(\theta_{(i,n)})| < 2\epsilon,\\
            &\text{for all }i \in \{1, 2\}, n \in \{1,\cdots,d-1\}.
        \end{aligned}
    \end{equation}
    Summing over all $d-1$ decision boundaries, and noting that $\mathcal{E}(r_i;f, D) = \sum_n D(E_n^{(i)})$ and $\cg{E}(r_i;f, D(S)) = \sum_n D(S)(E_n^{(i)})$, we get
    \[
        |\cg{E}(r_2;f, D) - \cg{E}(r_1;f, D)| < 2(d-1)\epsilon.
    \]
This holds true for all $S \in \cg{S}$. But we have $D^{\otimes m}(\cg{S}) > 1-2e^{-2m\eps^2}$. Therefore, we get
\[
       D^{\otimes m}(|\cg{E}(r_2;f, D) - \cg{E}(r_1;f, D)| < 2(d-1)\epsilon) > 1- 2e^{-2m\eps^2}.
    \]
as required.

\end{apdxproof}





\begin{apdxproof}(Proof of Lemma \ref{lem:same distribution})\label{proof:same distribution}
Let $r_1 := r^\star_{h,D}$, $r_2 := r^\star_{A(S),D}$, and $E(r,f) := E(r,s_r,f)$.
Since VC dim$(\cg{F}) < \infty$,  $\cg{F}$ is PAC learnable from Theorem \ref{thm:Fundamental theorem of statistical learning}. Furthermore, from Theorem  \ref{thm:Fundamental theorem of statistical learning} we have
\begin{equation}\label{eq:concentration bound}
    D^{\otimes m}(L_{D}(A(S);h) < \epsilon) > 1 - \delta,
\end{equation}
where $m \geq \underbar{m}(\eps,\delta)$ and $\underbar{m}(\eps,\delta)$ satisfies 
\begin{equation*}
    C_1 \frac{d + \log\left(\frac{1}{\delta}\right)}{\epsilon} \leq \underbar{m}(\eps,\delta) \leq C_2 \frac{d \log\left(\frac{1}{\epsilon}\right) + \log\left(\frac{1}{\delta}\right)}{\epsilon}.
\end{equation*}
Here $C_1, C_2 \geq 0$ are constants and $d$ is the VC dimension of $\mathcal{F}$.
Let 
\begin{equation}\label{eqn:good set}
    \cg{S} := \{S: L_D(A(S);h) < \eps\}.
\end{equation}
From \eqref{eq:concentration bound}, we know that $D^{\otimes m}(\cg{S}) > 1-\delta$.
Fix a $S \in \cg{S}$. Define $W(S) := \{x \mid h(x) \neq A(S)(x)\}$. For a receiver strategy $r$ and (an arbitrary) true function $f$ (recall from Definition \ref{defn:Best response} that $s_r$ will depends upon the true function $f$ v.i.a. utility function), the classification error set $E(r;A(S))$ (defined in \eqref{eqn:error set}) can be decomposed as
\[
E(r; A(S)) = (E(r; A(S)) \cap W(S)) \cup (E(r; A(S)) \cap W(S)^c).
\]
Since $S \in \cg{S}$, using \eqref{eqn:good set}, we get $L(A(S);h) < \eps$, and therefore it follows that $D(W(S)) < \epsilon$. Hence, we get
\[
\mathcal{E}(r; A(S),D) < \epsilon + D(E(r; A(S)) \cap W(S)^c).
\]
From the definition of $W(S)$, $h(x) = A(S)(x)$ for all $x \in W(S)^c$. Therefore,
\begin{equation}
    \begin{aligned}
        E(r;A(S)) \cap W(S)^c &= \{x:r(s_r(x)) = A(S)(x)\} \cap \{x: h(x) =A(S)(x)\} = \{x:r(s_r(x)) = A(S)(x) = h(x)\}\\  &= \{x: h(x) =A(S)(x)\} \cap\{x:r(s_r(x)) =h(x) \}      
        = E(r;h) \cap W(S)^c,\\
    \end{aligned}
\end{equation}
and thus
\begin{equation}\label{eqn:lem:change in ground truth 3}
    \mathcal{E}(r; A(S),D) < \epsilon + D(E(r; h)) = \epsilon + \mathcal{E}(r; h,D).
\end{equation}
Similarly, we can decompose the classification error set $E(r;h)$ as
\[E(r;h) := (E(r;h) \cap W(S)) \cup (E(r;h)\cap W(S)^c).\]
Repeating the argument used for deriving \eqref{eqn:lem:change in ground truth 3}, we obtain
\begin{equation}\label{eqn:lem:change in ground truth 6}
    \mathcal{E}(r; h,D) < \epsilon + \mathcal{E}(r; A(S),D).
\end{equation}
Combining \eqref{eqn:lem:change in ground truth 3} and \eqref{eqn:lem:change in ground truth 6}, we get
\[|\cg{E}(r;h,D) - \cg{E}(r,A(S),D)| < \eps.\] This holds true for all $S \in \cg{S}$. But we have $D^{\otimes m}(\cg{S}) > 1-\delta$, therefore we get 
\begin{equation}\label{eqn:lem:change in ground truth 5}
    D^{\otimes m}(|\mathcal{E}(r; h) - \mathcal{E}(r; A(S))|  < \epsilon) > 1-\delta \text{for all receiver strategy} r \text{and true function} f.
\end{equation}
This proves the $(a)$ part.
Using part $(a)$, we now establish
\begin{equation}
    D^{\otimes m}(\cg{E}(r_2;h)) - \cg{E}(r_1;h) < 2\eps) > 1-\delta,
\end{equation}
where $h$ is the ground true function of $r_1$ and $A(S)$ is the estimated true function of $r_2$. This is the same as the claim (given in \eqref{eqn:lem:vanilla ERM PAC bound}) we want to prove.
Substituting $r = r_1$ and $f=h$ (ground true function) in \eqref{eqn:lem:change in ground truth 3}, we get
\begin{equation}\label{eqn:lem:change in ground truth 3.1}
    \mathcal{E}(r_1; A(S),D) < \epsilon + \mathcal{E}(r_1; h,D).
\end{equation}
Substituting $r=r_2$ and $f = A(S)$ (estimated true function) in \eqref{eqn:lem:change in ground truth 6}, we get
\begin{equation}\label{eqn:lem:change in ground truth 6 1}
    \mathcal{E}(r_2; h,D) < \epsilon + \mathcal{E}(r_2; A(S),D).
\end{equation}
Consider the following:
\begin{equation}\label{eqn:lem:change in ground truth 4}
\begin{aligned}
    \mathcal{E}(r_2; h,D) - \mathcal{E}(r_1; h,D) &= \underbrace{\big(\mathcal{E}(r_2; h,D) - \mathcal{E}(r_2; A(S),D)\big)}_a  + \underbrace{\big(\mathcal{E}(r_2; A(S),D) - \mathcal{E}(r_1; A(S),D)\big)}_b \\
    &\quad + \underbrace{\big(\mathcal{E}(r_1; A(S),D) - \mathcal{E}(r_1; h,D)\big)}_c.
\end{aligned}
\end{equation}
From the definition of $r_2$, we know $b = \mathcal{E}(r_2; A(S),D) - \mathcal{E}(r_1; A(S),D) \leq 0$ as $\cg{E}(r_2;A(S),D)$ is the Stackelberg error. From \eqref{eqn:lem:change in ground truth 6 1}, we have $a < \eps$ and using \eqref{eqn:lem:change in ground truth 3.1}, we have $c < \eps$. Therefore we get
\[
\mathcal{E}(r_2; h,D) - \mathcal{E}(r_1; h,D) < 2\epsilon.
\]
This holds true for all $S \in \cg{S}$. But we have $D^{\otimes m}\cg{S} > 1-\delta$. Therefore we get 
\[D^{\otimes m}(\mathcal{E}(r_2; h,D) - \mathcal{E}(r_1; h,D)) < 2\epsilon
\]
as claimed in part $(b)$.
    
\end{apdxproof}


\bibliographystyle{IEEEtran}
\bibliography{Reference}


\end{document}
